\chapter{语音任务}
\label{chap:nlp-speech-tasks}

用户问“读书会是四月十八日吗”，系统即使写对了每个字，也可能在用户尚未说完时抢答，或者在回答时把正确的日期读成“四月十七日”。第\ref{chap:nlp-question-answering-dialogue}章已经区分回答依据与对话状态；声音接口又带来两类新问题。一类是连续的声学观测怎样对应离散的语言单位，另一类是正在进行的交流应在何时接收、确认和输出。识别、口语理解、合成、翻译与实时对话分别交付文字、语义、声音、跨语言结果和完整交互，其中一项正确不能代替其余各项。

围绕这句询问，本章要处理六项挑战，见表\ref{tab:nlp-speech-challenges}。其一，声音与文字长度不同：一秒声音有近百个特征帧，却只对应两三个字，转写也不标注每个字落在哪几帧。其二，同一内容的声音受口音、噪声和语速影响，变化很大。其三，带人工转写的录音昂贵，无标注录音却很丰富。其四，离线系统可以读完整句再决定，在线系统却要在低延迟下提交部分结果，后半句的更正可能推翻已经输出的内容。其五，合成与翻译要生成声音，同一句话没有唯一正确的读法，系统却必须保住日期、数字与地点等内容。其六，实时交互还要协调话轮、打断与业务状态，声音一旦播出就无法撤回。

\begin{table*}[htbp]
  \centering
  \small
  \begin{tabularx}{\linewidth}{@{}>{\raggedright\arraybackslash}p{40mm}
    >{\raggedright\arraybackslash}X>{\raggedright\arraybackslash}p{34mm}@{}}
    \toprule
    挑战 & 通知中的例子 & 主要处理位置\\
    \midrule
    声音与文字长度不同，没有帧级对齐 & “四月十八日”约1.5秒、上百帧，只对应五个字 & 第\ref{sec:nlp-speech-representation}、\ref{sec:nlp-speech-recognition}节\\
    同一内容的声学变化大 & “下午两点”被快速读出，或混入空调噪声 & 第\ref{sec:nlp-speech-features}、\ref{sec:nlp-speech-ssl}小节\\
    配对转写少而无标注录音多 & 积累了大量讲座录音，却只有少量人工转写 & 第\ref{sec:nlp-speech-ssl}小节\\
    离线全句可见与在线低延迟冲突 & “不是十七日，是十八日”的更正出现在后半句 & 第\ref{sec:nlp-speech-recognition}、\ref{sec:nlp-speech-postprocessing}、\ref{sec:nlp-speech-understanding}节\\
    生成声音无唯一答案却须保持内容 & “14:00”可读作“十四点”或“下午两点”；译成英语语音时日期不能变 & 第\ref{sec:nlp-speech-tts}、\ref{sec:nlp-speech-translation}节\\
    协调话轮、打断与业务状态 & 用户停顿时系统抢答；播报中途被打断 & 第\ref{sec:nlp-speech-realtime-dialogue}节\\
    \bottomrule
  \end{tabularx}
  \caption{本章的挑战清单。同一句日期询问在不同任务中暴露不同困难；最后一列给出主要讨论的位置，前两项挑战贯穿识别、合成与对话各节。}
  \label{tab:nlp-speech-challenges}
\end{table*}

本章的声音情境、概率、时长和运行记录均为教学构造，不表示实际录音、模型训练或性能测量。共同材料仍是2025年4月15日发布的机器学习读书会通知：活动\code{ML-0418}原定4月17日，改至4月18日（周五）14:00，地点为海棠楼201；校外来宾须在4月16日18:00前提交报名信息。这份通知只在数据、目标构造和内容核验处回用；讲结构与算法时，改用张量形状和小型数值算例。

声音任务还要把录音中的时间与程序运行的时间分开。设录音中某个单位的结束位置为$a_i$，在线系统开始接收这段录音的墙钟时刻为$t_0$；实时匀速送入时，该单位最早完整可见的时刻为$t_0+a_i$。文本单位首次不可撤回地提交于$c_i$，则其提交延迟为$c_i-(t_0+a_i)$。整轮用户结束于$t_{\mathrm{end}}$，最终转写、首段回复开始播放和业务完成分别发生于$t_{\mathrm{asr}}$、$t_{\mathrm{play}}$、$t_{\mathrm{done}}$。录音时间轴、稳定提交时刻与墙钟播放时刻可以相差很远；离线文件的快速解码也不能当作实时匀速输入的结果。表\ref{tab:nlp-speech-contracts}按这一约定列出各任务的产出、误差与计时起止点。

\begin{table*}[htbp]
  \centering
  \begin{tabularx}{\linewidth}{@{}p{15mm}p{28mm}X X p{36mm}@{}}
    \toprule
    任务 & 输入与输出 & 正确性及局部错误 & 端到端要求 & 时间起止点\\
    \midrule
    识别 & 声音到文字、时间区间 & 转写编辑距离；另查边界与说话者 & 该说的没有遗漏，未说的不新增 & 单位声学结束到稳定提交；整轮结束到最终转写\\
    口语理解 & 声音到意图、槽位或事实 & 日期、否定、更正及完整记录 & 当前有效要求正确，来源可追踪 & 所需证据可见到语义记录确认\\
    合成 & 文字及声音条件到波形 & 数量、读音、遗漏；自然度另评 & 听者获得规定内容和表达 & 请求到首段可播放；请求到全部声音生成\\
    翻译 & 源声音到目标文字或声音 & 原意、实体与额外新增内容 & 译文正确且实际声音也正确 & 源语证据可见到目标提交；另计首段播放和全段播放结束\\
    实时对话 & 双方声音、历史、工具到交互 & 抢话、漏听、状态错误、重复动作 & 合理轮次、可听回复及业务完成证据 & 用户结束到首段播放；打断到停止；请求到业务完成\\
    \bottomrule
  \end{tabularx}
  \caption{语音任务的产出、误差与时间口径。同一日期询问可以转写正确却回复过早。表中的起止点须由同一时钟或明确换算得到；波形生成结束、播放结束和业务完成不是同一个事件。}
  \label{tab:nlp-speech-contracts}
\end{table*}
\FloatBarrier

各项任务的方法沿一条共同的谱系演进。先是表示：用手工频谱特征或学到的连续表示描述声音，再用神经编解码器把声音变成离散词元。其次是对齐：HMM用隐状态在时间上的停留吸收长度差异；CTC与RNN-T对所有合法对齐路径求和，后者还让每一步读取已输出的前缀；注意力编码器—解码器与Whisper把识别写成以声音为条件的文本生成，由注意力隐式完成对齐。最后，声音词元并入语言模型：合成可以像生成文字一样逐个生成声学词元，翻译可以直接生成目标语言的离散单位，语音大模型与全双工对话模型则在同一个模型中读写文字与声音。合成本身也经历了从多模块流水线到联合训练、再到词元生成的变化，翻译则从级联走向直接转换。

各节的安排沿着这条谱系与声音的流向展开。第\ref{sec:nlp-speech-representation}节建立各任务共用的声音表示：声学特征、卷积下采样与Conformer编码器、自监督表示和离散声学词元。第\ref{sec:nlp-speech-recognition}节以“变长对齐”为主线比较各类识别模型，并讨论流式识别。第\ref{sec:nlp-speech-postprocessing}节把识别器的原始输出整理为可用转写：语言模型融合、强制对齐、说话者区分与反向文本规范化。第\ref{sec:nlp-speech-understanding}节从转写或声音中恢复意图、槽位与交流作用。第\ref{sec:nlp-speech-tts}节反转方向，从文字生成声音；第\ref{sec:nlp-speech-translation}节把识别、翻译与合成连起来，再比较直接转换。第\ref{sec:nlp-speech-realtime-dialogue}节把语音编码器接入大语言模型，并讨论全双工对话中的轮次、打断与行为协调。本章用到的Transformer语言模型见第\ref{chap:nlp-large-language-models}章，束搜索等解码方法见第\ref{chap:nlp-model-inference}章，文本翻译见第\ref{chap:nlp-text-generation}章；各节在首次用到时说明所用结论。

\section{语音信号与表示}
\label{sec:nlp-speech-representation}

识别、合成、翻译与对话都要先回答同一个问题：把波形变成什么样的序列，后续模型才能高效学习。16\,kHz采样的一秒录音有16000个数，相邻采样高度相关，单个采样几乎不携带“这是哪个音”的信息；直接把它交给注意力模型，序列长、冗余大。第二个困难来自声学变化：同一句“四月十八日”，换一位说话者、换一个房间或换一种语速，波形可以完全不同，而带人工转写的录音远少于没有转写的录音。第三个困难来自生成侧：语言模型按离散词元逐步预测，连续的频谱帧既不便直接作为它的输出，也不便低码率传输。

方法谱系沿这三个困难展开：手工短时频谱与卷积下采样缩短序列（第\ref{sec:nlp-speech-features}、\ref{sec:nlp-speech-encoder}小节），自监督预训练用无标注录音弥补转写不足（第\ref{sec:nlp-speech-ssl}小节），神经编解码器把波形变成语言模型可以读写的离散码流（第\ref{sec:nlp-speech-codec}小节）。节末表\ref{tab:nlp-speech-representations}比较这些表示的帧率、形式、监督需求与用途。

\subsection{声学特征与数据增强}
\label{sec:nlp-speech-features}

本小节解决两件事：用短时频谱表达“每一小段时间里有哪些频率成分”，以及在训练数据有限时让模型不过分依赖少数频带或帧。划分训练、开发与测试数据时应按录音或说话者分组，再截段和增强，否则同一录音的相邻切片会同时进入训练集与测试集。

后文会同时出现三种长度：特征帧数$n_{\mathrm{frame}}$（每10\,ms一帧）、编码器输出的编码步数$T$（下采样后），以及文字或子词符号数$U$。三者一般互不相等，识别模型要解决的对齐问题正是$T$步声音与$U$个符号之间的对应。

\subsubsection{log-Mel与MFCC}
\label{sec:nlp-speech-logmel-mfcc}

波形$x_n$记录的是瞬时振幅，直接观察它不容易区分某个音的频率成分；而语音的频率结构又随发音快速变化。短时分析的做法是把信号切成相互重叠的小窗，在每个窗内把信号近似看作平稳，观察局部频谱。设单声道波形采样率为$f_s$，窗长$w$个采样，帧移$h$个采样；统一采样率时应先作抗混叠滤波再重采样。

浊音的频谱能量随频率升高而下降，高频成分容易被低频淹没。\term{预加重}（pre-emphasis）先用一阶差分提升高频：$x'_n=x_n-\alpha x_{n-1}$，常取$\alpha=0.97$。随后对第$i$帧乘以窗函数$g_n$。直接截取一段信号相当于乘矩形窗，窗边缘的突变会把能量泄漏到其他频率；\term{汉明窗}（Hamming window）$g_n=0.54-0.46\cos\bigl(2\pi n/(w-1)\bigr)$让两端平滑衰减，减轻这种泄漏。加窗后补零到变换长度$k_{\mathrm{fft}}\geq w$，计算\term{短时傅里叶变换}（short-time Fourier transform, STFT）与功率谱：
\begin{equation}
 S_{i,k}=\sum_{n=0}^{w-1}x'_{ih+n}g_n
       \exp\!\left(-\mathrm{i}\frac{2\pi kn}{k_{\mathrm{fft}}}\right),
 \qquad P_{i,k}=\frac{|S_{i,k}|^2}{k_{\mathrm{fft}}}.
 \label{eq:nlp-speech-stft}
\end{equation}
这里帧号$i$从0开始，$\mathrm{i}$为虚数单位，第$k$个频率点对应$f_k=kf_s/k_{\mathrm{fft}}$；实信号的频谱共轭对称，只保留$k=0,\ldots,k_{\mathrm{fft}}/2$。取平方丢掉了相位，后文从频谱生成波形的声码器要设法补回它。

人耳对低频的频率差异比对高频敏感。\term{梅尔刻度}（Mel scale）用$\mathrm{mel}(f)=2595\log_{10}(1+f/700)$近似这种感知尺度：低频段近似线性，高频段近似对数。在梅尔刻度上把$[0,f_s/2]$等分，得到$m+2$个边界点，换回赫兹记为$b_0<b_1<\cdots<b_{m+1}$；第$j$个三角滤波器从$b_{j-1}$升到$b_j$再降到$b_{j+1}$：
\begin{equation}
 H_{j,k}=
 \begin{cases}
  \dfrac{f_k-b_{j-1}}{b_j-b_{j-1}}, & b_{j-1}\leq f_k\leq b_j,\\[2mm]
  \dfrac{b_{j+1}-f_k}{b_{j+1}-b_j}, & b_j< f_k\leq b_{j+1},\\[1mm]
  0, & \text{其他}.
 \end{cases}
 \label{eq:nlp-speech-mel-filter}
\end{equation}
功率谱经这组滤波器加权求和，再取对数，就得到\term{对数梅尔滤波器组特征}（log-Mel filterbank features，简称log-Mel）：用少量频带描述每帧的能量分布，
\begin{equation}
 \bm{X}_{i,j}=\log\!\left(\epsilon+
                   \sum_k H_{j,k}P_{i,k}\right),
 \quad j=1,\ldots,m.
 \label{eq:nlp-speech-logmel}
\end{equation}
$\epsilon>0$避免对零取对数。取对数一方面压缩动态范围，另一方面把信道的乘性增益变成加性偏移，便于再按句或按全体训练数据做均值—方差归一化。滤波器定义、对数底数和归一化方式随实现而异，须随模型一起保存。

以16\,kHz、$m=80$为教学例子：$\mathrm{mel}(8000)\approx2840$，82个边界点在梅尔刻度上的间距约35。第一个滤波器覆盖0至45\,Hz，最后一个覆盖约7475至8000\,Hz，宽约525\,Hz。$k_{\mathrm{fft}}=512$时频率点间隔为31.25\,Hz，所以低频滤波器只汇总一两个频率点，高频滤波器汇总十几个：频带越高，分辨越粗，这正是梅尔刻度要表达的感知特性。

再看帧数。给一秒、16\,kHz的教学波形，共16000个采样。取25\,ms窗长$w=400$、10\,ms帧移$h=160$、$k_{\mathrm{fft}}=512$，不做居中补边，也丢弃不足一窗的尾部，则帧数为
\begin{equation}
 n_{\mathrm{frame}}=
 1+\left\lfloor\frac{16000-400}{160}\right\rfloor=98.
 \label{eq:nlp-speech-frame-count}
\end{equation}
因此STFT单边谱为$98\times257$，log-Mel矩阵为$98\times80$。最后一窗从0.970\,s到0.995\,s，最后5\,ms未进入完整窗。改用20\,ms帧移得到49帧；保持10\,ms帧移、改用40\,ms窗，则得到97帧。更长的窗提高频率分辨率，却会混合更长时间内的变化；补零只加密离散频率采样，不增加真实分辨率。

\term{梅尔倒谱系数}（Mel-frequency cepstral coefficients, MFCC）在log-Mel上再做\term{离散余弦变换}（discrete cosine transform, DCT），保留前$q$项：
\begin{equation}
 c_{i,r}=\sum_{j=0}^{m-1}\bm{X}_{i,j+1}\cos\!\left[\frac{\pi r(j+1/2)}{m}\right],
 \qquad r=0,\ldots,q-1,
 \label{eq:nlp-speech-mfcc}
\end{equation}
归一化系数依实现约定。取$q=13$得到$98\times13$。相邻梅尔频带的对数能量高度相关，DCT把它们变换到近似不相关的坐标上，低阶系数描述谱包络，高阶系数描述细节，截断即丢弃细节。这种去相关服务于第\ref{sec:nlp-speech-hmm}小节用对角协方差高斯混合建模特征的HMM系统。传统系统还拼接一阶、二阶差分，$\Delta c_{i,r}=\sum_{\ell=1}^{2}\ell(c_{i+\ell,r}-c_{i-\ell,r})/(2\sum_{\ell=1}^{2}\ell^2)$，得到$13\times3=39$维；第\ref{sec:nlp-speech-ssl}小节的HuBERT第一轮聚类就用这种39维MFCC。神经编码器能自己学习频带之间的相关性，DCT截断反而丢信息，所以现代识别器多直接读取log-Mel。

\subsubsection{SpecAugment}
\label{sec:nlp-speech-specaugment}

识别器训练数据有限时，可能过分依赖某个频带或某几帧的局部细节。SpecAugment直接在训练输入的log-Mel矩阵上做三种变换：时间扭曲、频率遮蔽和时间遮蔽，转写标签保持不变\scite{nlp-speech-park2019}。它改变的是声学观测，不是重新录制一份独立录音；代价几乎只是一次矩阵操作，因此可以在每轮训练中重新采样。

设特征矩阵有$\tau$帧、$\nu$个频带。\term{频率遮蔽}（frequency masking）从$\{0,\ldots,F\}$均匀采样宽度$f$，再从$\{0,\ldots,\nu-f\}$采样起点$j_0$，令所有帧的第$j_0$至$j_0+f-1$列取遮蔽值。\term{时间遮蔽}（time masking）同样采样长度$r\leq R$与起点$i_0$，把第$i_0$至$i_0+r-1$行取遮蔽值；原方案还限制$r\leq p\tau$，避免短句被整段遮住。两种遮蔽可以各做若干次。以零均值化后的$98\times80$特征为例，若$f=8,j_0=20,r=10,i_0=30$，列20至27与行30至39被置零；索引均从0开始。置零在这里等于填入平均特征，取决于之前的归一化，不等于真实静音。

\term{时间扭曲}（time warping）模拟语速的局部变化。在$(W,\tau-W)$内随机选一个锚点帧$a$，再从$[-W,W]$采样位移$\omega$，把锚点移到$a+\omega$：锚点左侧的$a$帧被线性拉伸或压缩为$a+\omega$帧，右侧的$\tau-a$帧变为$\tau-a-\omega$帧，总帧数不变。输出第$i'$帧从原特征的位置
\begin{equation}
 s(i')=
 \begin{cases}
  i'\,\dfrac{a}{a+\omega}, & 0\leq i'\leq a+\omega,\\[2mm]
  a+(i'-a-\omega)\dfrac{\tau-a}{\tau-a-\omega}, & a+\omega< i'<\tau
 \end{cases}
 \label{eq:nlp-speech-time-warp}
\end{equation}
处插值取值。教学例子取$\tau=98$、$a=50$、$\omega=4$：输出前54帧来自原前50帧，输出第27帧取原位置25，后半段被压缩，输出第54帧对应原第50帧。原论文的LibriSpeech基本策略取$W=80$、$F=27$、$R=100$，并在消融中指出时间扭曲计算最贵、收益最小，算力紧张时可以首先去掉。

波形层面的增强作用在另一类对象上：加噪叠加其他声音，混响模拟房间传播后的拖尾，变速（常用0.9、1.0、1.1三档）同时改变时长与音高。它们要重新计算特征，成本更高，却能产生遮蔽无法模拟的声学条件。变速和时间扭曲之后文字可以沿用，时间戳须按同一映射变换。遮蔽过强若删去了日期的全部证据，标签未改，样本却已成为含歧义的条件预测；增强强度应由开发集选择。

\subsection{声学编码器结构}
\label{sec:nlp-speech-encoder}

10\,ms帧移下，一句10秒的话有约1000帧。自注意力对每层每头要计算$1000\times1000$个得分，代价随长度平方增长；语音还同时需要两类信息：几十毫秒内共振峰的过渡等局部频谱模式，以及跨越词和句子的长程上下文。双向LSTM能读整句，但沿时间串行计算，难以并行；纯卷积编码器（如ContextNet）擅长局部模式，需要很深的堆叠或额外的全局汇聚才能覆盖整句\cite{nlp-speech-han2020contextnet}；纯Transformer编码器一步就能读取全局，局部模式却只能靠注意力权重间接学习。本小节的路线是：先用卷积前端把帧率降到40\,ms，再用Conformer块把卷积与自注意力放进同一层\cite{nlp-speech-gulati2020conformer}，它也是第\ref{sec:nlp-speech-recognition}节各类识别器的默认骨干。缩放点积注意力、多头与位置编码沿用第\ref{sec:nlp-llm-decoder-layer}小节的记号，其一般推导可参考\BookChapterRef{chap:transformer}{卷五的Transformer章}；逐点卷积与深度可分离卷积的定义可参考\BookChapterRef{chap:convolutional-networks}{卷五的卷积神经网络章}。下文只写语音编码器特有的组合方式与张量形状。

\subsubsection{卷积下采样前端}
\label{sec:nlp-speech-subsampling}

\term{卷积下采样}（convolutional subsampling）前端把log-Mel矩阵看成单通道图像，用两层步长为2的二维卷积同时缩短时间轴与频率轴。一层卷积的输出长度为
\begin{equation}
 L_{\mathrm{out}}=\left\lfloor\frac{L_{\mathrm{in}}+2p-k}{s}\right\rfloor+1,
 \label{eq:nlp-speech-conv-length}
\end{equation}
其中$k$为核长、$s$为步长、$p$为单侧补边。本章约定核为$3\times3$、$s=2$、$p=1$，此时$L_{\mathrm{out}}=\lceil L_{\mathrm{in}}/2\rceil$。

沿用98帧的例子，取通道数$c$与模型宽度$d$都为144，各层形状为
\[
 \begin{aligned}
 &1\times98\times80
 \xrightarrow{\text{卷积}+\mathrm{ReLU}}
 144\times49\times40
 \xrightarrow{\text{卷积}+\mathrm{ReLU}}
 144\times25\times20\\
 &\qquad\xrightarrow{\text{重排}}
 25\times2880
 \xrightarrow{\text{线性}}
 25\times144 .
 \end{aligned}
\]
重排把每个编码步的144个通道、20个频带拼成一个2880维向量，线性层再投影到$d$维。于是一秒声音得到$T=25$个编码步，每步40\,ms；不补边的实现（$p=0$）则得到$98\to48\to23$。参数主要在第二层卷积（$9\times144\times144+144=186768$）与线性投影（$2880\times144+144=414864$）上。感受野按$r_2=r_1+(k-1)s_1=3+2\times2=7$计算，即每个编码步读取7个特征帧，对应$6\times10+25=85$\,ms的波形，频率方向覆盖7个梅尔频带。

下采样4倍让后续每层注意力的得分矩阵缩小为原来的$1/16$，也让第\ref{sec:nlp-speech-recognition}节CTC与RNN-T的对齐格点按比例缩小。下采样却不能无限加大。CTC要求编码步数不少于输出符号数（相邻重复符号还需额外的空白步）。英语按字母建模时，正常语速每秒约有十几个字母；若再下采样到80\,ms（每秒12.5步），快语速句子可能没有任何合法对齐。中文按字建模每秒只有四五个字，余量较大；改用子词单位也能缓解这一矛盾。

\subsubsection{Conformer块}
\label{sec:nlp-speech-conformer}

Conformer在Transformer块中加入一个卷积模块，并把前馈层拆成首尾两个半步。设输入$\bm H\in\mathbb R^{T\times d}$，每行是一个编码步的表示。一个Conformer块依次计算
\begin{equation}
 \begin{aligned}
  \tilde{\bm H}&=\bm H+\tfrac12\operatorname{FFN}(\bm H),\\
  \bm H'&=\tilde{\bm H}+\operatorname{MHSA}(\tilde{\bm H}),\\
  \bm H''&=\bm H'+\operatorname{Conv}(\bm H'),\\
  \bm Y&=\operatorname{LN}\bigl(\bm H''+\tfrac12\operatorname{FFN}(\bm H'')\bigr),
 \end{aligned}
 \label{eq:nlp-speech-conformer}
\end{equation}
其中$\operatorname{LN}$是层归一化，每个子层内部都先做层归一化（前置归一化），输出与输入同形，均为$T\times d$。图\ref{fig:nlp-speech-conformer}展开了整个编码器与卷积模块。

\paragraph{半步前馈。}前馈模块逐行计算
$\operatorname{FFN}(\bm H)=\operatorname{Dropout}\bigl(\operatorname{Swish}(\operatorname{LN}(\bm H)\bm W_1+\bm b_1)\bm W_2+\bm b_2\bigr)$，
$\bm W_1\in\mathbb R^{d\times4d}$，$\bm W_2\in\mathbb R^{4d\times d}$；\term{Swish函数}$\operatorname{Swish}(x)=x\,\sigma(x)$即第\ref{sec:nlp-llm-swiglu}小节SwiGLU中的SiLU\cite{ramachandran2017swish}。普通Transformer块是“注意力—前馈”两个子层，Conformer把一个前馈换成两个权重为$\tfrac12$的前馈，夹住注意力与卷积。这一\term{马卡龙式结构}（Macaron-style structure）来自把Transformer层看作微分方程数值解的分析：一步前馈拆成首尾两个半步，类似数值积分中把一个算子拆成两个半步夹住另一个算子的分裂格式\cite{nlp-speech-lu2019macaron}。Conformer原文在参数量不变的条件下把这对半步前馈换回单个前馈，LibriSpeech test-other的词错误率（不用外部语言模型）从4.3\%升到4.5\%。

\paragraph{相对位置多头自注意力。}语音句子长短差异很大，绝对位置编码在训练中没见过的长度上容易失效。Conformer沿用Transformer-XL的相对位置编码\cite{nlp-speech-dai2019transformerxl}。对某一头，记$\bm q_i,\bm k_j\in\mathbb R^{d_k}$为第$i$步的查询与第$j$步的键（由$\operatorname{LN}(\tilde{\bm H})$的对应行线性投影得到），$\bm p_{i-j}\in\mathbb R^{d}$为距离$i-j$的正弦编码（不含可学参数），注意力得分为
\begin{equation}
 s_{i,j}=\frac{(\bm q_i+\bm u)^{\top}\bm k_j+(\bm q_i+\bm v)^{\top}\bm W_R\,\bm p_{i-j}}{\sqrt{d_k}},
 \label{eq:nlp-speech-relpos}
\end{equation}
$\bm W_R\in\mathbb R^{d_k\times d}$把距离编码投影到键空间，$\bm u,\bm v\in\mathbb R^{d_k}$是各头可学习的全局偏置。展开后四项分别是内容对内容、内容对距离，以及与查询无关的内容偏好$\bm u^{\top}\bm k_j$和距离偏好$\bm v^{\top}\bm W_R\bm p_{i-j}$。与为每个截断距离学习一个向量的相对位置编码相比，这里距离先用正弦函数编码再投影，任意距离都有定义。其余计算与标准多头注意力相同。

\paragraph{卷积模块。}注意力负责全局依赖，卷积模块负责局部模式。记$\bm Z=\operatorname{LN}(\bm H')\in\mathbb R^{T\times d}$，模块依次计算
\begin{equation}
 \begin{aligned}
  \bm A&=\bm Z\bm W_{\mathrm{p1}}\in\mathbb R^{T\times2d},\qquad
  \bm G=\bm A_{:,1:d}\odot\sigma(\bm A_{:,d+1:2d})\in\mathbb R^{T\times d},\\
  \bm D_{t,c}&=\sum_{\ell=-(K-1)/2}^{(K-1)/2}w_{\ell,c}\,\bm G_{t+\ell,c},\\
  \operatorname{Conv}(\bm H')&=\operatorname{Dropout}\bigl(\operatorname{Swish}(\operatorname{BN}(\bm D))\,\bm W_{\mathrm{p2}}\bigr).
 \end{aligned}
 \label{eq:nlp-speech-conv-module}
\end{equation}
第一步是逐点卷积，即对每个编码步做同一个线性变换，把$d$维扩到$2d$维；随后的门控线性单元（GLU，与第\ref{sec:nlp-llm-swiglu}小节同一类门控）用后一半经sigmoid后的值逐元素调制前一半，回到$d$维。第二步是\term{深度卷积}（depthwise convolution）：每个通道$c$用自己的一维核$w_{:,c}$沿时间卷积，通道之间不混合，参数只有$Kd$个；越界位置补零。之后是批归一化$\operatorname{BN}$、Swish激活和第二个逐点卷积$\bm W_{\mathrm{p2}}\in\mathbb R^{d\times d}$，最后一步负责跨通道混合。原文的消融选了核长32；开源实现多取奇数（如31），使两侧补边对称。核长31、每步40\,ms时，深度卷积的局部窗口约为1.24秒，已经覆盖几个音节。

\paragraph{形状、参数与计算。}取$T=25$、$d=144$，一个块内的形状变化为：半步前馈$25\times144\to25\times576\to25\times144$；注意力的得分矩阵每头$25\times25$；卷积模块$25\times144\to25\times288\to25\times144$，此后保持$25\times144$。忽略偏置与归一化参数，两个前馈各$8d^2$，注意力的四个投影加$\bm W_R$共$5d^2$，卷积模块两个逐点卷积共$3d^2$、深度卷积$Kd$，合计约$24d^2+Kd$；$d=144$、$K=31$时约50万个参数，16块约800万个。普通Transformer块约$12d^2$，Conformer块的参数大约是它的两倍。每层计算中，注意力为$O(T^2d)$，卷积模块为$O(TKd+Td^2)$，前馈为$O(Td^2)$；下采样后的$T$较小，前馈与逐点卷积往往占主要部分。上述$d=144$、16块的配置对应原文的小型模型；原文大型模型为17块、$d=512$、约1.2亿参数。

\begin{figure*}[htbp]
  \centering
  % 第九卷 15.1.2 节：卷积下采样前端与Conformer块（跨栏，约140 mm宽）。
% 蓝色：输入特征；白底黄描边：模型部件；浅黄底：卷积模块；红色：编码输出。
\begin{tikzpicture}[
  x=1mm,y=1mm,
  figure picture,
  inp/.style={figure node,draw=FigureInput,fill=FigureInputFill,
    minimum height=7mm,minimum width=46mm,inner xsep=2pt,inner ysep=2pt,
    font=\figurefont\scriptsize},
  mod/.style={figure node,draw=FigureModel,fill=FigureModelFill,
    minimum height=7mm,minimum width=46mm,inner xsep=2pt,inner ysep=2pt,
    font=\figurefont\scriptsize},
  hot/.style={mod,fill=FigureCream},
  sub/.style={figure node,draw=FigureModel,fill=FigureModelFill,
    minimum height=6.5mm,minimum width=40mm,inner xsep=2pt,inner ysep=2pt,
    font=\figurefont\scriptsize},
  res/.style={figure node,draw=FigureOutput,fill=FigureOutputFill,
    minimum height=7mm,minimum width=46mm,inner xsep=2pt,inner ysep=2pt,
    font=\figurefont\scriptsize},
  plus/.style={circle,draw=FigureStroke,fill=white,line width=.6pt,
    inner sep=0pt,minimum size=3.6mm,font=\figurefont\scriptsize},
  note/.style={font=\figurefont\scriptsize,text=FigureMuted,anchor=west},
  fl/.style={figure flow},
  skip/.style={figure flow,draw=FigureMuted,line width=.6pt},
  guide/.style={figure link,dash pattern=on 1.5pt off 1.5pt}
]
  % ---------- 前端 ----------
  \node[inp] (x) at (0,0) {log-Mel特征$\bm X$};
  \node[note] at (25,0) {$T_0\times80$};
  \node[mod] (cv) at (0,12) {两层$3\times3$卷积，步长2，ReLU};
  \node[note] at (25,12) {$c\times\lceil T_0/4\rceil\times20$};
  \node[mod] (ln0) at (0,24) {重排、线性投影、Dropout};
  \node[note] at (25,24) {$T\times d$};
  \draw[fl] (x) -- (cv);
  \draw[fl] (cv) -- (ln0);

  % ---------- Conformer块 ----------
  \draw[figure subgroup] (-33,31) rectangle (24,111);
  \node[font=\figurefont\scriptsize,text=FigureMuted,anchor=north west] at (-32.5,110.5) {$\times N$};

  \node[mod] (f1) at (0,39) {前馈模块（$d\to4d\to d$，Swish）};
  \node[plus] (p1) at (0,48) {$+$};
  \node[mod] (at) at (0,56) {多头自注意力（相对位置编码）};
  \node[plus] (p2) at (0,65) {$+$};
  \node[hot] (cm) at (0,73) {卷积模块};
  \node[plus] (p3) at (0,82) {$+$};
  \node[mod] (f2) at (0,90) {前馈模块};
  \node[plus] (p4) at (0,99) {$+$};
  \node[mod] (ln) at (0,106) {层归一化};
  \node[res] (y) at (0,120) {编码输出$\bm H$};
  \node[note] at (25,120) {$T\times d$，$T\approx T_0/4$};

  \draw[fl] (ln0) -- (f1);
  \draw[fl] (f1) -- node[right,font=\figurefont\scriptsize,text=FigureMuted]{$\times\frac12$} (p1);
  \draw[fl] (p1) -- (at);
  \draw[fl] (at) -- (p2);
  \draw[fl] (p2) -- (cm);
  \draw[fl] (cm) -- (p3);
  \draw[fl] (p3) -- (f2);
  \draw[fl] (f2) -- node[right,font=\figurefont\scriptsize,text=FigureMuted]{$\times\frac12$} (p4);
  \draw[fl] (p4) -- (ln);
  \draw[fl] (ln) -- (y);

  % 残差支路
  \foreach \a/\b in {33.5/p1,51.5/p2,68.5/p3,85.5/p4}{
    \draw[skip] (0,\a) -- (-28,\a) |- (\b.west);
  }

  % ---------- 卷积模块展开 ----------
  \draw[figure subgroup] (52,36) rectangle (128,112);
  \node[font=\figurefont\scriptsize,text=FigureMuted,anchor=north west] at (52.5,111.5) {卷积模块展开};
  \node[sub] (c1) at (75,42) {层归一化};
  \node[sub] (c2) at (75,52) {逐点卷积$d\to2d$};
  \node[sub] (c3) at (75,62) {GLU门控$2d\to d$};
  \node[sub] (c4) at (75,72) {深度卷积，核长$K$};
  \node[sub] (c5) at (75,82) {批归一化};
  \node[sub] (c6) at (75,92) {Swish};
  \node[sub] (c7) at (75,102) {逐点卷积$d\to d$，Dropout};
  \foreach \i/\j in {c1/c2,c2/c3,c3/c4,c4/c5,c5/c6,c6/c7}{\draw[fl] (\i) -- (\j);}
  \node[note] at (96,42) {$T\times d$};
  \node[note] at (96,52) {$T\times2d$};
  \node[note] at (96,62) {$T\times d$};
  \node[note] at (96,72) {逐通道沿时间滤波};
  \node[note] at (96,82) {批内统计量};
  \node[note] at (96,92) {$x\,\sigma(x)$};
  \node[note] at (96,102) {$T\times d$};

  \draw[guide] (cm.east) -- (52,36);
  \draw[guide] (cm.east) -- (52,112);
\end{tikzpicture}

  \caption{卷积下采样前端与Conformer编码器。左列为主链路：$T_0$帧80维log-Mel经两层步长2的卷积降为$T\approx T_0/4$步、$d$维，再堆叠$N$个Conformer块；每个块按“半步前馈—相对位置多头自注意力—卷积模块—半步前馈—层归一化”的次序计算，各子层带残差（圆圈加号），两个前馈模块的输出乘以$\frac12$。右列展开卷积模块：逐点卷积把通道扩为$2d$，GLU门控回到$d$，深度卷积沿时间逐通道滤波，经批归一化与Swish后再由逐点卷积混合通道。右侧标注各步的张量形状。}
  \label{fig:nlp-speech-conformer}
\end{figure*}

Conformer的代价主要落在流式与小批量场景：深度卷积读取后方$(K-1)/2$步，注意力默认读整句，批归一化依赖批内统计量，三者都要改造才能流式运行（第\ref{sec:nlp-speech-chunk-attention}小节）。后续结构沿“降低计算”的方向继续演进，例如Zipformer让中间的若干段编码器在更低帧率上运行，再上采样合并\cite{nlp-speech-yao2024zipformer}。表\ref{tab:nlp-speech-encoders}汇总几类编码器的取舍。

\begin{table*}[htbp]
  \centering
  \small
  \begin{tabularx}{\linewidth}{@{}p{22mm}>{\raggedright\arraybackslash}X>{\raggedright\arraybackslash}X>{\raggedright\arraybackslash}X>{\raggedright\arraybackslash}X@{}}
    \toprule
    编码器 & 局部模式 & 全局上下文 & 并行与计算 & 主要代价\\
    \midrule
    双向LSTM & 由循环状态隐式累积 & 读整句，但远距离信息要逐步传递 & 沿时间串行 & 训练慢；流式时只能用单向\\
    纯卷积（ContextNet类） & 卷积核直接建模 & 靠深层堆叠或全局汇聚 & 全并行，$O(TKd)$ & 全局依赖弱；感受野须设计\\
    Transformer & 只能由注意力权重间接学习 & 一层即可读整句 & 全并行，$O(T^2d)$ & 局部模式学得慢；长序列显存大\\
    Conformer & 深度卷积模块 & 相对位置自注意力 & 全并行，两者之和 & 参数约为Transformer块两倍；批归一化不利于流式\\
    Zipformer & 卷积与注意力 & 多帧率分段 & 中间段降帧率 & 结构与训练配方更复杂\\
    \bottomrule
  \end{tabularx}
  \caption{语音编码器的比较。“局部模式”指几十至几百毫秒内的频谱变化，“全局上下文”指跨词与跨句的依赖；复杂度中$T$为编码步数，$d$为宽度，$K$为卷积核长。}
  \label{tab:nlp-speech-encoders}
\end{table*}

\subsection{自监督语音表示}
\label{sec:nlp-speech-ssl}

带人工转写的录音昂贵：LibriSpeech约有960小时带转写的有声书录音，而Libri-Light提供约6万小时同类无转写录音。若能只用录音学到可迁移的帧表示，再用少量转写训练识别头，就能缓解配对数据不足。这属于\term{自监督学习}（self-supervised learning）：从数据自身构造预测任务，不依赖人工标签；一般框架可参考\BookChapterRef{chap:machine-learning-paradigms}{卷三的机器学习范式章}。

文本上的遮蔽语言模型（第\ref{sec:nlp-classical-bert}小节）遮住部分词元，再从上下文预测它们。搬到语音上有两处不同：语音没有现成的离散词表，预测目标须先构造；相邻帧几乎相同，只遮一帧时模型靠左右插值就能恢复，所以要遮住连续的若干帧。下面三种方法都遮住若干段编码步、由上下文恢复，区别在恢复目标：wav2vec 2.0在同时学习的量化码本上做对比判别\cite{nlp-speech-baevski2020wav2vec2}；HuBERT先离线聚类得到离散伪标签，再做遮蔽分类\cite{nlp-speech-hsu2021hubert}；WavLM在HuBERT的基础上混入干扰声音并改进位置偏置\cite{nlp-speech-chen2022wavlm}。更早的对比预测编码（CPC）用InfoNCE预测未来若干步的潜表示\cite{nlp-matching-oord2018cpc}，wav2vec 2.0把“预测未来”改成了“恢复被遮住的部分”。

\subsubsection{wav2vec 2.0遮蔽对比学习}
\label{sec:nlp-speech-wav2vec}

wav2vec 2.0由三部分组成：从波形提取潜表示的卷积特征编码器、读取遮蔽后序列的Transformer上下文网络，以及把未遮蔽的潜表示变成离散目标的量化模块，如图\ref{fig:nlp-speech-wav2vec}所示。

\paragraph{特征编码器。}\term{特征编码器}（feature encoder）直接读取16\,kHz波形，由7层一维卷积组成，每层512个通道，核长依次为$(10,3,3,3,3,2,2)$，步长依次为$(5,2,2,2,2,2,2)$，每层后接归一化与GELU。总步长为$5\times2^6=320$个采样，即每20\,ms输出一步；感受野按$1+\sum_l(k_l-1)\prod_{m<l}s_m=1+9+10+20+40+80+80+160=400$个采样，即25\,ms，与手工特征的窗长相当。一秒波形依次变为$3199,1599,799,399,199,99,49$步，输出$\bm Z=(\bm z_1,\ldots,\bm z_{49})$，每步512维。这一前端代替了STFT与梅尔滤波器，滤波方式由数据学出。

\paragraph{遮蔽。}从全部编码步中以概率$p=0.065$独立选出起点，每个起点连同其后共$M=10$步被遮蔽，相邻跨度可以重叠。某一步不被任何跨度覆盖的概率为$(1-p)^M=0.935^{10}\approx0.51$，所以约49\%的编码步被遮蔽，平均跨度约15步、300\,ms左右。被遮蔽位置的$\bm z_t$替换为同一个可学习的向量，再送入上下文网络；量化模块读取的是遮蔽前的$\bm Z$。

\paragraph{上下文网络。}Transformer读取遮蔽后的序列，输出上下文表示$\bm c_t$。基础版为12层、宽768、8头，大型版为24层、宽1024、16头。位置信息不用绝对编码，而由一层核长128的分组卷积从输入序列算出，加到输入上，起到相对位置编码的作用。

\paragraph{乘积量化与Gumbel-softmax。}目标$\bm q_t$从未遮蔽的$\bm z_t$量化得到。\term{乘积量化}（product quantization）准备$G$个码本，每个码本有$V$个码字$\bm e_{g,v}$；每一步从每个码本各选一个码字，拼接后经线性变换得到$\bm q_t=\bm W_q[\bm e_{1,i_1};\ldots;\bm e_{G,i_G}]$。原文取$G=2$、$V=320$，只存640个码字，却能组合出$320^2=102400$个目标。选码字是离散操作，无法求导。设线性层由$\bm z_t$给出第$g$组对第$v$个码字的对数几率$l_{g,v}$。Gumbel-max技巧指出：给每个$l_{g,v}$加独立的Gumbel噪声$n_v=-\log(-\log u_v)$，$u_v\sim\mathrm U(0,1)$，再取最大者，恰好等价于按$\operatorname{softmax}(l_{g,:})$抽样。Gumbel-softmax把这个取最大换成带温度$\tau$的softmax\cite{nlp-speech-jang2017gumbel}：
\begin{equation}
 p_{g,v}=\frac{\exp\bigl((l_{g,v}+n_v)/\tau\bigr)}{\sum_{k=1}^{V}\exp\bigl((l_{g,k}+n_k)/\tau\bigr)}.
 \label{eq:nlp-speech-gumbel}
\end{equation}
前向计算取$i_g=\argmax_v p_{g,v}$的硬选择，反向传播则沿软概率$p_{g,v}$求梯度，这种前向离散、反向连续的做法称为\term{直通估计器}（straight-through estimator）。温度从2按每步乘以0.999995退火，基础版下限为0.5、大型版为0.1：早期分布较平，各码字都能得到梯度；后期接近硬选择，训练与使用一致。

\paragraph{对比损失与多样性损失。}对每个被遮蔽的步$t$，模型要从一组候选$\mathcal Q_t$中认出本步的量化向量$\bm q_t$。候选还包括从同一句其他被遮蔽步均匀抽取的$K$个干扰项。记$\operatorname{sim}(\bm a,\bm b)=\bm a^{\top}\bm b/(\lVert\bm a\rVert\lVert\bm b\rVert)$为余弦相似度，$\kappa$为温度，$\bm c_t$先经线性层投影到$\bm q_t$的维数，\term{对比损失}（contrastive loss）为
\begin{equation}
 L_m=-\log\frac{\exp\bigl(\operatorname{sim}(\bm c_t,\bm q_t)/\kappa\bigr)}
 {\sum_{\tilde{\bm q}\in\mathcal Q_t}\exp\bigl(\operatorname{sim}(\bm c_t,\tilde{\bm q})/\kappa\bigr)},
 \qquad |\mathcal Q_t|=K+1.
 \label{eq:nlp-speech-wav2vec}
\end{equation}
原文取$K=100$、$\kappa=0.1$。这正是第\ref{chap:nlp-classification-matching-organization}章式\eqref{eq:nlp-matching-infonce}的InfoNCE损失：查询是遮蔽位置的上下文输出，正例是同一位置的量化目标，负例是同一句其他位置的量化目标。负例取自同一句，意味着正负例共享说话者与信道条件，模型无法靠音色区分，只能依据该位置的内容。

与文本匹配不同，这里的正负例本身也是学出来的：量化模块若只使用少数几个码字，许多位置的目标相同，对比任务就失去区分力。\term{多样性损失}（diversity loss）鼓励各码字被均匀使用。记$\bar p_{g,v}$为一批句子中各步不加Gumbel噪声的softmax概率的平均，
\begin{equation}
 L_d=\frac{1}{GV}\sum_{g=1}^{G}\sum_{v=1}^{V}\bar p_{g,v}\log\bar p_{g,v}
    =-\frac{1}{GV}\sum_{g=1}^{G}H(\bar{\bm p}_g),
 \label{eq:nlp-speech-diversity}
\end{equation}
最小化$L_d$即最大化平均分布的熵$H$，在全部码字均匀使用时取最小值。总目标为$L=L_m+\alpha L_d$，原文取$\alpha=0.1$；在较小的LibriSpeech预训练集上，另对特征编码器输出加L2惩罚。原文消融比较了输入与目标是否量化，连续输入配离散目标效果最好；原文的解释是连续输入保留更多信息，便于上下文网络建模，而离散目标使训练更稳健。算法\ref{alg:nlp-speech-wav2vec}汇总一步预训练。

\begin{algorithm}[wav2vec 2.0的一步预训练]
  \label{alg:nlp-speech-wav2vec}
  \begin{algorithmic}[1]
    \Require 一批波形，特征编码器$f$，上下文网络$g$，量化模块，遮蔽参数$p,M$，干扰项数$K$，温度$\kappa,\tau$及下限$\tau_{\min}$
    \Ensure 更新后的参数
    \For{批中每条波形$\bm x$}
      \State $\bm Z=f(\bm x)$，$\bm Z\in\mathbb R^{T\times512}$ \Comment{每20\,ms一步}
      \State 对每步按式\eqref{eq:nlp-speech-gumbel}选码字，得到量化目标$\bm q_1,\ldots,\bm q_T$ \Comment{读取未遮蔽的$\bm Z$}
      \State 以概率$p$抽起点，每个起点遮蔽$M$步，得到遮蔽集合$\mathcal M$；把$\bm z_t,t\in\mathcal M$换成可学习向量，得$\tilde{\bm Z}$
      \State $\bm C=g(\tilde{\bm Z})$
      \For{$t\in\mathcal M$}
        \State 从$\mathcal M\setminus\{t\}$均匀抽$K$个位置，取其量化向量作干扰项
        \State 按式\eqref{eq:nlp-speech-wav2vec}计算$L_m^{(t)}$
      \EndFor
    \EndFor
    \State 按式\eqref{eq:nlp-speech-diversity}在全批计算$L_d$；$L\gets\operatorname{mean}_t L_m^{(t)}+\alpha L_d$
    \State 反向传播（量化选择用直通估计器），更新$f$、$g$、码本与投影层；$\tau\gets\max(\tau_{\min},\,0.999995\,\tau)$
  \end{algorithmic}
\end{algorithm}

\begin{figure*}[htbp]
  \centering
  % 第九卷 15.1.3 节：wav2vec 2.0 预训练结构（跨栏，约145 mm宽）。
% 蓝色：输入与潜表示；灰色：被遮蔽的步；白底黄描边：模型部件；红色：损失。
\begin{tikzpicture}[
  x=1mm,y=1mm,
  figure picture,
  inp/.style={figure node,draw=FigureInput,fill=FigureInputFill,
    minimum height=7mm,minimum width=90mm,inner xsep=2pt,inner ysep=2pt,
    font=\figurefont\scriptsize},
  mod/.style={figure node,draw=FigureModel,fill=FigureModelFill,
    minimum height=7mm,minimum width=90mm,inner xsep=2pt,inner ysep=2pt,
    font=\figurefont\scriptsize},
  qmod/.style={figure node,draw=FigureModel,fill=FigureModelFill,
    minimum height=11mm,minimum width=40mm,inner xsep=2pt,inner ysep=2pt,
    font=\figurefont\scriptsize},
  cell/.style={rectangle,rounded corners=.6mm,draw=FigureInput,
    fill=FigureInputFill,line width=.6pt,minimum size=6.5mm,inner sep=0pt,
    font=\figurefont\scriptsize},
  mcell/.style={cell,draw=FigureMuted,fill=FigureGrid},
  ccell/.style={cell,draw=FigureModel,fill=FigureCream},
  loss/.style={figure node,draw=FigureOutput,fill=FigureOutputFill,
    minimum height=8mm,minimum width=74mm,inner xsep=2pt,inner ysep=2pt,
    font=\figurefont\scriptsize},
  note/.style={font=\figurefont\scriptsize,text=FigureMuted,align=center},
  fl/.style={figure flow},
  thin/.style={figure flow,line width=.6pt}
]
  \node[inp] (x) at (43,0) {波形$\bm x$（16\,kHz）};
  \node[mod] (f) at (43,12) {卷积特征编码器$f$：7层一维卷积，总步长320（20\,ms）};
  \draw[fl] (x) -- (f);

  % 潜表示 z：第3--5步被遮蔽（从0计）
  \foreach \i in {0,...,7}{
    \pgfmathtruncatemacro{\t}{\i+1}
    \ifnum\i>2 \ifnum\i<6
      \node[mcell] (z\i) at (4+11*\i,25) {$\bm m$};
    \else
      \node[cell] (z\i) at (4+11*\i,25) {$\bm z_{\t}$};
    \fi\else
      \node[cell] (z\i) at (4+11*\i,25) {$\bm z_{\t}$};
    \fi
    \draw[thin] (4+11*\i,15.5) -- (z\i.south);
  }
  \node[note,anchor=east] at (-2,25) {遮蔽后};

  \node[mod] (g) at (43,39) {Transformer上下文网络$g$（卷积相对位置编码）};
  \foreach \i in {0,...,7}{\draw[thin] (z\i.north) -- (4+11*\i,35.5);}

  \foreach \i in {0,...,7}{
    \pgfmathtruncatemacro{\t}{\i+1}
    \node[ccell] (c\i) at (4+11*\i,53) {$\bm c_{\t}$};
    \draw[thin] (4+11*\i,42.5) -- (c\i.south);
  }

  % 量化支路：读取遮蔽前的 Z
  \node[qmod] (q) at (122,25) {乘积量化：$G=2$个码本\\各$V=320$个码字，Gumbel-softmax};
  \draw[fl] (f.east) -| node[pos=.25,above,note]{遮蔽前的$\bm Z$} (q.south);
  \node[cell] (qt) at (110,45) {$\bm q_{5}$};
  \node[cell] (qa) at (122,45) {$\tilde{\bm q}$};
  \node[cell] (qb) at (134,45) {$\tilde{\bm q}$};
  \foreach \n in {qt,qa,qb}{\draw[thin] (q.north) -- (\n.south);}
  \node[note] at (122,53) {本位置目标与$K$个干扰项};

  % 损失
  \node[loss] (l) at (78,68) {对比损失$L_m$：$\operatorname{sim}(\bm c_5,\bm q_5)$高于干扰项；总目标$L_m+\alpha L_d$};
  \draw[fl] (c4.north) -- (c4.north |- l.south);
  \draw[fl] (qt.north) -- (qt.north |- l.south);
  \draw[figure link,dash pattern=on 1.5pt off 1.5pt] (q.east) -- ++(4,0) |- (l.east)
    node[pos=.25,right,note]{多样性\\损失$L_d$};
\end{tikzpicture}

  \caption{wav2vec 2.0的预训练结构。卷积特征编码器把波形变成每20\,ms一步的潜表示$\bm z_t$；一路经乘积量化（两个码本各选一个码字后拼接投影）得到目标$\bm q_t$，另一路把被遮蔽的步（灰色）换成同一个可学习向量后送入Transformer，得到$\bm c_t$。对比损失要求遮蔽位置的$\bm c_t$与本位置的$\bm q_t$最相似，干扰项取自同一句其他遮蔽位置；多样性损失作用于码字的平均使用分布。}
  \label{fig:nlp-speech-wav2vec}
\end{figure*}

预训练之后去掉量化模块，在上下文网络顶部加一个随机初始化的线性层，输出字符与词边界符号，用第\ref{sec:nlp-speech-ctc}小节的CTC损失微调。微调时特征编码器通常冻结，输入仍做类似SpecAugment的遮蔽。原文在约5.3万小时无标注录音上预训练后，只用10分钟带转写数据微调（配合语言模型），在LibriSpeech的test-clean与test-other上得到4.8\%与8.2\%的词错误率。代价在于预训练计算量大，Gumbel温度、多样性权重与遮蔽比例都要调；目标与表示同时学习，训练早期目标本身也不稳定。

\subsubsection{HuBERT与WavLM遮蔽预测}
\label{sec:nlp-speech-hubert}

wav2vec 2.0的目标在训练中不断变化，还需要负样本、多样性损失与温度退火来维持。HuBERT（hidden-unit BERT）把两件事拆开：先离线给每个编码步分配一个离散的“隐藏单元”标签，再像BERT那样只在被遮蔽的步上做分类。标签由K均值聚类给出：在一组帧特征上找$C$个聚类中心，使每帧到最近中心的平方距离之和最小，每帧的标签就是最近中心的编号；算法细节可参考\BookChapterRef{chap:clustering}{卷四的聚类章}。

HuBERT的编码器结构与wav2vec 2.0相同：卷积特征编码器加Transformer，遮蔽方式也相同（原文取$p=0.08$、$M=10$）。记$\tilde{\bm X}$为遮蔽后的输入，$\bm o_t$为Transformer在第$t$步的输出，$z_t\in\{1,\ldots,C\}$为该步的聚类标签，$\mathcal M$为遮蔽位置集合。模型为每个类别$c$学一个嵌入$\bm e_c$，并用投影矩阵$\bm A$把$\bm o_t$映到嵌入空间，预测分布为
\[
 p(c\mid\tilde{\bm X},t)=\frac{\exp\bigl(\operatorname{sim}(\bm A\bm o_t,\bm e_c)/\tau\bigr)}{\sum_{c'=1}^{C}\exp\bigl(\operatorname{sim}(\bm A\bm o_t,\bm e_{c'})/\tau\bigr)},
\]
$\operatorname{sim}$为余弦相似度，$\tau=0.1$。损失分为遮蔽位置与未遮蔽位置两部分：
\begin{equation}
 \begin{gathered}
 L=\beta L_{\mathrm m}+(1-\beta)L_{\mathrm u},\\
 L_{\mathrm m}=-\sum_{t\in\mathcal M}\log p(z_t\mid\tilde{\bm X},t),\qquad
 L_{\mathrm u}=-\sum_{t\notin\mathcal M}\log p(z_t\mid\tilde{\bm X},t).
 \end{gathered}
 \label{eq:nlp-speech-hubert}
\end{equation}
原文发现$\beta=1$即只在遮蔽位置计损失效果最好。机制上的理由是：未遮蔽位置的输入就在眼前，模型只需模仿聚类函数“由本帧特征算出最近中心”，用不到上下文；只有被遮住的位置，才迫使模型从周围推断这里应是哪个单位。

对比式\eqref{eq:nlp-speech-wav2vec}，HuBERT的预测分布同样是余弦相似度除以温度后的softmax，区别在候选集合：wav2vec 2.0从同一句抽取$K$个量化向量作负例，HuBERT以全部$C$个类别嵌入为候选，相当于固定类别上的分类交叉熵。与BERT相比，BERT的词表事先给定，HuBERT的词表由聚类构造。

标签的质量来自迭代。第一轮在39维MFCC上做$C=100$的K均值聚类，这些标签与音素的对应很粗；模型学会从上下文预测它们之后，中间层表示比MFCC更接近语音内容。第二轮取第一轮基础模型第6层Transformer的输出重新聚类（$C=500$），再从头训练；大型模型又用第二轮基础模型第9层的特征聚类。这个过程与自训练相似：每轮用当前最好的表示产生目标，再学习预测目标。标签不必正确，只需在相同声音上前后一致，模型就能从大量上下文中平均掉标签噪声。原文还用多组不同$C$的聚类结果同时作为目标，进一步降低对单次聚类的依赖。代价是要训练多轮，每轮要在海量帧特征上聚类，初始特征差时第一轮目标信息量少。

\paragraph{WavLM：去噪遮蔽预测。}HuBERT与wav2vec 2.0主要为识别内容而设计；说话者验证、说话者日志与语音分离则需要保留“谁在说”，并在多人交叠时仍能工作。WavLM沿用HuBERT式的聚类伪标签与遮蔽预测，做了两处修改。第一处是\term{话语混合}（utterance mixing）：训练时以一定概率从同一批中取另一条话语或从噪声库中取一段噪声，随机截取不超过主话语一半长度的片段，叠加到主话语的随机位置；能量比对话语从$[-5,5]$\,dB、对噪声从$[-5,20]$\,dB均匀抽取。预测目标仍是主话语在干净条件下的伪标签。模型因此要在遮蔽之外再去掉干扰，学到的表示须区分主说话者与其他声音。

第二处是\term{门控相对位置偏置}（gated relative position bias）。设$d_{i-j}$为按距离分桶的可学习标量偏置，$\bm q_i$为第$i$步的查询，$\bm a,\bm b\in\mathbb R^{d_k}$与标量$\gamma$可学习，门控由查询内容决定：
\begin{equation}
 g^{\mathrm u}_i=\sigma(\bm q_i^{\top}\bm a),\quad
 g^{\mathrm r}_i=\sigma(\bm q_i^{\top}\bm b),\quad
 r_{i-j}=d_{i-j}+g^{\mathrm u}_i d_{i-j}+(1-g^{\mathrm u}_i)\,\gamma\,g^{\mathrm r}_i\,d_{i-j},
 \label{eq:nlp-speech-wavlm-gate}
\end{equation}
得分为$\bm q_i^{\top}\bm k_j/\sqrt{d_k}+r_{i-j}$。固定的距离偏置对所有位置一视同仁；门控让偏置强度随当前帧内容变化，例如静音帧与语音帧可以对邻近上下文给出不同权重。Base+与大型版在约9.4万小时英语录音上预训练。原文对下游各层加权权重的分析显示，识别、意图分类等内容类任务主要使用较高层；说话者类任务在基础版中主要使用较低层，在大型版中主要使用中间层。

\paragraph{使用方式。}下游使用这些表示时，常冻结预训练模型，像式\eqref{eq:nlp-classical-elmo}的ELMo那样对各层输出学一组softmax权重加权求和，再接轻量的任务头；识别任务也可整体微调。表\ref{tab:nlp-speech-ssl}比较三种方法。

\begin{table*}[htbp]
  \centering
  \small
  \begin{tabularx}{\linewidth}{@{}p{20mm}>{\raggedright\arraybackslash}X>{\raggedright\arraybackslash}X>{\raggedright\arraybackslash}X>{\raggedright\arraybackslash}X@{}}
    \toprule
    方法 & 预测目标与来源 & 损失 & 解决的问题 & 新代价\\
    \midrule
    wav2vec 2.0 & 同一网络在线学习的乘积量化码字 & 对比损失（同句负例）加多样性损失 & 无须转写即可学习帧表示 & 目标与表示同时变化；需负样本与温度退火\\
    HuBERT & 离线K均值聚类标签，逐轮由更好的特征重聚 & 遮蔽位置交叉熵 & 去掉负样本与码本训练，目标稳定 & 多轮预训练；大规模聚类\\
    WavLM & 主话语的聚类伪标签，输入混入干扰 & 遮蔽位置交叉熵 & 兼顾说话者、交叠与噪声条件 & 混合策略与数据规模需设计\\
    \bottomrule
  \end{tabularx}
  \caption{三种自监督语音表示的比较。三者共享卷积特征编码器、Transformer与跨度遮蔽，区别在恢复目标怎样构造以及输入是否加入干扰。}
  \label{tab:nlp-speech-ssl}
\end{table*}

\subsection{神经音频编解码与离散声学词元}
\label{sec:nlp-speech-codec}

连续特征适合作为识别输入，却不便让语言模型直接生成声音：语言模型按离散词元逐步预测，log-Mel既不是离散的，丢掉的相位也要另由声码器补回。传输和存储同样希望用少量比特表示声音。\term{神经音频编解码器}（neural audio codec）用编码器把波形压成低帧率的潜向量，把每个向量量化为若干个码字编号，再由解码器从编号重建波形。这些编号组成的离散序列称为\term{声学词元}（acoustic token），可以像文本词元一样交给语言模型。其核心技术是残差向量量化。

\subsubsection{残差向量量化}
\label{sec:nlp-speech-rvq}

\term{向量量化}（vector quantization, VQ）用一个有限的\term{码本}（codebook）$\mathcal E=\{\bm e_1,\ldots,\bm e_N\}\subset\mathbb R^{D}$近似连续向量：输入$\bm z$被替换为最近的码字，
\[
 k(\bm z)=\argmin_{j}\lVert\bm z-\bm e_j\rVert_2^2,\qquad Q(\bm z)=\bm e_{k(\bm z)},
\]
只需传送编号$k(\bm z)$，每个向量占$\log_2N$比特。码本的学习本身就是在编码器输出上做K均值聚类，只是聚类空间随编码器一起变化。

最近邻选择不可求导。VQ-VAE的做法是：解码器读取$\bm z+\operatorname{sg}(Q(\bm z)-\bm z)$，其中$\operatorname{sg}$表示停止梯度；前向值等于$Q(\bm z)$，反向时梯度原样传给$\bm z$，这与wav2vec 2.0的直通估计器思路相同\cite{nlp-speech-oord2017vqvae}。码本与编码器另由两项损失连接：
\begin{equation}
 L_{\mathrm{vq}}=\bigl\lVert\operatorname{sg}(\bm z)-\bm e_{k(\bm z)}\bigr\rVert_2^2
 +\beta\bigl\lVert\bm z-\operatorname{sg}(\bm e_{k(\bm z)})\bigr\rVert_2^2 .
 \label{eq:nlp-speech-vq}
\end{equation}
第一项把被选中的码字拉向编码器输出；第二项称为\term{承诺损失}（commitment loss），让编码器输出靠近已选码字，避免它在码字之间来回跳动，原文取$\beta=0.25$。码字也可以不用梯度更新，而按指数滑动平均设为分配给它的编码器输出的均值，即在线K均值；长期没有被选中的“死码字”可以用当前批中的随机向量重新初始化。

单码本的问题在于比特数与码本大小呈对数关系。若每帧要用80比特表示，单码本须有$2^{80}$个码字，既存不下也训不动。\term{残差向量量化}（residual vector quantization, RVQ）改为串联$N_q$个各有$N$个码字的小码本，每一级量化上一级剩下的误差：
\begin{equation}
 \bm r^{(1)}=\bm z,\qquad
 \bm r^{(k+1)}=\bm r^{(k)}-Q_k\bigl(\bm r^{(k)}\bigr),\qquad
 \hat{\bm z}=\sum_{k=1}^{N_q}Q_k\bigl(\bm r^{(k)}\bigr),
 \label{eq:nlp-speech-rvq}
\end{equation}
$Q_k$是第$k$级码本上的最近邻量化。每帧输出$N_q$个编号$(i_1,\ldots,i_{N_q})$，共$N_q\log_2N$比特，需要存储的码字只有$N_qN$个。各级的作用由粗到细：第一级给出大致位置，后续各级逐步修正残差。承诺损失对每一级分别计算，$\sum_k\lVert\bm r^{(k)}-Q_k(\bm r^{(k)})\rVert_2^2$。

用二维向量做一个可复核的教学例子。设两级码本与输入分别为
\[
 \begin{gathered}
 \mathcal E_1=\{(0,0),(4,0),(0,4),(4,4)\},\qquad
 \mathcal E_2=\{(0,1),(0,-1),(1,0),(-1,0)\},\\
 \bm z=(4.6,1.3).
 \end{gathered}
\]
第一级到四个码字的平方距离为$22.85,2.05,28.45,7.65$，选$(4,0)$，残差$\bm r^{(2)}=(0.6,1.3)$。第二级到四个码字的平方距离为$0.45,5.65,1.85,4.25$，选$(0,1)$，重建$\hat{\bm z}=(4,1)$，误差$(0.6,0.3)$，平方误差从一级时的2.05降到0.45。两级各2比特，合计4比特，存储8个码字；单码本要达到同样的4比特需要16个码字，比特越多，这一差距按指数扩大。贪心的逐级选择不保证每一级都缩小误差（码本中有零向量时才有此保证），实际训练中后续码本学到的尺度逐级变小，误差随级数下降。

码率按“帧率$\times$码本数$\times\log_2$码本大小”计算。24\,kHz波形经总步长320的编码器后为每秒75帧；码本大小$N=1024$即每个编号10比特；$N_q=8$时码率为$75\times8\times10=6000$比特每秒，即6\,kbps。同样一秒的16比特线性PCM为$24000\times16=384$\,kbps，压缩了64倍。对语言模型而言，6\,kbps意味着每秒600个词元，远多于文字的每秒几个词元；第\ref{sec:nlp-speech-codec-lm-tts}小节的合成模型与第\ref{sec:nlp-speech-moshi}小节的Moshi都要专门处理这种多路码流。

RVQ与wav2vec 2.0用的乘积量化都用几个小码本组合出巨大的码字空间，组合方式不同：乘积量化把向量切成几段或几组，各自独立量化后拼接，各组地位平等；RVQ对整个向量逐级相加，编号有先后与主次，只截取前几级就得到较粗的重建。正是这一性质使下文的可变码率成为可能。

\subsubsection{SoundStream与EnCodec}
\label{sec:nlp-speech-soundstream}

SoundStream与EnCodec是两种代表性的神经编解码器，都采用“卷积编码器—RVQ—镜像解码器”的结构，并用重建与对抗目标联合训练\cite{nlp-speech-zeghidour2022soundstream,nlp-speech-defossez2023encodec}。图\ref{fig:nlp-speech-codec}画出这一结构及训练信号。

\paragraph{编码器与解码器。}以EnCodec的24\,kHz模型为例，编码器先用一层32通道的一维卷积，再接4个卷积块；每块由带残差的卷积单元与一层步长分别为$2,4,5,8$的下采样卷积组成，每次下采样通道数加倍。其后是两层LSTM与一层输出卷积，得到$D=128$维的潜向量。总步长$2\times4\times5\times8=320$，一秒波形$\bm x\in\mathbb R^{1\times24000}$变为$\bm Z\in\mathbb R^{75\times128}$，经RVQ得到编号矩阵$\{1,\ldots,1024\}^{75\times N_q}$，再由镜像结构的转置卷积解码器重建$\hat{\bm x}\in\mathbb R^{1\times24000}$。SoundStream的结构类似，但全部使用因果卷积，可以边录边编码。

\paragraph{训练目标。}只用时域误差训练，重建声音往往发闷，高频细节被平均掉。EnCodec的生成器目标由五项组成。时域重建项为$\ell_t=\lVert\bm x-\hat{\bm x}\rVert_1$。频域重建项在多种窗长上比较梅尔谱，记$S_i(\cdot)$为窗长$2^i$、64个梅尔频带的梅尔谱，
\[
 \ell_f=\frac{1}{|\mathcal I|}\sum_{i\in\mathcal I}\Bigl(\lVert S_i(\bm x)-S_i(\hat{\bm x})\rVert_1+\lVert S_i(\bm x)-S_i(\hat{\bm x})\rVert_2\Bigr),
 \qquad \mathcal I=\{5,\ldots,11\},
\]
短窗约束时间细节，长窗约束频率细节。重建误差仍不足以保证听感自然，于是再加对抗训练：判别器区分真实与重建声音，编解码器则学习使重建难以分辨，生成对抗网络的一般原理可参考\BookChapterRef{chap:neural-generative-models}{卷五的神经生成模型章}。EnCodec使用\term{多尺度短时傅里叶变换判别器}（multi-scale STFT discriminator）：$K=5$个子判别器分别读取窗长为2048、1024、512、256、128的复数STFT（实部与虚部作为两个通道），用二维卷积网络输出真假得分$D_k(\cdot)$。判别器与生成器采用合页形式的对抗损失：
\[
 \begin{gathered}
 L_D=\frac1K\sum_{k=1}^{K}\Bigl[\max\bigl(0,1-D_k(\bm x)\bigr)+\max\bigl(0,1+D_k(\hat{\bm x})\bigr)\Bigr],\\
 \ell_g=\frac1K\sum_{k=1}^{K}\max\bigl(0,1-D_k(\hat{\bm x})\bigr).
 \end{gathered}
\]
\term{特征匹配损失}（feature matching loss）比较判别器各层$l=1,\ldots,L$对真实与重建声音的内部特征，$\ell_{\mathrm{feat}}=\frac{1}{KL}\sum_{k,l}\lVert D_k^l(\bm x)-D_k^l(\hat{\bm x})\rVert_1/\operatorname{mean}\lVert D_k^l(\bm x)\rVert_1$；判别器特征相当于学到的感知度量，这一项也使生成器的训练比单用对抗项稳定。最后是各级RVQ的承诺损失$\ell_w$。生成器总目标为
\begin{equation}
 L_G=\lambda_t\ell_t+\lambda_f\ell_f+\lambda_g\ell_g+\lambda_{\mathrm{feat}}\ell_{\mathrm{feat}}+\lambda_w\ell_w,
 \label{eq:nlp-speech-encodec}
\end{equation}
原文24\,kHz模型取$\lambda_t=0.1$、$\lambda_f=1$、$\lambda_g=3$、$\lambda_{\mathrm{feat}}=3$，并用一个“损失平衡器”按各项对输出$\hat{\bm x}$的梯度范数重新缩放，使这四个权重近似表示各项在总梯度中所占的比例；承诺损失不依赖$\hat{\bm x}$，不进入平衡器。SoundStream的组合相同而形式略异：判别器为多分辨率波形判别器加STFT判别器，频谱项为多尺度梅尔谱距离；两者的码本都按滑动平均更新。训练按算法\ref{alg:nlp-speech-codec}交替更新两方。

\begin{algorithm}[神经音频编解码器的一步训练（EnCodec式）]
  \label{alg:nlp-speech-codec}
  \begin{algorithmic}[1]
    \Require 一批波形$\bm x$，编码器$E$，解码器$G$，$N_q^{\max}$级码本，判别器$\{D_k\}_{k=1}^{K}$
    \Ensure 更新后的编解码器与判别器
    \State $\bm Z\gets E(\bm x)$ \Comment{$75\times128$每秒}
    \State 从$\{1,\ldots,N_q^{\max}\}$均匀抽取本批使用的级数$n_q$ \Comment{量化器丢弃}
    \State 按式\eqref{eq:nlp-speech-rvq}只用前$n_q$级量化，得到$\hat{\bm Z}$与各级残差；前向用$\hat{\bm Z}$，反向直通传给$\bm Z$
    \State $\hat{\bm x}\gets G(\hat{\bm Z})$
    \State 计算$\ell_t,\ell_f,\ell_g,\ell_{\mathrm{feat}},\ell_w$，按式\eqref{eq:nlp-speech-encodec}更新$E$、$G$（码本按梯度或滑动平均更新）
    \State 计算$L_D$，更新判别器 \Comment{$\hat{\bm x}$停止梯度}
  \end{algorithmic}
\end{algorithm}

\paragraph{量化器丢弃与可变码率。}若每个码率都单独训练一个模型，部署要保存多套参数。SoundStream提出\term{量化器丢弃}（quantizer dropout）：每个训练样本从$1$至$N_q^{\max}$中均匀抽取$n_q$，只用前$n_q$级码本量化与重建。由于RVQ的各级按由粗到细排列，模型学会让前几级承载最主要的信息，推理时截取不同级数就得到不同码率。EnCodec的24\,kHz模型据此支持1.5、3、6、12、24\,kbps，分别对应$N_q=2,4,8,16,32$。码率越低，重建越粗，高频与音色细节损失越多。SoundStream原文的主观听测显示，它在3\,kbps下优于12\,kbps的传统编解码器Opus。EnCodec还在码字序列上训练一个小型Transformer语言模型做熵编码，进一步压缩比特数。代价是：对抗训练需要多个判别器与细致调参；码流有$N_q$路，交给语言模型时序列很长；各级码字同时编码内容、音色、韵律与录音条件，彼此纠缠。

\begin{figure*}[htbp]
  \centering
  % 第九卷 15.1.4 节：神经音频编解码器（编码器—RVQ—解码器）与训练信号（跨栏，约145 mm宽）。
% 蓝色：波形；白底黄描边：模型部件；浅黄底：量化器；红色：损失。
\begin{tikzpicture}[
  x=1mm,y=1mm,
  figure picture,
  inp/.style={figure node,draw=FigureInput,fill=FigureInputFill,
    minimum height=8mm,minimum width=16mm,inner xsep=2pt,inner ysep=2pt,
    font=\figurefont\scriptsize},
  mod/.style={figure node,draw=FigureModel,fill=FigureModelFill,
    minimum height=11mm,minimum width=20mm,inner xsep=2pt,inner ysep=2pt,
    font=\figurefont\scriptsize},
  qnt/.style={figure node,draw=FigureModel,fill=FigureCream,
    minimum height=6.5mm,minimum width=20mm,inner xsep=2pt,inner ysep=2pt,
    font=\figurefont\scriptsize},
  sum/.style={circle,draw=FigureStroke,fill=white,line width=.6pt,
    inner sep=0pt,minimum size=5mm,font=\figurefont\scriptsize},
  loss/.style={figure node,draw=FigureOutput,fill=FigureOutputFill,
    minimum height=9mm,minimum width=42mm,inner xsep=2pt,inner ysep=2pt,
    font=\figurefont\scriptsize},
  note/.style={font=\figurefont\scriptsize,text=FigureMuted,align=center},
  fl/.style={figure flow},
  thin/.style={figure flow,line width=.6pt}
]
  \node[inp] (x) at (0,34) {波形$\bm x$\\$1\times24000$};
  \node[mod] (e) at (26,34) {编码器$E$\\卷积，总步长320};
  \draw[fl] (x) -- (e);

  % RVQ
  \draw[figure subgroup] (51,6) rectangle (79,62);
  \node[note] at (65,66) {残差向量量化};
  \node[qnt] (q1) at (65,54) {$Q_1$};
  \node[qnt] (q2) at (65,40) {$Q_2$};
  \node[note] (qd) at (65,28.5) {$\vdots$};
  \node[qnt] (qn) at (65,14) {$Q_{N_q}$};
  \draw[fl] (e.east) -- ++(6,0) |- node[pos=.75,above,note]{$\bm r^{(1)}=\bm z$} (q1.west);
  \draw[thin] (q1) -- node[right,note]{$\bm r^{(2)}$} (q2);
  \draw[thin] (q2) -- (qd);
  \draw[thin] (qd) -- node[right,note]{$\bm r^{(N_q)}$} (qn);

  \node[sum] (s) at (92,34) {$\sum$};
  \draw[thin] (q1.east) -- node[above,note]{$i_1$} (s);
  \draw[thin] (q2.east) -- node[above,note,pos=.4]{$i_2$} (s);
  \draw[thin] (qn.east) -- node[below,note]{$i_{N_q}$} (s);
  \node[note] at (65,0) {每帧$N_q$个编号：声学词元，$75\times N_q$每秒};

  \node[mod] (d) at (112,34) {解码器$G$\\转置卷积};
  \node[inp] (xh) at (136,34) {重建$\hat{\bm x}$\\$1\times24000$};
  \draw[fl] (s) -- node[above,note]{$\hat{\bm z}$} (d);
  \draw[fl] (d) -- (xh);

  % 训练信号
  \node[loss] (rl) at (118,10) {重建损失：时域$\ell_t$\\多尺度梅尔谱$\ell_f$};
  \node[loss] (dl) at (118,58) {多尺度STFT判别器$D_k$\\对抗损失$\ell_g$、特征匹配$\ell_{\mathrm{feat}}$};
  \draw[fl] (xh.south) -- (xh.south |- rl.north);
  \draw[fl] (xh.north) -- (xh.north |- dl.south);
  \draw[thin,draw=FigureMuted] (x.south) -- ++(0,-36) -| ([xshift=-12mm]rl.south);
  \draw[thin,draw=FigureMuted] (x.north) -- ++(0,34) -| ([xshift=-12mm]dl.north);
  \node[note,anchor=west] at (2,-3) {原波形};
\end{tikzpicture}

  \caption{神经音频编解码器的结构与训练信号（以24\,kHz、每秒75帧为例）。编码器以总步长320把波形压为128维潜向量序列；残差向量量化逐级量化上一级残差，每级输出一个码字编号，$N_q$路编号即声学词元；各级码字相加后由镜像解码器重建波形。训练时，重建声音与原声音在时域和多尺度梅尔谱上比较，同时送入多尺度STFT判别器，产生对抗损失与特征匹配损失。}
  \label{fig:nlp-speech-codec}
\end{figure*}

\subsubsection{语义词元与声学词元}
\label{sec:nlp-speech-semantic-acoustic}

离散单位还有另一个来源：对自监督表示做K均值聚类。取HuBERT某一中间层的输出，聚成几百至一千类，每个编码步得到一个编号，这类编号称为\term{语义词元}（semantic token）。它通常每秒25至50个、只有一路，与音素等语言内容高度相关，说话者、韵律与录音条件的信息大多被丢弃。于是语义词元适合让语言模型学习“说了什么”，却不能单独还原原声：从它合成波形需要另训单位声码器，并另外提供音色条件。编解码器的声学词元正相反，保真度高、能还原音色与录音条件，但多路码流中内容与声学细节纠缠，仅在声学词元上训练的语言模型容易生成音质好、语言内容却不连贯的声音。

AudioLM把两者分工使用：先由语言模型生成语义词元，负责长程的内容与结构；再以语义词元为条件生成SoundStream前几级的粗粒度声学词元；最后生成其余各级的细粒度词元，负责音色与音质\cite{nlp-speech-borsos2023audiolm}。另一种折中是让编解码器的第一级码本直接承载语义：SpeechTokenizer在RVQ训练中加入蒸馏损失，使第一级量化输出经投影后与HuBERT表示的余弦相似度尽量高，后续各级则编码剩余的音色与韵律\cite{nlp-speech-zhang2024speechtokenizer}；Moshi所用的Mimi编解码器也把自监督表示蒸馏进第一级码本（第\ref{sec:nlp-speech-moshi}小节）。这样一套码流同时提供“先读第一级得到内容”与“读全部级得到原声”两种用法。

\paragraph{本节对照。}表\ref{tab:nlp-speech-representations}按同一组维度比较本节的各种表示。后文的词元化合成（第\ref{sec:nlp-speech-codec-lm-tts}小节）以声学词元为目标，直接语音翻译（第\ref{sec:nlp-speech-direct-translation}小节）以HuBERT聚类单位为目标，语音大模型（第\ref{sec:nlp-speech-realtime-dialogue}节）同时读写文字与声学词元。这些帧或词元与文字之间怎样对应，留给下一节的识别模型回答。

\begin{table*}[htbp]
  \centering
  \small
  \begin{tabularx}{\linewidth}{@{}p{26mm}p{20mm}p{14mm}>{\raggedright\arraybackslash}X>{\raggedright\arraybackslash}X>{\raggedright\arraybackslash}X@{}}
    \toprule
    表示 & 帧率 & 形式 & 是否需要标注 & 主要用途 & 代价\\
    \midrule
    波形 & 16000或24000\,Hz & 连续 & 否 & 声码器与编解码器的输入输出 & 序列极长，单个采样信息少\\
    log-Mel、MFCC & 100\,Hz & 连续 & 否（手工设计） & 识别输入；合成的中间目标 & 丢弃相位，还原波形需声码器\\
    Conformer输出 & 25\,Hz & 连续 & 是（随识别目标训练） & CTC、RNN-T、注意力识别器的骨干 & 与任务绑定；流式需改造\\
    自监督表示 & 50\,Hz & 连续 & 预训练否，微调需少量 & 低资源识别、说话者等下游任务 & 预训练计算量大\\
    语义词元 & 25至50\,Hz，一路 & 离散 & 否 & 语言模型的内容建模；单位翻译 & 难以还原音色与录音条件\\
    声学词元 & 75\,Hz，$N_q$路 & 离散 & 否 & 词元化合成；语音大模型的输出 & 码流多、序列长；内容与声学纠缠\\
    \bottomrule
  \end{tabularx}
  \caption{本节各种声音表示的比较。帧率按本节例子的配置给出（语义词元依所取的自监督模型而定），实际系统可以不同；“是否需要标注”指学习该表示本身是否需要转写。}
  \label{tab:nlp-speech-representations}
\end{table*}
\FloatBarrier

\section{语音识别}
\label{sec:nlp-speech-recognition}

\term{自动语音识别}（automatic speech recognition, ASR）把一段声音转换为规定形式的文字。沿用第\ref{sec:nlp-speech-representation}节的记号：声学特征矩阵$\bm{X}$经编码器得到$T$个编码步的表示$\bm{h}_1,\ldots,\bm{h}_T$，目标文本为$y=(y_1,\ldots,y_U)$，符号来自词表$V$，可以是字符、子词或词。识别器要给出条件概率$p(y\mid\bm{X})$，并在推断时找出概率高的文本。难点在于训练数据只给出整段声音与整句文字，没有说明哪个字对应哪几帧。本节各方法的根本差别，就在于怎样处理这种未给定的对应关系。

\subsection{识别的挑战与模型谱系}
\label{sec:nlp-speech-asr-challenges}

一个识别样本至少包含声音、采样率、语种和转写规范。转写规范规定输出单位，以及数字、语气词、标点和大小写如何书写；这些规定不改变声音，却决定“正确答案”是什么。语音识别面临的主要挑战可以归为五项。

\begin{enumerate}
  \item 长度不等且没有帧级对齐。第\ref{sec:nlp-speech-features}小节的例子中，一秒声音有98个特征帧，降采样后约25个编码步，对应的中文可能只有四五个字。语速改变时，同一个字占用的帧数随之改变。训练标签通常只有整句文字，模型必须自己决定对应关系，或者把所有可能的对应都考虑进去。
  \item 上下文与流式的矛盾。听到“四月十”时还不能确定后面是“七日”还是“八日”，同音词也常要靠后文区分，所以看见右侧声音能提高准确率。在线系统却必须在用户说完之前陆续交出文字，可等待的右侧声音受延迟预算限制。
  \item 配对标注稀缺。人工转写一小时录音往往需要数倍于录音时长的工作量。网络字幕数量大，但含有错位、删节、意译和非语音文字，属于弱监督。
  \item 噪声、口音与设备差异。背景噪声、混响、口音和麦克风改变了观测分布而不改变文字，在一种条件下训练的模型换到另一种条件时错误率会上升。
  \item 多说话者。会议录音中多人交替甚至同时说话，系统还要回答“谁在何时说话”。这一项在第\ref{sec:nlp-speech-xvector}小节单独处理。
\end{enumerate}

评价识别结果先要比较两个符号序列。设参考序列为$r_1,\ldots,r_N$，识别结果为$\hat r_1,\ldots,\hat r_M$，$d_{i,j}$表示把参考前$i$项变成识别结果前$j$项所需的最少编辑次数，则
\begin{equation}
 d_{i,j}=\min\{d_{i-1,j}+1,\ d_{i,j-1}+1,\
              d_{i-1,j-1}+\mathbf{1}[r_i\ne\hat r_j]\},
 \quad d_{i,0}=i,\ d_{0,j}=j.
 \label{eq:nlp-speech-edit-distance}
\end{equation}
三项依次对应删除、插入与替换或匹配。回溯最短路径得到替换数$N_{\mathrm{S}}$、删除数$N_{\mathrm{D}}$和插入数$N_{\mathrm{I}}$，\term{词错误率}（word error rate, WER）定义为$(N_{\mathrm{S}}+N_{\mathrm{D}}+N_{\mathrm{I}})/N$。中文通常按字计算，称为\term{字符错误率}（character error rate, CER）。插入可以使错误率超过1；参考为空时只能报告插入数。

计分前必须固定规范化口径。参考写“14:00”而识别结果写“下午两点”，按字比较会产生多次编辑，语义却相同；繁简、全半角和标点是否计分也要事先约定。另一方面，平均错误率会掩盖关键字段的失败：参考“四月十八日”被识别为“四月十七日”，CER为1/5；放在一条五十字的通知中只占2\%，通知的核心信息却已经错了。因此，日期、数量和专名这类字段宜在CER之外单独统计。

图\ref{fig:nlp-speech-asr-lineage}给出本节的方法谱系。HMM路线用显式的声学状态描述“在某个音上停留，再前进到下一个音”，需要帧级对齐和多阶段训练。CTC在所有单调对齐上求和，只需整句标签，但假设各步输出条件独立。注意力编码器—解码器读取已输出的前缀，去掉了条件独立假设，却失去了单调性保证，也难以流式运行。联合CTC与注意力训练用CTC分支找回单调约束；RNN-T则在同一个格点上同时保留路径求和与输出历史。Whisper不改变基本结构，而是用大规模弱监督数据和任务词元扩大覆盖。流式识别最后限制编码器可见的声音范围。第\ref{sec:nlp-speech-asr-comparison}小节把这些差别汇总为对照表。

\begin{figure*}[htbp]
  \centering
  % 第九卷 15.2.1 节：语音识别模型谱系（跨栏，约156 mm）。
% 蓝色：显式状态对齐；黄色：对齐路径求和；红色：注意力条件生成。
% 实线箭头表示“改变了前一方法的哪个机制”，虚线表示另一来源的组成部分。
\begin{tikzpicture}[
  x=1mm,y=1mm,
  figure picture,
  base/.style={figure node,minimum width=30mm,minimum height=11mm,
    inner ysep=2pt,font=\figurefont\scriptsize},
  nst/.style={base,draw=FigureInput,fill=FigureInputFill},
  nps/.style={base,draw=FigureModel,fill=FigureCream},
  nat/.style={base,draw=FigureOutput,fill=FigureOutputFill},
  fl/.style={figure flow},
  dl/.style={figure flow,draw=FigureMuted,dash pattern=on 2pt off 1.5pt},
  el/.style={font=\figurefont\scriptsize,text=FigureMuted,inner sep=1pt,align=center}
]
  \node[nst] (hmm)   at (15,40) {HMM--GMM\\帧级状态与GMM似然};
  \node[nst] (hyb)   at (15,18) {混合HMM--DNN\\网络给出状态后验};
  \node[nps] (ctc)   at (57,40) {CTC\\对单调路径求和};
  \node[nps] (rnnt)  at (57,18) {RNN-T\\路径求和＋输出历史};
  \node[nps] (chunk) at (57,-4) {分块流式编码器\\限制右侧可见范围};
  \node[nat] (aed)   at (99,40) {注意力编码器—解码器\\读取历史，无单调约束};
  \node[nat] (joint) at (99,18) {联合CTC/注意力\\CTC分支补单调约束};
  \node[nat] (wsp)   at (141,40) {Whisper\\弱监督与任务词元};
  \node[nat] (wss)   at (141,18) {Whisper-Streaming\\稳定前缀提交};

  \draw[fl] (hmm) -- node[el,left] {神经网络\\替换GMM} (hyb);
  \draw[fl] (hyb.north east) -- node[el,below right=-1mm and 0mm,pos=.45] {去掉帧级\\对齐标签} (ctc.south west);
  \draw[fl] (ctc) -- node[el,above] {引入\\输出历史} (aed);
  \draw[fl] (ctc) -- node[el,right] {加入\\预测网络} (rnnt);
  \draw[dl] (aed.south west) -- (rnnt.north east);
  \draw[fl] (aed) -- node[el,right] {共享编码器\\加CTC分支} (joint);
  \draw[fl] (aed) -- node[el,above] {扩大\\数据覆盖} (wsp);
  \draw[fl] (wsp) -- node[el,right] {缓冲重识别} (wss);
  \draw[fl] (rnnt) -- node[el,right] {因果化\\与分块} (chunk);

  % 图例
  \node[nst,minimum width=6mm,minimum height=4mm] (l1) at (88,-4) {};
  \node[anchor=west,font=\figurefont\scriptsize] at (l1.east) {显式状态对齐};
  \node[nps,minimum width=6mm,minimum height=4mm] (l2) at (111,-4) {};
  \node[anchor=west,font=\figurefont\scriptsize] at (l2.east) {对齐路径求和};
  \node[nat,minimum width=6mm,minimum height=4mm] (l3) at (134,-4) {};
  \node[anchor=west,font=\figurefont\scriptsize] at (l3.east) {注意力条件生成};
\end{tikzpicture}

  \caption{语音识别模型谱系。颜色表示处理声文对应的方式：蓝色用显式状态对齐，黄色对所有单调对齐路径求和，红色用注意力在生成时读取声音。实线箭头旁注明后一方法改变了前一方法的哪个机制；虚线表示RNN-T同时吸收了注意力模型读取输出历史的做法。CTC同样可以接分块流式编码器，图中只画出RNN-T一支。}
  \label{fig:nlp-speech-asr-lineage}
\end{figure*}

\subsection{HMM识别路线}
\label{sec:nlp-speech-hmm}

最早系统处理变长对应的思路，是让每个发音单位在若干声学状态中停留若干帧，再按顺序前进到下一个状态。设$W$为词序列，识别要找后验概率最大的词序列。由贝叶斯公式，$p(\bm{X})$对所有候选相同，可以略去：
\begin{equation}
 \hat W=\argmax_W p(W\mid\bm{X})
       =\argmax_W p(\bm{X}\mid W)\,P(W).
 \label{eq:nlp-speech-noisy-channel}
\end{equation}
$p(\bm{X}\mid W)$由\term{声学模型}（acoustic model）和\term{发音词典}（pronunciation lexicon）共同给出：词典把词展开成音素串，每个音素对应一个小HMM，声学模型评价每帧特征在状态下的似然。$P(W)$由第\ref{sec:nlp-classical-ngram-models}节的$n$元语言模型给出。三者最终组合成一张统一的解码图，在图上搜索最优路径。

HMM的前向算法、Viterbi算法和Baum--Welch估计已在第\ref{sec:nlp-classical-hmm}小节以词性标注为例推导过；带缩放的前向后向算法和期望计数的一般推导可参考\BookChapterRef{chap:pgm-dynamic}{卷六动态有向模型章}。本小节只写识别特有的部分：音素状态的拓扑、高斯混合似然、上下文相关状态的绑定，以及解码图的组合。本小节以$t$索引特征帧（帧移10\,ms，不降采样），$T$为帧数。

\subsubsection{HMM--GMM声学模型}
\label{sec:nlp-speech-hmm-gmm}

每个音素用一个三状态从左到右的HMM表示，三个状态大致对应音的起始、稳定段与过渡段。状态$j$只允许自环或前进到$j+1$，转移概率为$a_{j,j}$和$a_{j,j+1}=1-a_{j,j}$。自环让一个状态持续任意帧：停留恰好$k$帧的概率为$a_{j,j}^{k-1}(1-a_{j,j})$，平均停留$1/(1-a_{j,j})$帧。例如$a_{j,j}=0.8$时平均停留5帧，即50\,ms。这种几何分布的时长模型很粗糙，但它让同一个HMM能适应不同语速。按词典把音素HMM首尾相接，就得到词的HMM；再按词序列相接，词间可插入可选的静音模型，就得到整句的HMM。

状态$j$下观测特征$\bm{x}_t\in\mathbb{R}^{d}$的似然用\term{高斯混合模型}（Gaussian mixture model, GMM）表示。传统系统常取$d=39$，即13维MFCC及其一阶、二阶差分。设状态$j$有$M$个分量，权重$c_{j,m}\geq0$、$\sum_mc_{j,m}=1$，均值$\bm{\mu}_{j,m}$，对角协方差$\operatorname{diag}(\bm{\sigma}_{j,m}^2)$，则
\begin{equation}
\begin{aligned}
 b_j(\bm{x}_t)&=\sum_{m=1}^{M}c_{j,m}\,
   \mathcal{N}\bigl(\bm{x}_t;\bm{\mu}_{j,m},\operatorname{diag}(\bm{\sigma}_{j,m}^2)\bigr),\\
 \log\mathcal{N}(\cdot)&=-\frac12\sum_{k=1}^{d}
   \left[\log(2\pi\sigma_{j,m,k}^2)+\frac{(x_{t,k}-\mu_{j,m,k})^2}{\sigma_{j,m,k}^2}\right].
\end{aligned}
\label{eq:nlp-speech-gmm}
\end{equation}
对角协方差把每个分量的参数从$O(d^2)$降到$O(d)$，计算也只需逐维累加。这一简化依赖特征各维近似去相关；第\ref{sec:nlp-speech-features}小节MFCC中的离散余弦变换正起这个作用，这也是MFCC长期与GMM搭配的原因。

识别时在整句HMM上运行Viterbi算法。用对数分数表示，到第$t$帧处于状态$s$的最优路径分数为
\begin{equation}
 v_t(s)=\log b_s(\bm{x}_t)+
       \max_{r\to s}\{v_{t-1}(r)+\log a_{r,s}\},
 \label{eq:nlp-speech-hmm-viterbi}
\end{equation}
$r\to s$取遍允许转移到$s$的前驱。起始状态取对数概率0，其余取$-\infty$；保存取最大值的前驱，在允许的终态中选最佳分数并回溯。词边界处再加上语言模型的对数分数。若教学词典把“hi”映射成两个音素/h/、/ay/，一条可能的回溯路径是$h_1,h_1,h_2,h_3,ay_1,ay_2,ay_2,ay_3$；沿词典读出的是一个词“hi”，八个状态并不对应八个输出符号。

同一个音素在不同邻音之间的发音差别很大，例如/ay/在/h/之后与在/m/之后的起始段不同。\term{三音素}（triphone）模型为每个“左邻—中心—右邻”组合单独建模。设音素集有40个，三音素有$40^3=64000$种，每种3个状态，共约19万个状态；若每个状态用16个39维对角高斯，参数约2.4亿，且大多数组合在训练数据中从未出现。\term{状态绑定}（state tying）用决策树解决这个问题：对同一中心音素同一位置的所有三音素状态，从根节点开始，用“左邻是否为鼻音”“右邻是否为元音”这类语音学问题逐层二分，每次选择使训练数据对数似然增益最大的问题，增益低于阈值或节点数据量过少时停止\scite{nlp-speech-young1994}。叶节点就是绑定后的状态，通常只有数千个；训练中没见过的三音素沿树的问题走到某个叶节点，也能得到声学模型。

参数用Baum--Welch算法估计。每条训练语句按转写和词典拼成整句HMM，路径只在这张图中展开，这种做法称为嵌入式训练。E步用前向后向算法算出状态占用$\gamma_t(j)$，再按分量似然分配到各个分量：$\gamma_t(j,m)=\gamma_t(j)\,c_{j,m}\mathcal{N}_{j,m}(\bm{x}_t)/b_j(\bm{x}_t)$。M步用加权统计更新参数：
\begin{equation}
\begin{aligned}
 \bm{\mu}_{j,m}&=\frac{\sum_t\gamma_t(j,m)\bm{x}_t}{\sum_t\gamma_t(j,m)},\qquad
 c_{j,m}=\frac{\sum_t\gamma_t(j,m)}{\sum_t\gamma_t(j)},\\
 \bm{\sigma}_{j,m}^2&=\frac{\sum_t\gamma_t(j,m)(\bm{x}_t-\bm{\mu}_{j,m})^{\odot2}}{\sum_t\gamma_t(j,m)}.
\end{aligned}
 \label{eq:nlp-speech-gmm-update}
\end{equation}
转移概率按第\ref{sec:nlp-classical-baum-welch}小节的期望转移计数更新。实际流程从单音素模型开始，先把各段均匀切分作为初始对齐，迭代收敛后再展开为三音素、建决策树并继续迭代。

HMM--GMM的优点是每一帧属于哪个状态都可以检查，对齐和时长可解释，所需数据量也较小。它的局限来自两个假设：给定状态后各帧独立，相邻帧的强相关被忽略；GMM在对角协方差下难以刻画复杂的特征分布，也难以直接使用多帧拼接的高维输入。

\subsubsection{混合HMM--DNN与解码图}
\label{sec:nlp-speech-hybrid}

GMM表达力的不足可以用判别式神经网络弥补，同时保留HMM的状态结构和解码流程。混合系统的深度网络读取以第$t$帧为中心的一个窗口，例如前后各5帧、每帧40维log-Mel，拼成440维输入，经过多个隐藏层后对数千个绑定状态做softmax，输出状态后验$p(s\mid\bm{x}_t)$\scite{nlp-speech-hinton2012}。帧级训练标签来自已训练好的HMM--GMM在转写上的强制对齐（第\ref{sec:nlp-speech-forced-alignment}小节），损失为逐帧交叉熵；之后还可以用整句的序列判别目标继续训练。

HMM的解码需要似然$p(\bm{x}_t\mid s)$，网络给出的却是后验。由贝叶斯公式，
\begin{equation}
 \log p(\bm{x}_t\mid s)=\log p(s\mid\bm{x}_t)-\log p(s)+\log p(\bm{x}_t).
 \label{eq:nlp-speech-hybrid-likelihood}
\end{equation}
最后一项对同一帧的所有状态相同，比较候选路径时可以略去；状态先验$p(s)$用对齐中各状态出现的频率估计。因此，混合系统用后验除以先验得到的伪似然替换式\eqref{eq:nlp-speech-gmm}，其余部分不变。网络后验与HMM转移、语言模型分数的尺度不一致，解码时还要给声学分数乘一个小于1的缩放系数，在开发集上选定。

接下来要把HMM状态、三音素、词典和语言模型组合成一张可以搜索的图。\term{加权有限状态转导器}（weighted finite-state transducer, WFST）提供了统一的表示：它是一个有向图，每条弧带有输入符号、输出符号和权重，一条从初态到终态的路径把它读入的输入串映射成输出串，路径权重为各弧权重之和（权重取负对数概率）。两个转导器$T_1$与$T_2$的\term{复合}（composition）$T_1\circ T_2$把$T_1$的输出接到$T_2$的输入：若$T_1$把$a$映射为$b$、代价为$w_1$，$T_2$把$b$映射为$c$、代价为$w_2$，则复合结果把$a$直接映射为$c$，代价为对所有中间串$b$取$\min(w_1+w_2)$。复合的含义就是把两步映射合成一步，并把两步的代价相加。

识别中依次使用四个转导器：$H$把HMM转移标识映射为上下文相关音素，其中编码了各状态的拓扑与转移代价；$C$把上下文相关音素映射为上下文无关音素；$L$是发音词典，把音素串映射为词；$G$是语言模型，作为只接受词序列的接受器，权重为$n$元模型的负对数概率。Kaldi工具包把它们组合为
\begin{equation}
 \mathit{HCLG}=\min\bigl(\det(H\circ C\circ L\circ G)\bigr),
 \label{eq:nlp-speech-hclg}
\end{equation}
这里略去了实现中逐级插入的确定化与最小化步骤\scite{nlp-speech-mohri2002}。\term{确定化}（determinization）使每个状态出发的弧不出现相同的输入符号，每个输入串至多对应一条路径，搜索时不必在等价路径上重复计算；\term{最小化}（minimization）合并后续行为完全相同的状态，缩小图的规模。词典中的同音词会使$L$无法确定化：若“hi”和“high”都读作/h ay/，同一音素串对应两个词。实现时在同音词的音素串末尾加上不同的消歧符号，确定化后再去掉。复合后的图输入转移标识、输出词，权重中已经包含转移与语言模型分数。

解码在$\mathit{HCLG}$上做带剪枝的Viterbi搜索。每读入一帧，活动路径沿输入非空的弧前进一步，代价加上$-\kappa\log p(\bm{x}_t\mid s)$，其中$s$由弧上的转移标识确定，$\kappa$是上述声学缩放系数；代价比当前最优路径高出束宽以上的路径被删除。经过输出词的弧时记录该词，最后回溯得到词序列，也可以保留一个词格供第\ref{sec:nlp-speech-fusion}小节重打分\scite{nlp-speech-kaldi}

这条路线的错误可以分成三类定位。发音词典没有“海棠楼”时，解码图里根本没有这个词，无论声学模型多好都无法输出，属于词典缺项；声学模型给错误状态更高的分数，属于评分错误；正确路径的总分其实更高，却在搜索中被剪掉，属于搜索错误。区分后两类的办法是计算参考文本在图中的强制路径分数：若它比输出路径分数高，就是搜索错误，应放宽束宽而不是重训模型。混合系统的主要负担在工程上：单音素、三音素、决策树、特征变换、说话者自适应和神经网络各阶段都需要单独训练和调参，前一阶段的对齐错误还会传给后一阶段。

\subsection{对齐路径求和：CTC}
\label{sec:nlp-speech-ctc}

混合系统仍依赖帧级对齐标签和多个独立训练的阶段。如果只有声音与整句文字，又不想先指定一个唯一的帧级对齐，可以把所有合法对应的概率加起来作为文本的概率。\term{连接时序分类}（connectionist temporal classification, CTC）就是这样定义似然的\scite{nlp-speech-graves2006}。它要求每个编码步的输出在给定声音后相互条件独立，因此结构简单、天然单调、容易流式；代价是不显式使用输出历史，常需外部语言模型补足拼写与语言规律。CTC在全书其他位置没有正文，本小节完整写出递推、梯度与解码。

\subsubsection{扩展标签与前向--后向递推}
\label{sec:nlp-speech-ctc-forward-backward}

CTC在词表$V$外增加\term{空白符}（blank）$\varnothing$，编码器之后接一个线性层与softmax，在每个编码步给出$|V|+1$维分布$q_t(k)=p(k\mid\bm{X},t)$。空白表示此步不输出新符号，它不等于静音。长度为$T$的路径$\pi=(\pi_1,\ldots,\pi_T)$每步取一个符号；记$\mathcal{B}$为“先合并相邻相同符号，再删除空白”的折叠变换，则
\begin{equation}
 p(y\mid\bm{X})=
 \sum_{\pi:\mathcal{B}(\pi)=y}\prod_{t=1}^{T}q_t(\pi_t),
 \qquad L_{\mathrm{CTC}}=-\log p(y\mid\bm{X}).
 \label{eq:nlp-speech-ctc-probability}
\end{equation}
乘积形式假定给定$\bm{X}$后各步的符号选择条件独立。这里独立的是路径上的输出，不是声音帧：双向编码器或自注意力编码器仍可以让每个$q_t$依赖整段声音。但$q_t$不读取已经输出的文字，模型无法直接学到“上一个字是‘四’，下一个字更可能是‘月’”这类规律。

折叠的顺序不能颠倒：路径$(a,a,\varnothing,a)$先合并为$(a,\varnothing,a)$，再删空白得到$aa$；若先删空白再合并，就会错误地得到$a$。因此目标中相邻的相同符号之间至少要有一个空白隔开。令$r(y)=\sum_{u=2}^{U}\mathbf{1}[y_u=y_{u-1}]$为相邻重复数，非空目标存在合法路径的条件是$T\geq U+r(y)$。这是对编码器输出长度的约束：第\ref{sec:nlp-speech-encoder}小节的降采样倍数过大时，这个条件可能不满足，样本的损失为无穷大。

直接枚举路径的数量随$T$指数增长，需要动态规划。在目标首尾和相邻符号之间插入空白，得到\term{扩展标签}（extended label）
$z=(\varnothing,y_1,\varnothing,\ldots,y_U,\varnothing)$，
长度为$2U+1$。定义前向变量$\alpha_{t,s}$为前$t$步、终于扩展标签位置$s$的所有合法部分路径的概率之和。对$U>0$，初值为
\begin{equation}
 \alpha_{1,1}=q_1(\varnothing),\quad
 \alpha_{1,2}=q_1(y_1),\quad
 \alpha_{1,s}=0\quad(s>2).
 \label{eq:nlp-speech-ctc-initial}
\end{equation}
越界位置的概率取0。到达位置$s$有三种来源：停在原位置，从前一位置进入，或跳过中间的空白从$s-2$进入。最后一种只在当前标签非空白、且与$s-2$处的标签不同时允许，否则目标$aa$中的两个$a$会被走成一个持续的$a$：
\begin{equation}
\begin{aligned}
 \delta_s&=\mathbf{1}[s>2,\ z_s\ne\varnothing,\ z_s\ne z_{s-2}],\\
 \alpha_{t,s}
 &=q_t(z_s)\bigl(\alpha_{t-1,s}+\alpha_{t-1,s-1}
                    +\delta_s\alpha_{t-1,s-2}\bigr),\\
 p(y\mid\bm{X})&=\alpha_{T,2U}+\alpha_{T,2U+1}.
\end{aligned}
\label{eq:nlp-speech-ctc-forward}
\end{equation}
合法的结束位置只有最后一个标签和尾部空白两个。空目标只有全空白路径，概率为$\prod_tq_t(\varnothing)$。

表\ref{tab:nlp-speech-ctc-numbers}给出三步完整分布及目标$ab$的前向变量，分布为教学构造。第三步到达位置4时，三个来源的前向变量为$0.04,0.16,0.45$，乘以$q_3(b)=0.1$得$0.065$；再加尾空白的$0.020$，目标概率为$0.085$。

\begin{table}[htbp]
  \centering
  \begin{tabular}{@{}crrrrrrrr@{}}
    \toprule
    & \multicolumn{3}{c}{逐步分布} & \multicolumn{5}{c}{$z=(\varnothing,a,\varnothing,b,\varnothing)$}\\
    $t$ & $\varnothing$ & $a$ & $b$ & $\alpha_{t,1}$ & $\alpha_{t,2}$ & $\alpha_{t,3}$ & $\alpha_{t,4}$ & $\alpha_{t,5}$\\
    \midrule
    1 & .5 & .4 & .1 & .5 & .4 & 0 & 0 & 0\\
    2 & .4 & .5 & .1 & .2 & .45 & .16 & .04 & 0\\
    3 & .5 & .4 & .1 & .1 & .26 & .305 & .065 & .020\\
    \bottomrule
  \end{tabular}
  \caption{CTC前向变量的教学计算。逐步分布为构造数值，前向变量由式\eqref{eq:nlp-speech-ctc-forward}递推得到；每行的前向变量是部分路径概率之和，不是归一化分布。}
  \label{tab:nlp-speech-ctc-numbers}
\end{table}

同一结果可以完全枚举核对。折叠为$ab$的五条路径是
$(a,a,b)$、$(a,b,b)$、$(a,\varnothing,b)$、$(\varnothing,a,b)$和$(a,b,\varnothing)$，
概率分别为$0.020,0.004,0.016,0.025,0.020$，总和$0.085$，损失$-\log0.085=2.4651$。目标$aa$在这三步中只有$(a,\varnothing,a)$一条路径，概率为$0.064$。

求梯度还需要后向变量。定义$\beta_{t,s}$为第$t$步位于位置$s$时，后续第$t+1$至$T$步能够合法走到终点的所有路径后缀的概率之和；它不含第$t$步本身的输出概率。终点条件和逆序递推为
\begin{equation}
\begin{aligned}
 \beta_{T,s}&=\mathbf{1}[s\in\{2U,2U+1\}],\\
 \beta_{t,s}&=q_{t+1}(z_s)\beta_{t+1,s}+q_{t+1}(z_{s+1})\beta_{t+1,s+1}
   +\delta_{s+2}\,q_{t+1}(z_{s+2})\beta_{t+1,s+2},
\end{aligned}
\label{eq:nlp-speech-ctc-backward}
\end{equation}
三项与前向递推的三种来源一一对应，越界项取0。按这个定义，$\alpha_{t,s}\beta_{t,s}$恰好是所有在第$t$步经过位置$s$的完整路径概率之和。每条合法路径在每一步都恰好经过一个位置，于是对任意$t$都有$\sum_s\alpha_{t,s}\beta_{t,s}=p(y\mid\bm{X})$，这是检查实现的常用等式。在上例中，$\beta_{3,:}=(0,0,0,1,1)$，$\beta_{2,:}=(0,.1,.1,.6,.5)$，$\beta_{1,:}=(.05,.15,.10,.26,.20)$。第2步的检查为$.45\times.1+.16\times.1+.04\times.6=0.085$，第1步为$.5\times.05+.4\times.15=0.085$。

实际计算在对数域进行，以免长序列的连乘下溢。令$\ell_{t,s}=\log\alpha_{t,s}$，把式\eqref{eq:nlp-speech-ctc-forward}中的加法换成
$\operatorname{LSE}(v_1,\ldots,v_k)=m+\log\sum_j\exp(v_j-m)$，
其中$m=\max_jv_j$；不可达项取$-\infty$，全部不可达时直接返回$-\infty$。后向变量同样处理。前向与后向各需$O(TU)$次状态更新；只求损失时可以滚动保存两行，求梯度则要保存或重算全部状态。

\subsubsection{损失梯度与尖峰现象}
\label{sec:nlp-speech-ctc-gradient}

有了前向与后向变量，就能算出每一步每个符号被“占用”的后验概率。第$t$步位于位置$s$的后验为$\alpha_{t,s}\beta_{t,s}/p(y\mid\bm{X})$；把同一符号的所有位置加起来，得到符号$k$在第$t$步的\term{后验占用}（posterior occupancy）
\begin{equation}
 \gamma_t(k)=\frac{1}{p(y\mid\bm{X})}\sum_{s:\,z_s=k}\alpha_{t,s}\beta_{t,s},
 \qquad \sum_{k}\gamma_t(k)=1.
 \label{eq:nlp-speech-ctc-occupancy}
\end{equation}
它表示：在所有能折叠为目标的路径中按概率加权，第$t$步输出$k$的比例。

设第$t$步softmax之前的未归一化分数为$o_{t,k}$，$q_t=\operatorname{softmax}(\bm{o}_t)$。$q_t(k)$在每条经过$(t,s)$且$z_s=k$的路径中恰好出现一次，因此$\partial p/\partial q_t(k)=\sum_{s:z_s=k}\alpha_{t,s}\beta_{t,s}/q_t(k)$，即$\partial\log p/\partial q_t(k)=\gamma_t(k)/q_t(k)$。再用softmax的导数$\partial q_t(j)/\partial o_{t,k}=q_t(j)(\mathbf{1}[j=k]-q_t(k))$求和：
\begin{equation}
 \frac{\partial L_{\mathrm{CTC}}}{\partial o_{t,k}}
 =-\sum_j\frac{\gamma_t(j)}{q_t(j)}\,q_t(j)\bigl(\mathbf{1}[j=k]-q_t(k)\bigr)
 =q_t(k)-\gamma_t(k).
 \label{eq:nlp-speech-ctc-gradient}
\end{equation}
最后一步用到$\sum_j\gamma_t(j)=1$。这个结果与普通分类的交叉熵梯度形式相同，只是“标签”换成了由当前模型自己算出的后验占用：CTC相当于在每一步用软目标$\gamma_t$训练分类器，而软目标随模型更新而改变。这与EM的思路一致，E步算对齐的后验，M步向它拟合，区别在于CTC用梯度一步完成两者。

在表\ref{tab:nlp-speech-ctc-numbers}的例子中，第2步的位置后验为$(0,.529,.188,.282,0)$，于是$\gamma_2(\varnothing)=.188$、$\gamma_2(a)=.529$、$\gamma_2(b)=.282$。梯度$q_2-\gamma_2=(.212,-.029,-.182)$：梯度下降会降低这一步空白的分数，主要提高$b$的分数。

这个梯度也解释了训练后CTC输出的\term{尖峰现象}（peaky behavior）：多数步以接近1的概率输出空白，每个非空白符号只在一两步上出现高概率的尖峰\scite{nlp-speech-zeyer2021}。原因在于两类符号在路径中的地位不同。每个标签至少占一步，而空白可以填充标签之间任意多步，并且扩展标签中有$U+1$个空白位置。训练初期各符号概率接近均匀，后验占用中空白占很大比例，梯度便在多数步上推高空白。空白概率升高后，每个标签只需在其相对概率最高的一步出现即可完成折叠，占用进一步集中到这一步，形成正反馈。尖峰只表示“在这一步输出该符号”，CTC目标从未要求它落在发音的起点、中心或终点；单向流式编码器还常把尖峰推迟到看见足够声音之后。图\ref{fig:nlp-speech-alignment}用两条折叠为同一文本的路径说明，符号出现在哪一步本身就是不确定的。

\begin{figure*}[htbp]
  \centering
  % Adapted from the academic-drawing TikZ template; dimensions are final size.
\begin{tikzpicture}[
  x=1mm,y=1mm,
  every node/.style={font=\figurefont\small,align=center,text=BookInk},
  cell/.style={draw=BookRule,line width=.8pt,minimum width=13mm,
    minimum height=9mm,inner sep=1mm},
  label/.style={anchor=east},
  flow/.style={book arrow}
]
  \foreach \i in {1,...,6} {
    \node[cell,fill=FigureInput!12] at ({27+17*\i},0) {$\bm{h}_{\i}$};
  }
  \node[label] at (31,0) {编码步};
  \node[label] at (31,-15) {路径A};
  \node[label] at (31,-30) {路径B};
  \foreach \i/\v in {1/八,2/八,3/{$\varnothing$},4/八,5/{$\varnothing$},6/{$\varnothing$}} {
    \node[cell,fill=FigureModel!10] at ({27+17*\i},-15) {\v};
  }
  \foreach \i/\v in {1/{$\varnothing$},2/八,3/{$\varnothing$},4/八,5/八,6/{$\varnothing$}} {
    \node[cell,fill=FigureModel!10] at ({27+17*\i},-30) {\v};
  }
  \draw[flow] (137,-15) -- (145,-15);
  \draw[flow] (137,-30) -- (145,-30);
  \node[anchor=west,text=FigureOutput] at (146,-15) {八八};
  \node[anchor=west,text=FigureOutput] at (146,-30) {八八};
  \draw[BookRule,line width=.8pt] (5,-42) -- (160,-42);
  \node[label] at (31,-54) {独立标注};
  \node[anchor=west,align=left] at (37,-54)
    {第一个“八”：0.02--0.10\,s，说话者甲\\
     第二个“八”：0.13--0.22\,s，说话者甲};
\end{tikzpicture}

  \caption{同一声学序列的两条CTC教学路径均折叠为“八八”，但符号出现的位置不同。空白为$\varnothing$，相邻重复先合并，再删除空白。最下方的时间区间与说话者是另行给定的示意标注，不是由两条路径证明的精确时间。}
  \label{fig:nlp-speech-alignment}
\end{figure*}
\FloatBarrier

\subsubsection{贪心解码与前缀束搜索}
\label{sec:nlp-speech-ctc-decoding}

最简单的解码是每步取概率最大的符号再折叠。但这样得到的是概率最大的单条路径所对应的文本，而不是总概率最大的文本，因为一个文本的概率是许多路径之和。构造两步相同分布$q(\varnothing)=0.40$、$q(a)=0.35$、$q(b)=0.25$。贪心选出全空白路径，对应空文本，概率$0.16$；文本$a$却汇总了$(a,a)$、$(a,\varnothing)$、$(\varnothing,a)$三条路径，概率$0.1225+0.14+0.14=0.4025$，远高于空文本。

前缀束搜索把折叠后对应同一文字前缀的路径合并，但必须区分路径的末步是空白还是非空白。令$b_t(l)$与$n_t(l)$分别为前$t$步折叠为$l$、且末步为空白与非空白的概率之和；初值$b_0(\epsilon)=1$，其余为0。区分两者的原因是：若末字为$a$的前缀再遇到$a$，非空白结尾的路径只能延续同一个$a$，空白结尾的路径才能形成新的$a$。算法\ref{alg:nlp-speech-prefix}给出这一过程\scite{nlp-speech-hannun2014}

\begin{algorithm}[CTC前缀束搜索]
  \label{alg:nlp-speech-prefix}
  \begin{algorithmic}[1]
    \Require 分布$q_1,\ldots,q_T$，词表$V$，束宽$K$
    \Ensure 按合并概率选择的文字序列
    \State 初始化$b_0(\epsilon)=1,\ n_0(\epsilon)=0$，束为$\{\epsilon\}$
    \For{$t=1,\ldots,T$}
      \State 新映射$b',n'$全部置0
      \For{旧束中的每个前缀$l$}
        \State $s\gets b_{t-1}(l)+n_{t-1}(l)$
        \State $b'(l)\mathrel{+}=s\,q_t(\varnothing)$ \Comment{补一个空白，前缀不变}
        \For{$c\in V$}
          \If{$l\ne\epsilon$且$c=\operatorname{last}(l)$}
            \State $n'(l)\mathrel{+}=n_{t-1}(l)q_t(c)$ \Comment{延续同一个字}
            \State $n'(lc)\mathrel{+}=b_{t-1}(l)q_t(c)$ \Comment{空白之后形成新字}
          \Else
            \State $n'(lc)\mathrel{+}=s\,q_t(c)$
          \EndIf
        \EndFor
      \EndFor
      \State 合并同名前缀后，按$b'(l)+n'(l)$保留前$K$项，赋给$b_t,n_t$
    \EndFor
    \State \Return $\argmax_l\{b_T(l)+n_T(l)\}$
  \end{algorithmic}
\end{algorithm}

对上述两步例子，前缀$a$在第一步为$(b_1,n_1)=(0,0.35)$；第二步$b_2(a)=0.35\times0.40=0.14$，$n_2(a)=0.35\times0.35+0.40\times0.35=0.2625$，合计$0.4025$，与枚举一致。不剪枝时算法精确地合并了所有前缀，保留前$K$项后才成为近似搜索。

外部语言模型在形成新前缀的位置加入。把算法中更新$n'(lc)$的两行里的$q_t(c)$换成$q_t(c)\,p_{\mathrm{lm}}(c\mid l)^{\lambda}e^{\gamma}$，其中$\lambda$是语言模型权重，$\gamma$是补偿语言模型偏好短文本的插入奖励；延续同一个字的那一行不乘语言模型因子。同一个新前缀$lc$可能在不同的步上形成，但每次乘的都是同一个因子，所以等价于对整句文本只计一次语言模型分数。第\ref{sec:nlp-speech-fusion}小节在更一般的形式下讨论这种融合。

CTC用路径求和去掉了帧级对齐标签和多阶段训练，训练只需整句转写，输出天然单调，解码可以逐步进行。它的弱点同样来自条件独立假设：每一步的输出不读取已输出文字，同音字、拼写一致性等依赖输出历史的规律只能靠编码器从声音中间接学到，或交给外部语言模型。

\subsection{注意力编码器—解码器}
\label{sec:nlp-speech-aed}

要让后一个字的概率直接依赖前面已经输出的字，可以把识别看作以声音为条件的序列生成。\term{注意力编码器—解码器}（attention-based encoder--decoder, AED）由编码器读入整段声音，解码器在每一步读取已输出的前缀，并用注意力从编码序列中取出与当前输出相关的上下文，直接最大化整句的条件概率。这种结构与第\ref{sec:nlp-generation-encoder-decoder}小节的机器翻译模型相同，教师强制训练与束搜索的一般机制分别见第\ref{sec:nlp-generation-training}小节与第\ref{sec:nlp-generation-beam}小节。它解决了CTC的条件独立问题，但带来两项新代价：注意力没有单调约束，长语音上可能跳读、重复或过早结束；注意力需要整段编码，难以流式运行。

\subsubsection{LAS金字塔编码与注意力解码}
\label{sec:nlp-speech-las}

Listen, Attend and Spell（LAS）是这类模型的代表，由编码器listener和解码器speller组成\scite{nlp-speech-chan2015}。语音特征序列比文字长得多：10\,ms帧移下，10秒声音有1000帧。注意力每步要在全部编码位置上计算权重，位置太多时不仅计算量大，训练也难以收敛。listener因此采用\term{金字塔双向LSTM}（pyramidal bidirectional LSTM, pBLSTM）：每一层把下层相邻两步的输出拼接后作为本层的一步输入，使长度减半。记第$j$层第$i$步的输出为$\bm{h}^{(j)}_i$，则
\begin{equation}
 \bm{h}^{(j)}_i=\operatorname{BLSTM}\bigl(\bm{h}^{(j)}_{i-1},\,
   [\bm{h}^{(j-1)}_{2i-1};\bm{h}^{(j-1)}_{2i}]\bigr).
 \label{eq:nlp-speech-pblstm}
\end{equation}
LSTM与双向循环网络的结构可参考\BookChapterRef{chap:recurrent-networks}{卷五循环神经网络章}。设输入为$\bm{X}\in\mathbb{R}^{T_0\times80}$，底层BLSTM每个方向$d$维，输出$T_0\times2d$；第一层金字塔把相邻两步拼成$4d$维输入，输出$(T_0/2)\times2d$。叠加三层后长度为$T_0/8$，每步对应80\,ms；10秒声音的1000帧缩为125步。降低长度同时让每个编码步汇总更长一段声音，注意力更容易找到对应位置。

speller是一个LSTM解码器。第$u$步的状态$\bm{s}_u$由前一状态、前一个输出字符$y_{u-1}$和前一步上下文$\bm{c}_{u-1}$更新；然后用$\bm{s}_u$与每个编码步打分，得到注意力权重与上下文向量：
\begin{equation}
 e_{u,t}=\langle\phi(\bm{s}_u),\psi(\bm{h}_t)\rangle,\qquad
 a_{u,t}=\frac{\exp e_{u,t}}{\sum_{t'}\exp e_{u,t'}},\qquad
 \bm{c}_u=\sum_{t}a_{u,t}\bm{h}_t,
 \label{eq:nlp-speech-las-attention}
\end{equation}
其中$\phi,\psi$是两个多层感知机，把解码状态与编码表示映射到同一空间后做内积。输出层以$[\bm{s}_u;\bm{c}_u]$经多层感知机与softmax给出字符分布。目标序列末尾加结束符，损失为
\begin{equation}
 L_{\mathrm{AED}}=-\sum_{u=1}^{U+1}
       \log p_\theta(y_u\mid\bm{X},y_{<u}),
 \qquad y_{U+1}=\texttt{EOS}.
 \label{eq:nlp-speech-las}
\end{equation}
它与式\eqref{eq:nlp-speech-ctc-probability}的根本区别在于每一项都条件于$y_{<u}$，没有任何对齐变量需要求和：对齐被注意力权重隐式地处理了。

内积打分只看内容，不知道上一步读到了哪里。语音的对应是单调的，相邻输出读取的位置应该相邻，但式\eqref{eq:nlp-speech-las-attention}并不保证这一点。\term{位置感知注意力}（location-aware attention）把上一步的注意力权重$\bm{a}_{u-1}$经一维卷积得到每个位置的特征$\bm{f}_{u,t}$，加入加性打分$e_{u,t}=\bm{w}^{\top}\tanh(\bm{W}\bm{s}_{u-1}+\bm{V}\bm{h}_t+\bm{U}\bm{f}_{u,t}+\bm{b})$\scite{nlp-speech-chorowski2015}。加性打分的形式与式\eqref{eq:nlp-generation-bahdanau}相同，新增的$\bm{f}_{u,t}$让模型知道上一步读到了哪里，倾向于在附近继续读取。这一做法减轻了跳读，但它仍是学到的倾向，而不是硬约束。

仍用“八八”短录音，两个“八”由两个输出步分别生成，不需要CTC的空白分隔。假设候选“八／八／EOS”的三个条件概率为$0.7,0.6,0.8$，候选概率为$0.336$；另一候选必须按自己的前缀重新计算各步条件概率，不能只替换一个字的概率。训练时教师强制让每一步读取正确前缀，推理时却只能读取自己生成的前缀，这种不一致使早期错误可能连锁传播；LAS原文在训练中以一定概率用模型自己的采样代替正确前缀来缓解这一点。

AED的两项代价都来自注意力的全局性。注意力权重可以落在任何位置，模型在长语音上可能跳过一段、反复读同一段，或者在说完之前就输出EOS；训练集中没有的长度尤其容易出问题。双向编码与全局注意力又要求整段声音都已到达，不能边听边出字。

\subsubsection{CTC与注意力的联合训练}
\label{sec:nlp-speech-joint-ctc-attention}

CTC有单调约束但条件独立，AED读取历史但没有单调约束，二者的优缺点恰好互补。联合训练在共享编码器上同时接两个输出分支：一个线性层加softmax的CTC分支，以及一个注意力解码器，损失取加权和\scite{nlp-speech-kim2017}：
\begin{equation}
 L=\lambda L_{\mathrm{CTC}}+(1-\lambda)L_{\mathrm{AED}},\qquad 0\leq\lambda\leq1.
 \label{eq:nlp-speech-joint-loss}
\end{equation}
$\lambda$在开发集上选取。CTC损失只在单调路径上有概率，它迫使共享编码器的表示与文字顺序对齐，注意力解码器因而更快学到合理的对齐，训练早期收敛也更快。

解码时也可以让两个分支共同打分\scite{nlp-speech-watanabe2017}。注意力解码器按束搜索逐字扩展前缀，每个前缀同时计算CTC分支给出的\term{CTC前缀分数}（CTC prefix score），即所有折叠后以该前缀开头的文本的总概率。对前缀$h=gc$，$c$是新加入的字，按$c$第一次出现的步$t$分类求和：
\begin{equation}
 \Psi(gc)=\sum_{t=1}^{T}\Phi_{t-1}(g)\,q_t(c),\qquad
 \Phi_{t-1}(g)=\begin{cases}
   b_{t-1}(g), & c=\operatorname{last}(g),\\
   b_{t-1}(g)+n_{t-1}(g), & \text{其他},
 \end{cases}
 \label{eq:nlp-speech-ctc-prefix-score}
\end{equation}
其中$b,n$是算法\ref{alg:nlp-speech-prefix}中的空白结尾与非空白结尾概率，按完整的$T$步计算而不剪枝。若候选在$g$之后输出EOS，CTC分数取整句概率$b_T(g)+n_T(g)$。联合分数为
\begin{equation}
 S(h)=\lambda\log\Psi(h)+(1-\lambda)\log p_{\mathrm{AED}}(h\mid\bm{X})
      +\mu\log p_{\mathrm{lm}}(h),
 \label{eq:nlp-speech-joint-score}
\end{equation}
最后一项是可选的外部语言模型。仍用两步分布$q(\varnothing)=.40$、$q(a)=.35$、$q(b)=.25$：$\Psi(a)=1\times.35+.40\times.35=.49$，它等于以$a$开头的两个文本$a$（$.4025$）与$ab$（$.0875$）的概率之和；若候选在$a$之后结束，CTC分数取$.4025$。

CTC前缀分数的作用在于否决不单调的候选。若注意力解码器跳过了一段声音，跳过的字无法与CTC的单调路径对上，$\Psi$会很小；若解码器过早输出EOS，剩余声音找不到对应的文字，整句CTC概率也很小。另一种做法是先用其中一个分支生成$n$个完整候选，再用另一个分支重打分，实现简单，但只能在已有候选中选择。联合模型的代价是两个分支都要计算，且解码仍需完整编码，并未解决流式问题；第\ref{sec:nlp-speech-chunk-attention}小节的两遍解码把它与流式编码器结合起来。

\subsection{历史条件与路径求和：RNN-T}
\label{sec:nlp-speech-rnnt}

CTC可以流式但不读输出历史，AED读取历史却不单调、难以流式。\term{循环神经网络转导器}（recurrent neural network transducer, RNN-T）把两者的优点放进同一个概率模型\scite{nlp-speech-graves2012}。它在声学步$t$与输出位置$u$组成的二维格点上定义两类边：空白边推进时间，标签边推进输出，同一个声学步可以连续输出多个符号。每个格点的输出分布同时依赖编码器的$\bm{h}_t$和已输出的前缀，路径求和仍保证单调对齐。代价是联合网络要在$T\times U$个格点上各计算一次词表分布，训练显存很大。

\subsubsection{预测网络与联合网络}
\label{sec:nlp-speech-rnnt-networks}

RNN-T由三部分组成，结构见图\ref{fig:nlp-speech-rnnt-structure}。声学编码器与CTC相同，把$\bm{X}$变成$\bm{h}_1,\ldots,\bm{h}_T$，默认采用第\ref{sec:nlp-speech-encoder}小节的Conformer。\term{预测网络}（prediction network）只读取已输出的非空白符号：以开始符号$y_0$和$y_1,\ldots,y_u$为输入，经嵌入层和一两层LSTM得到$\bm{g}_u$。它的作用类似一个只看文字的语言模型，为下一个符号提供历史条件。\term{联合网络}（joint network）把两路表示合并：
\begin{equation}
 \bm{z}_{t,u}=\bm{W}_o\tanh(\bm{W}_a\bm{h}_t+\bm{W}_p\bm{g}_u+\bm{b}),\qquad
 q_{t,u}=\operatorname{softmax}(\bm{z}_{t,u})\in\mathbb{R}^{|V|+1},
 \label{eq:nlp-speech-rnnt-joint}
\end{equation}
其中$\bm{W}_a\in\mathbb{R}^{d_{\mathrm{joint}}\times d_{\mathrm{enc}}}$、$\bm{W}_p\in\mathbb{R}^{d_{\mathrm{joint}}\times d_{\mathrm{pred}}}$把两路投影到同一维度，$\bm{W}_o\in\mathbb{R}^{(|V|+1)\times d_{\mathrm{joint}}}$输出含空白的词表分数。Graves的原始形式直接把两路输出的分数相加再做softmax，式\eqref{eq:nlp-speech-rnnt-joint}在相加之后多了一层非线性。

\begin{figure}[htbp]
  \centering
  % 第九卷 15.2.5.1 节：RNN-T的编码器、预测网络与联合网络（正文栏宽，约108 mm）。
% 蓝色：输入；黄色描边：网络模块；红色：输出分布。虚线为推断时的反馈。
\begin{tikzpicture}[
  x=1mm,y=1mm,
  figure picture,
  rinp/.style={figure node,draw=FigureInput,fill=FigureInputFill,
    minimum width=40mm,minimum height=8mm,inner ysep=2pt,font=\figurefont\scriptsize},
  rmod/.style={figure node,draw=FigureModel,fill=FigureModelFill,
    minimum width=40mm,minimum height=10mm,inner ysep=2pt,font=\figurefont\scriptsize},
  rout/.style={figure node,draw=FigureOutput,fill=FigureOutputFill,
    minimum width=86mm,minimum height=8mm,inner ysep=2pt,font=\figurefont\scriptsize},
  fl/.style={figure flow},
  dl/.style={figure flow,draw=FigureMuted,dash pattern=on 2pt off 1.5pt},
  el/.style={font=\figurefont\scriptsize,text=FigureMuted,inner sep=1pt,align=center}
]
  \node[rinp] (x) at (22,0) {声学特征$\bm{X}$};
  \node[rinp] (y) at (72,0) {已输出的非空白符号$y_0,y_1,\ldots,y_u$};
  \node[rmod] (enc) at (22,15) {声学编码器（如Conformer）\\$\bm{H}\in\mathbb{R}^{B\times T\times d_{\mathrm{enc}}}$};
  \node[rmod] (pred) at (72,15) {预测网络（嵌入＋LSTM）\\$\bm{G}\in\mathbb{R}^{B\times(U+1)\times d_{\mathrm{pred}}}$};
  \node[rmod,minimum width=86mm] (joint) at (47,33)
    {联合网络：$\bm{z}_{t,u}=\bm{W}_o\tanh(\bm{W}_a\bm{h}_t+\bm{W}_p\bm{g}_u)$\\
     按$(t,u)$广播相加，隐藏张量$B\times T\times(U+1)\times d_{\mathrm{joint}}$};
  \node[rout] (q) at (47,49) {softmax：$q_{t,u}=\operatorname{softmax}(\bm{z}_{t,u})$，
    张量$B\times T\times(U+1)\times(|V|+1)$};

  \draw[fl] (x) -- (enc);
  \draw[fl] (y) -- (pred);
  \draw[fl] (enc) -- node[el,left] {$\bm{h}_t$} (enc |- joint.south);
  \draw[fl] (pred) -- node[el,right] {$\bm{g}_u$} (pred |- joint.south);
  \draw[fl] (joint) -- (q);
  \draw[dl] (q.east) -- ++(6,0) |- node[el,pos=.25,right] {推断时\\非空白\\输出送回} (y.east);
\end{tikzpicture}

  \caption{RNN-T的三部分结构。编码器只读声音，预测网络只读已输出的非空白符号，联合网络把$\bm{h}_t$与$\bm{g}_u$投影后相加、经$\tanh$与输出层得到格点$(t,u)$上含空白的分布。训练时$\bm{H}$与$\bm{G}$按$(t,u)$广播成四维张量；推断时非空白输出沿虚线送回预测网络，空白则让编码器前进一步。}
  \label{fig:nlp-speech-rnnt-structure}
\end{figure}

张量形状说明了计算量落在哪里。批大小为$B$时，编码器输出$\bm{H}\in\mathbb{R}^{B\times T\times d_{\mathrm{enc}}}$，预测网络输出$\bm{G}\in\mathbb{R}^{B\times(U+1)\times d_{\mathrm{pred}}}$，多出的一行对应只输出了开始符号的位置$u=0$。两路投影后分别扩展为$B\times T\times1\times d_{\mathrm{joint}}$与$B\times1\times(U+1)\times d_{\mathrm{joint}}$，广播相加得到$B\times T\times(U+1)\times d_{\mathrm{joint}}$，最后输出$B\times T\times(U+1)\times(|V|+1)$。编码器和预测网络都只算一次，四维张量完全来自联合网络。

预测网络不读声音，是有意的设计。这样声学信息只经编码器进入，编码器可以独立地做成流式；预测网络也可以像语言模型一样在推断时缓存状态，只在输出非空白符号时更新。由于空白只推进时间、不进入预测网络，RNN-T的路径恢复只需删除空白，不需要合并重复：目标$aa$可以在同一时间步连续发出两个$a$，CTC的长度条件$T\geq U+r(y)$在这里不适用。若令$\bm{W}_p=\bm{0}$，分布不再依赖输出历史，模型就退化为一种条件独立的转导器，与CTC相近，但同一步仍允许输出多个符号，两者并不等同。实践中预测网络甚至可以只看最近一两个符号，称为无状态预测网络，准确率下降不多\scite{nlp-speech-ghodsi2020}；这说明RNN-T的历史条件主要提供局部的拼写与搭配信息。

\subsubsection{格点前向--后向与损失}
\label{sec:nlp-speech-rnnt-loss}

格点上的索引取$t=1,\ldots,T$、$u=0,\ldots,U$。空白边从$(t,u)$到$(t+1,u)$，消耗一个声学步；标签$y_{u+1}$从$(t,u)$到$(t,u+1)$，增加一个输出而不推进声音。固定目标后，记$b_{t,u}=q_{t,u}(\varnothing)$、$e_{t,u}=q_{t,u}(y_{u+1})$。前向变量$\alpha_{t,u}$汇总到达格点$(t,u)$的所有部分路径概率，不含该格点本身的输出。初值$\alpha_{1,0}=1$，越界项为0：
\begin{equation}
\begin{aligned}
 \alpha_{t,u}
 &=\alpha_{t-1,u}b_{t-1,u}
   +\alpha_{t,u-1}e_{t,u-1},\\
 p(y\mid\bm{X})&=\alpha_{T,U}b_{T,U}.
\end{aligned}
\label{eq:nlp-speech-rnnt-forward}
\end{equation}
首列$u=0$只能由空白进入，首行$t=1$只能由标签进入。最后必须从$(T,U)$再发出一次空白，消耗最后一个声学步，所以终点概率要乘$b_{T,U}$。

后向变量$\beta_{t,u}$汇总从格点$(t,u)$出发、含该格点输出在内走到终点的剩余概率。初值$\beta_{T,U}=b_{T,U}$，逆序递推为
\begin{equation}
 \beta_{t,u}=b_{t,u}\beta_{t+1,u}
       +e_{t,u}\beta_{t,u+1},
 \label{eq:nlp-speech-rnnt-backward}
\end{equation}
只保留合法边：$u=U$时没有标签边，$t=T$时没有通往下一声学步的空白边。于是$\beta_{1,0}=p(y\mid\bm{X})$，而$\alpha_{t,u}\beta_{t,u}$是所有经过格点$(t,u)$的路径概率之和。

构造$T=2$、$y=aa$的例子，词表只有$a$和空白。第一列三个格点的$(b,q(a))$依次为$(.4,.6)$、$(.5,.5)$、$(.7,.3)$，第二列依次为$(.3,.7)$、$(.2,.8)$、$(.9,.1)$，合法边见图\ref{fig:nlp-speech-rnnt-lattice}。三条完整路径
$(a,a,\varnothing,\varnothing)$、
$(a,\varnothing,a,\varnothing)$、
$(\varnothing,a,a,\varnothing)$
的概率为$0.189$、$0.216$、$0.2016$，总和$0.6066$。前向递推得$\alpha_{1,:}=(1,.6,.3)$、$\alpha_{2,:}=(.4,.58,.674)$，末尾乘$.9$得到同一结果；后向递推得$\beta_{2,:}=(.504,.72,.9)$、$\beta_{1,:}=(.6066,.675,.63)$。CTC在两步内不能产生$aa$，这个例子显示了两种路径空间的区别。

\begin{figure}[htbp]
  \centering
  % Target-specific probability lattice, not a measured alignment.
\begin{tikzpicture}[
  x=1mm,y=1mm,
  every node/.style={font=\figurefont\small,text=BookInk},
  state/.style={circle,draw=FigureModel,fill=FigureModel!10,
    minimum size=5mm,inner sep=0pt,line width=.9pt},
  time/.style={book arrow,draw=BookBlue},
  token/.style={book arrow,draw=BookTeal},
  edge label/.style={fill=white,inner sep=1pt}
]
  \foreach \u in {0,1,2} {
    \node[state] (a\u) at (20,{19*\u}) {};
    \node[state] (b\u) at (62,{19*\u}) {};
    \node[anchor=east] at (9,{19*\u}) {$u=\u$};
  }
  \node at (20,-10) {$t=1$};
  \node at (62,-10) {$t=2$};
  \draw[time] (a0) -- node[edge label,above] {$\varnothing:\ .4$} (b0);
  \draw[time] (a1) -- node[edge label,above] {$\varnothing:\ .5$} (b1);
  \draw[time] (a2) -- node[edge label,above] {$\varnothing:\ .7$} (b2);
  \draw[token] (a0) -- node[edge label,right] {$a:\ .6$} (a1);
  \draw[token] (a1) -- node[edge label,right] {$a:\ .5$} (a2);
  \draw[token] (b0) -- node[edge label,right] {$a:\ .7$} (b1);
  \draw[token] (b1) -- node[edge label,right] {$a:\ .8$} (b2);
  \node[draw=FigureOutput,fill=FigureOutput!12,minimum height=8mm,
    inner sep=2mm] (end) at (97,38) {终点};
  \draw[time] (b2) -- node[edge label,above] {$\varnothing:\ .9$} (end);
  \node[anchor=west] at (5,-22) {起点$(1,0)$概率为1；终点概率为0.6066};
\end{tikzpicture}

  \caption{两编码步、目标“$aa$”的RNN-T格点。横向空白推进时间，纵向$a$推进输出；标注为教学边概率。只有从右上角发出的末尾空白进入终点，重复标签没有折叠。}
  \label{fig:nlp-speech-rnnt-lattice}
\end{figure}

梯度的推导与CTC类似。$q_{t,u}(k)$只出现在从格点$(t,u)$发出的边上，因此$\partial p/\partial q_{t,u}(k)=\alpha_{t,u}\beta'_{t,u}(k)$，其中$\beta'_{t,u}(\varnothing)=\beta_{t+1,u}$（约定$\beta_{T+1,U}=1$），$\beta'_{t,u}(y_{u+1})=\beta_{t,u+1}$，其余符号为0。经过softmax导数，得到
\begin{equation}
 \frac{\partial L}{\partial z_{t,u,k}}
 =\frac{\alpha_{t,u}\beta_{t,u}}{p}\,q_{t,u}(k)
  -\frac{\alpha_{t,u}\,q_{t,u}(k)\,\beta'_{t,u}(k)}{p}.
 \label{eq:nlp-speech-rnnt-gradient}
\end{equation}
第一项是格点占用乘以模型分布，第二项是沿符号$k$的边离开该格点的后验；形式仍是“模型输出减后验”，只是多了格点占用作为权重，模型很少经过的格点几乎不受更新。在上例的格点$(1,1)$，占用为$.6\times.675/.6066=.668$；空白边后验为$.6\times.5\times.72/.6066=.356$，标签边后验为$.6\times.5\times.63/.6066=.312$。空白的梯度为$.668\times.5-.356=-.022$，$a$的梯度为$+.022$，更新会略微提高这一格点的空白概率，因为经空白到达的$(2,1)$有更大的后续概率。

推断时，贪心解码在当前格点取概率最大的符号：空白使$t$加1，非空白追加到前缀、更新预测网络状态并保持$t$不变。束搜索需要处理同一前缀可经不同路径到达的情况。算法\ref{alg:nlp-speech-rnnt-beam}是Graves原文算法的简化形式，省略了原文中把短前缀的概率转移给长前缀的合并步骤。

\begin{algorithm}[RNN-T束搜索（简化形式）]
  \label{alg:nlp-speech-rnnt-beam}
  \begin{algorithmic}[1]
    \Require 编码输出$\bm{h}_1,\ldots,\bm{h}_T$，预测网络与联合网络，束宽$K$
    \Ensure 识别文本
    \State $B\gets\{\epsilon\}$，$P(\epsilon)=1$，缓存预测网络输出$\bm{g}(\epsilon)$
    \For{$t=1,\ldots,T$}
      \State $A\gets B$，再把$B$置空
      \While{$A$中最大概率高于$B$中第$K$大的概率（$B$不足$K$项时视为0）}
        \State 从$A$中取出概率最大的前缀$y^*$
        \State 由$\bm{h}_t$与缓存的$\bm{g}(y^*)$按式\eqref{eq:nlp-speech-rnnt-joint}计算分布$q$
        \State 把$y^*$以概率$P(y^*)q(\varnothing)$加入$B$，已存在则概率相加 \Comment{空白：进入下一步}
        \For{$k\in V$}
          \State 把$y^*k$以概率$P(y^*)q(k)$加入$A$，已存在则相加；计算并缓存$\bm{g}(y^*k)$
        \EndFor
      \EndWhile
      \State 保留$B$中概率最大的$K$个前缀
    \EndFor
    \State \Return $B$中$\log P(y)/|y|$最大的$y$
  \end{algorithmic}
\end{algorithm}

每次扩展都乘以小于1的概率，$A$中的最大概率不断下降，循环必然结束。实际系统通常还限制每个声学步最多输出的符号数，并只对概率最高的若干个$k$扩展；达到上限后强制进入下一步，会改变搜索结果，这是实现上的折中。预测网络的状态只取决于前缀，所以相同前缀共享一份缓存；预测网络很小，这部分计算远少于编码器。

\subsubsection{显存开销与剪枝训练}
\label{sec:nlp-speech-rnnt-pruned}

RNN-T训练的主要瓶颈是联合网络的四维输出。取批大小$B=16$，10秒声音经4倍降采样后$T=250$，目标长$U=60$，子词词表500个加空白，输出张量有$16\times250\times61\times501\approx1.22\times10^8$个元素，32位浮点约0.49\,GB；反向传播还要保存同样大小的梯度，以及$d_{\mathrm{joint}}=512$时约同样规模的隐藏张量。若改用约5000个汉字的词表，仅输出张量就约4.9\,GB。显存随$T$、$U$与$|V|$的乘积增长，这是CTC与AED都没有的代价。

一类做法从实现入手。函数合并把softmax、损失与梯度合在一个核函数里计算，梯度直接覆盖原来的输出缓冲，不再单独保存概率张量\scite{nlp-speech-li2019rnnt}；按长度排序组批则减少填充造成的浪费。这些做法降低常数，但不改变$T\times U\times|V|$的规模。

\term{剪枝RNN-T}（pruned RNN-T）从格点本身入手\scite{nlp-speech-kuang2022}。在整个格点中，大多数$(t,u)$的占用接近0：第$t$个声学步附近实际可能输出的只是目标中少数几个位置。它分两步计算。第一步用一个简化的加性联合网络：编码器与预测网络各自投影到词表维度，得到$\bm{l}^{\mathrm{enc}}_t$与$\bm{l}^{\mathrm{pred}}_u$，令
\begin{equation}
 \tilde{z}_{t,u,k}=l^{\mathrm{enc}}_{t,k}+l^{\mathrm{pred}}_{u,k}.
 \label{eq:nlp-speech-pruned-simple}
\end{equation}
格点递推只需要空白和目标符号这两个分数，以及softmax的归一化项$\log\sum_k\exp(l^{\mathrm{enc}}_{t,k}+l^{\mathrm{pred}}_{u,k})$；后者可以写成两个矩阵$\exp(\bm{L}^{\mathrm{enc}})\in\mathbb{R}^{T\times|V|}$与$\exp(\bm{L}^{\mathrm{pred}})^{\top}$的乘积，不必展开$T\times U\times|V|$的张量。用这个简化格点算出损失及其对各格点的梯度，就知道每个$t$上哪些$u$的占用较大。第二步为每个$t$选一个起点$p_t$，要求$p_t$随$t$不减，以保持单调；只在宽度为$S$的带$u\in[p_t,p_t+S)$内计算式\eqref{eq:nlp-speech-rnnt-joint}的完整联合网络及其损失。总损失是剪枝损失与简化损失的加权和，原文发现把简化损失乘以0.5效果最好，它同时训练用于选带的投影。

带宽$S$通常取个位数。上例取$S=5$、汉字词表时，输出张量从约4.9\,GB降到$16\times250\times5\times5001\times4$字节，约0.40\,GB。代价是选带依赖简化模型的质量：训练初期简化模型尚不可靠，带选得不准会使真正占用大的格点落在带外。开源实现icefall因此在最初几千步只用简化损失，随后以较小权重加入剪枝损失，再升到完整权重；带宽过窄也会影响收敛。

RNN-T同时保留了单调路径求和与输出历史，编码器做成流式后就能逐步输出，是目前在线识别的主流结构之一。它的代价集中在训练显存和实现复杂度上；预测网络只读文字，外部语言模型融合时还要考虑它已经学到的内部语言先验（第\ref{sec:nlp-speech-fusion}小节）。

\subsection{大规模弱监督识别：Whisper}
\label{sec:nlp-speech-whisper}

前述模型通常在一两个人工转写的数据集上训练。它们在同分布测试集上错误率很低，换到另一种口音、噪声或领域时却明显变差，而人工转写难以扩大到覆盖这些条件。Whisper换了一个方向：结构仍是标准的Transformer编码器—解码器，训练数据改为约68万小时从网络收集的音频与字幕配对\scite{radford2023whisper}。其中约11.7万小时覆盖英语以外的96种语言，约12.5万小时是其他语言音频配英语译文。这些字幕没有经过人工核对，属于弱监督；数据处理用启发式规则删去疑似机器生成的转写（例如全大写或全无标点的文本），并删去字幕语种与音频语种不符的样本。

模型输入是30秒窗口。音频重采样到16\,kHz，计算25\,ms窗长、10\,ms帧移的80维log-Mel，不足30秒的补零，得到$\bm{X}\in\mathbb{R}^{3000\times80}$。卷积前端由两层宽度为3的一维卷积组成，激活函数为GELU：第一层步长1，输出$3000\times d$；第二层步长2，输出$1500\times d$，每个位置对应20\,ms。加上正弦位置编码后进入预归一化的Transformer编码器，输出$1500\times d$。解码器使用可学习的位置编码，输入与输出词嵌入共享参数，每层通过交叉注意力（式\eqref{eq:nlp-llm-cross-attention}）读取编码器输出。每层的自注意力、前馈与残差结构与第\ref{sec:nlp-llm-decoder-layer}小节的解码器层相同，只是编码器不用因果掩码。模型从约3900万参数（4层，$d=384$）到约15.5亿参数（32层，$d=1280$）分为多档。与第\ref{sec:nlp-speech-encoder}小节Conformer常用的4倍二维卷积降采样相比，Whisper只降采样2倍，编码序列更长，但30秒的固定窗口限制了总长度。

同一段声音可以要求转写、译成英语或只判断语种，Whisper把任务写成解码器开头的特殊词元。一个带时间戳的中文转写目标形如
\begin{center}
\small
\texttt{<|startofprev|>} 机器学习读书会通知 \texttt{<|startoftranscript|><|zh|><|transcribe|>}\\
\texttt{<|0.00|>} 活动改到四月十八日 \texttt{<|2.40|><|2.40|>} 下午两点在海棠楼 \texttt{<|4.60|><|endoftext|>}
\end{center}
\texttt{<|startofprev|>}之后是前一窗口的转写，作为上下文；\texttt{<|zh|>}是语种，\texttt{<|transcribe|>}与\texttt{<|translate|>}区分转写和译成英语；每个文本片段前后各有一个时间词元，以20\,ms为间隔量化窗口内的相对时间。不需要时间时改用\texttt{<|notimestamps|>}；窗口内没有语音时，语种位置预测\texttt{<|nospeech|>}。训练损失是目标词元的条件负对数概率，与式\eqref{eq:nlp-speech-las}相同，但上下文部分不计损失。这样，语种识别、有无语音判断、转写、翻译和时间戳都成为同一个解码器的生成任务。

长录音按30秒窗口依次处理：用当前窗口最后一个完整片段的结束时间决定下一窗口的起点，把已解码文字作为下一窗口的上下文，局部时间戳加上窗口在录音中的偏移。原文还采用几项启发式：束宽为5的束搜索；输出的平均对数概率过低或文本压缩率过高（提示重复）时，提高采样温度重新解码；\texttt{<|nospeech|>}概率高且平均对数概率低时，把该窗口判为无语音。这些规则说明，固定窗口的自回归解码在长音频上需要额外的保护。

大规模数据带来的主要收益是跨领域的鲁棒性。原文比较了在LibriSpeech上错误率相近的模型：Whisper不在任何目标数据集上微调，在其他多个英语识别数据集上的平均相对错误率比在LibriSpeech上有监督训练的模型低约55\%。代价同样明显。第一，解码器本质上是以声音为条件的语言模型，声音证据弱时语言先验可能占上风：静音或背景音乐段可能生成“谢谢观看”一类网络字幕中常见的句子，或在窗口之间重复上一段文字\scite{nlp-speech-koenecke2024}。CTC与RNN-T的每个非空白输出都要在某个声学步上与空白竞争，没有声音支持时默认输出空白，这类无依据的长段插入少得多，但并非不可能。第二，固定30秒窗口、全局注意力和自回归解码都要求先拿到整段声音，Whisper不能直接流式运行，需要第\ref{sec:nlp-speech-whisper-streaming}小节的外部组织。第三，把任务词元设错，例如把转写设为翻译，输出语言就会改变，这不是听错了声音，评价时应与声学错误分开统计。

\subsection{流式识别}
\label{sec:nlp-speech-streaming}

在线识别要求每个单位在其声音结束后不久提交，并且已提交的结果不再频繁改动。延迟有两个来源：模型为了看见右侧声音而等待，以及计算本身所需的时间。本小节分两条路线。第一条从模型结构入手，限制编码器的右侧可见范围，使CTC或RNN-T能逐块输出。第二条处理Whisper这类整段模型，反复重新识别不断增长的缓冲，只提交多次假设一致的前缀。两者都要同时报告提交延迟、修订量和最终错误率，因为结果稳定不等于结果正确。

\subsubsection{分块注意力与流式编码器}
\label{sec:nlp-speech-chunk-attention}

自注意力编码器的每个位置默认读取整段序列，最直接的流式改造是用注意力掩码限制可见范围。若每层允许看右侧$r$步，$N$层叠加后最上层的位置实际依赖右侧$Nr$步：12层、每层2步、每步40\,ms时，前瞻达到960\,ms，延迟随层数累积。\term{分块注意力}（chunk attention）改为把编码序列切成长为$C$的块，每个位置可以看本块全部位置和左侧$L$个历史块。以从1开始的位置索引，记$c(i)=\lfloor(i-1)/C\rfloor$为位置$i$所在的块号，掩码为
\begin{equation}
 M_{i,j}=\begin{cases}
   0, & c(i)-L\leq c(j)\leq c(i),\\
   -\infty, & \text{其他},
 \end{cases}
 \qquad
 \operatorname{Attn}=\operatorname{softmax}\!\left(\frac{\bm{Q}\bm{K}^{\top}}{\sqrt{d_k}}+\bm{M}\right)\bm{V}.
 \label{eq:nlp-speech-chunk-mask}
\end{equation}
图\ref{fig:nlp-speech-chunk-mask}画出了$C=2$、$L=1$时的掩码。由于每层的可见范围都止于本块末尾，叠加多少层，右侧都不会越过块边界。块内第一个位置要等到块末尾才能计算，最多等待$C-1$步，块大小因而直接决定了模型带来的延迟下界：$C=16$个40\,ms编码步时，最长等待约640\,ms，与层数无关。

\begin{figure}[htbp]
  \centering
  % 第九卷 15.2.7.1 节：分块注意力掩码（正文栏宽，约60 mm）。
% 8个编码步，块大小C=2，左侧保留L=1个历史块。蓝色格为可见的键位置。
\begin{tikzpicture}[x=1mm,y=1mm,figure picture]
  \foreach \i in {0,...,7} {
    \foreach \j in {0,...,7} {
      \pgfmathtruncatemacro{\ci}{floor(\i/2)}
      \pgfmathtruncatemacro{\cj}{floor(\j/2)}
      \pgfmathtruncatemacro{\vis}{and(\cj<=\ci,\cj>=\ci-1)}
      \ifnum\vis=1
        \path[fill=FigureInputFill,draw=FigureInput,line width=.4pt]
          ({5*\j},{-5*\i}) rectangle ++(5,-5);
      \else
        \path[fill=white,draw=FigureGrid,line width=.4pt]
          ({5*\j},{-5*\i}) rectangle ++(5,-5);
      \fi
    }
    \pgfmathtruncatemacro{\k}{\i+1}
    \node[font=\figurefont\scriptsize,anchor=east] at (-1,{-5*\i-2.5}) {$\k$};
    \node[font=\figurefont\scriptsize,anchor=south] at ({5*\i+2.5},1) {$\k$};
  }
  \foreach \k in {0,2,4,6,8} {
    \draw[FigureStroke,line width=.8pt] ({5*\k},0) -- ({5*\k},-40);
    \draw[FigureStroke,line width=.8pt] (0,{-5*\k}) -- (40,{-5*\k});
  }
  \node[font=\figurefont\scriptsize,anchor=south] at (20,6) {键位置$j$};
  \node[font=\figurefont\scriptsize,rotate=90,anchor=south] at (-6,-20) {查询位置$i$};
  \node[font=\figurefont\scriptsize,anchor=west,align=left,text=FigureMuted]
    at (43,-20) {块大小$C=2$\\左侧历史$L=1$块\\粗线为块边界};
\end{tikzpicture}

  \caption{分块注意力掩码，8个编码步，块大小$C=2$，左侧保留$L=1$个历史块。第$i$行的蓝色格表示查询位置$i$可以读取的键位置。每个位置都能看到本块内右侧的位置，但不能越过块边界；各层使用同一掩码，叠加后的可见范围仍止于块末尾。}
  \label{fig:nlp-speech-chunk-mask}
\end{figure}

块大小是准确率与延迟的权衡，而不同应用的延迟预算不同。U2的\term{动态块训练}（dynamic chunk training）为每个批次随机选择块大小：以一半概率使用整句（全上下文），否则从1至25个编码步中均匀抽取，使同一个模型在推断时可以按延迟预算选择块大小\scite{nlp-speech-zhang2020u2}。该工作还把它与第\ref{sec:nlp-speech-joint-ctc-attention}小节的联合模型结合成两遍解码：CTC分支在分块编码器上流式输出暂定结果，整句结束后由注意力解码器对CTC的$n$个最佳候选重打分，作为最终结果。

Conformer中的其他模块也要随之改造。卷积模块的深度卷积若两侧补零，会读到右侧未来的帧，需改为只在左侧补零的因果卷积；跨块时缓存上一块末尾的$k-1$帧（$k$为卷积核宽度），作为本块卷积的左侧输入。注意力同样缓存左侧历史块的键与值，与第\ref{sec:nlp-inference-cache}小节文本生成中的键值缓存作用相同。批归一化在推断时使用训练中累积的统计量，本身不读未来帧，但分块训练与整句训练的统计差异可能造成不一致，一些实现改用层归一化。CTC与RNN-T接上这种编码器即可逐块输出；AED的全局注意力解码不能直接这样改造，这是流式场景偏向CTC与RNN-T的主要原因。

\subsubsection{稳定前缀提交：Whisper-Streaming}
\label{sec:nlp-speech-whisper-streaming}

整段识别器也可以通过缓冲和反复识别接入实时声音，难点是决定哪些结果已经可以提交。Whisper-Streaming采用\term{局部一致性}（LocalAgreement）策略：对不断增长的声音缓冲反复识别，比较连续两次的结果，只提交它们一致的公共前缀，其余部分留待更多声音到来后再确认\scite{nlp-speech-machacek2023}。这是模型之外的调度规则，Whisper本身的概率模型并不包含它；其思路与第\ref{sec:nlp-generation-adaptive-policy}小节同声翻译中的自适应读写策略相似。

每轮收集到至少规定长度的新声音后，识别当前缓冲；跳过已经提交的文字，对上次与本次的剩余部分求最长公共前缀，扩展提交边界。已提交部分仍留在缓冲内为识别提供上下文，但它在新识别中的变化不再对外撤回。缓冲过长时，在已提交的句末位置裁掉旧声音，并把已提交文字作为下一轮的提示。若一轮计算慢于新声音的积累，下一轮会处理更多声音，实际等待由输入速度与计算速度共同决定。

设三次教学假设依次为“改到／四月／十七日”“改到／四月／十八日”“改到／四月／十八日／下午”。第一次没有可比较的结果；第二次提交“改到／四月”，日期两次不同而保留；第三次才提交“十八日”。若第一次的“十七日”被直接播放给听者，后来修改屏幕文字也无法撤销已经听到的日期。反过来，若连续两次都误识为“十七日”，局部一致性同样会提交错误，它判断的是稳定性而不是正确性。评价时固定块长、缓冲裁剪和计分规范，分别统计最终字符错误率、暂定文字的修订量和各单位的提交延迟；被漏识的字没有提交时刻，须同时报告遗漏，不能只平均已识别字的延迟。

\subsection{识别模型的比较}
\label{sec:nlp-speech-asr-comparison}

表\ref{tab:nlp-speech-asr-comparison}汇总了本节的六类识别模型。各列依次回答：训练需要什么标签；输出之间是否假设条件独立；对齐是否被约束为单调；能否在声音到达时逐步输出；外部语言模型怎样接入；能否直接得到每个单位的时间位置；主要的计算或工程代价是什么。

\begin{table*}[htbp]
  \centering
  \footnotesize
  \begin{tabularx}{\linewidth}{@{}>{\raggedright\arraybackslash}p{17mm}>{\raggedright\arraybackslash}p{16mm}>{\raggedright\arraybackslash}p{13mm}>{\raggedright\arraybackslash}p{14mm}>{\raggedright\arraybackslash}p{17mm}>{\raggedright\arraybackslash}p{19mm}>{\raggedright\arraybackslash}p{17mm}>{\raggedright\arraybackslash}X@{}}
    \toprule
    方法 & 训练标签 & 输出条件独立 & 单调对齐 & 流式 & 外部语言模型 & 对齐与时间 & 主要代价\\
    \midrule
    混合HMM--DNN & 帧级状态对齐（来自HMM--GMM） & 是，给定状态各帧独立 & 是，由HMM拓扑保证 & 可以，受网络输入窗口限制 & 解码图中的$G$，可词格重打分 & 状态级对齐直接可得 & 多阶段训练，需发音词典\\
    CTC & 整句文字 & 是 & 是 & 配合分块或因果编码器 & 前缀束搜索中按字加分 & 尖峰位置，偏离发音边界 & 需外部语言模型补历史\\
    AED（LAS） & 整句文字 & 否 & 否，注意力可跳读 & 难，需整段编码 & 束搜索中浅层融合 & 注意力权重，不可靠 & 跳读、重复与过早结束\\
    联合CTC/\allowbreak 注意力 & 整句文字 & 注意力分支否 & CTC分支约束 & 两遍解码时首遍可流式 & 联合分数中加一项 & 可由CTC分支取得 & 两个分支都要计算\\
    RNN-T & 整句文字 & 否 & 是 & 配合分块或因果编码器 & 浅层融合，宜扣除内部语言先验 & 发出位置，常有延迟 & $B\times T\times U\times|V|$显存\\
    Whisper & 弱监督字幕，约68万小时 & 否 & 否 & 不能直接流式，需外部调度 & 可融合，但解码器已含强语言先验 & 20\,ms时间词元，片段级 & 无语音段幻觉，30秒固定窗口\\
    \bottomrule
  \end{tabularx}
  \caption{六类语音识别模型的对照。“流式”指在声音到达过程中逐步给出稳定输出的能力，取决于编码器可见范围与解码方式，不由损失函数单独决定；“对齐与时间”指不另做强制对齐时能直接取得的时间信息。}
  \label{tab:nlp-speech-asr-comparison}
\end{table*}

按表中的列可以读出各方法依次解决了什么问题。CTC用路径求和去掉了帧级对齐标签与多阶段训练；AED引入输出历史，却丢掉了单调约束；联合训练用CTC分支把单调约束找回来；RNN-T在一个模型里同时得到输出历史、单调对齐和流式能力，代价转移到训练显存；Whisper不改变结构，而是用弱监督数据扩大覆盖，换来鲁棒性，同时带来幻觉与非流式的问题。选型取决于延迟预算、可用数据与部署资源：需要低延迟在线输出时，分块编码器上的RNN-T或CTC更合适；离线转写、语种多样且没有领域内标注时，大规模弱监督模型更合适；需要精确时间边界或只有少量数据时，HMM路线仍有用处。没有一种方法在所有条件下都占优。

\section{识别结果的解码、分段与整理}
\label{sec:nlp-speech-postprocessing}

识别器输出的字符串还不能直接交付。声音含混时，需要借助大量文本中的语言知识在候选之间选择；字幕和数据清洗需要每个词的时间位置；会议录音要分清哪段话是谁说的；书面使用又需要标点和规范的数字写法。本节依次处理这四项，说明各自的输入、监督来源与评价口径。四项整理都可能引入新的错误，评价时应与识别本身的错误分开报告。

\subsection{语言模型浅层融合与候选重排}
\label{sec:nlp-speech-fusion}

端到端识别器只从配对数据中学到语言规律。一千小时转写大约只有一千万字的文本，而纯文本语料可以多出几个数量级，独立训练的语言模型能补充大量词汇与搭配知识，例如在“海棠楼”与发音相近的候选之间作出选择。\term{浅层融合}（shallow fusion）在解码搜索中把识别分数与语言模型分数相加\scite{nlp-speech-kannan2018}：
\begin{equation}
 S(y)=\log p_{\mathrm{asr}}(y\mid\bm{X})
       +\lambda\log p_{\mathrm{lm}}(y)+\gamma |y|.
 \label{eq:nlp-speech-fusion}
\end{equation}
$\lambda$为语言模型权重，$\gamma$为长度奖励，用来抵消语言模型偏好短文本的倾向；两者在独立开发集上按识别错误率选定。对逐字生成的AED与RNN-T，这一分数在束搜索中逐步累加：每输出一个非空白符号$y_u$，就加上$\log p_{\mathrm{asr}}(y_u\mid\cdot)+\lambda\log p_{\mathrm{lm}}(y_u\mid y_{<u})+\gamma$。CTC中多条路径对应同一文本，要先按概率合并同文本路径，再对该文本计一次语言模型分数；第\ref{sec:nlp-speech-ctc-decoding}小节已经给出在前缀束搜索中实现这一点的位置。

教学候选甲、乙的识别对数分数为$-2.0$与$-2.4$，语言模型对数分数为$-3.0$与$-1.0$，长度相同。取$\lambda=0.3$时总分分别为$-2.9$与$-2.7$，乙胜出。语言先验改变了选择，但它只说明乙在文本中更常见，并不能证明录音中说的就是乙。专业词表也可以用受控的加分或词典约束接入，代价是可能把发音相近的其他词吸附到常用专名上。

识别模型自身也学到了一个语言先验，浅层融合会把它与外部语言模型的先验重复计入。这一问题在RNN-T中尤为明显，因为预测网络本身就是一个文字模型。\term{内部语言模型估计}（internal language model estimation）先估计识别模型在去掉声学输入时给出的文字概率$p_{\mathrm{ilm}}(y)$，融合时减去它：$S(y)=\log p_{\mathrm{asr}}-\mu\log p_{\mathrm{ilm}}+\lambda\log p_{\mathrm{lm}}$，在识别器训练领域与目标领域不同时尤其有用\scite{nlp-speech-meng2021}。深度融合等方法则把语言模型的隐藏状态接入解码器再联合训练，不在本节展开。

候选重排把搜索分为两遍：第一遍用识别器和较小的语言模型产生$n$个完整候选，HMM系统则输出词格；第二遍用更大的语言模型对这些候选重新打分，再按式\eqref{eq:nlp-speech-fusion}的同类组合选出结果。第二遍可以使用计算量大的神经语言模型，因为只需对有限个候选打分；第\ref{sec:nlp-inference-reranking}小节的一般重排方法同样适用。限制在于它只能在已有集合中选择：若“海棠楼”在第一遍已被剪掉，第二遍的语言模型再强也找不回来。比较两种方法时应固定候选数量，并区分错误是候选中缺少正确答案，还是有正确答案却选错了。

\subsection{给定文本的CTC与HMM强制对齐}
\label{sec:nlp-speech-forced-alignment}

字幕制作、语音数据清洗和发音分析都需要知道已知文本中的每个字或词在录音中的位置。\term{强制对齐}（forced alignment）固定文本，只搜索它与声音的最佳对应，不需要重新训练声学模型。

CTC对齐复用第\ref{sec:nlp-speech-ctc-forward-backward}小节的扩展标签。把式\eqref{eq:nlp-speech-ctc-forward}中的求和换成取最大，并记录取得最大值的前驱：
\begin{equation}
 v_{t,s}=q_t(z_s)\max\{v_{t-1,s},\ v_{t-1,s-1},\ \delta_sv_{t-1,s-2}\},
 \label{eq:nlp-speech-ctc-viterbi}
\end{equation}
在两个合法终态$2U$与$2U+1$中取较大者回溯，就得到每一步所处的扩展标签位置。HMM对齐则按词典把给定词序列展开成音素HMM的串联，词间插入可选的静音状态，在这张图上运行式\eqref{eq:nlp-speech-hmm-viterbi}的Viterbi算法。蒙特利尔强制对齐器就是基于Kaldi的HMM--GMM对齐工具\scite{nlp-speech-mcauliffe2017}。

从路径到时间还需要换算。设编码步间隔为40\,ms、起点偏移为0，回溯路径为$(\varnothing,a,a,\varnothing,b,b)$。若约定以占用网格单元的外边界表示区间，可报告$a:[0.04,0.12)$秒、$b:[0.16,0.24)$秒；若改用窗口中心，或把相邻空白分配给左右的字，会得到不同的边界。特征窗长、卷积感受野和降采样带来的偏移也要计入换算。

两条路线的边界性质不同。CTC训练后的尖峰现象（第\ref{sec:nlp-speech-ctc-gradient}小节）使每个字通常只占一两个编码步，其余时间归给空白，所以CTC对齐得到的更接近“在何时确认了这个字”，而不是这个音从哪里开始到哪里结束；单向模型的尖峰还可能整体滞后。HMM的状态要覆盖音素的全部帧，边界通常更贴近发音，分辨率也是10\,ms的特征帧而非降采样后的编码步。目标含重复字符时，CTC对齐必须保留中间空白；停顿则由空白或HMM静音状态吸收。

强制对齐总会给出一条路径，即使文本本身有错。给错的文字可能对上一条分数较低的近似路径，因此对齐成功不能证明转写正确。数据清洗中常按归一化后的路径分数和异常时长筛出可疑样本，再人工核对。

\subsection{说话者嵌入与聚类：x-vector}
\label{sec:nlp-speech-xvector}

会议转写要回答“谁在何时说话”，说话者验证要判断两段声音是否出自同一人。两者都需要把变长的声音片段映射为便于比较的固定长度向量，即\term{说话者嵌入}（speaker embedding），希望同一说话者的片段彼此接近、不同说话者的片段相距较远。x-vector用说话者分类任务学习这种嵌入\scite{nlp-speech-snyder2018}。

帧级部分采用\term{时延神经网络}（time-delay neural network, TDNN）。一个TDNN层在每个时刻$t$读取下层若干个偏移位置的输出，拼接后做仿射变换：
\begin{equation}
 \bm{h}^{(l)}_t=\operatorname{ReLU}\bigl(\bm{W}^{(l)}[\bm{h}^{(l-1)}_{t+o_1};\ldots;\bm{h}^{(l-1)}_{t+o_K}]+\bm{b}^{(l)}\bigr),
 \label{eq:nlp-speech-tdnn}
\end{equation}
其中$o_1,\ldots,o_K$是该层的时间偏移。它等价于一维扩张卷积：偏移集合$\{-2,0,2\}$即卷积核宽3、扩张2。逐层叠加后，上层每个输出覆盖的帧数不断扩大。表\ref{tab:nlp-speech-xvector}列出原文的结构：输入为24维滤波器组特征，前三层把上下文扩大到15帧，后两层逐帧变换，第五层输出1500维。

\begin{table}[htbp]
  \centering
  \small
  \begin{tabular}{@{}llrr@{}}
    \toprule
    层 & 时间偏移 & 累计上下文 & 输入$\times$输出\\
    \midrule
    帧层1 & $\{-2,\ldots,2\}$ & 5 & $120\times512$\\
    帧层2 & $\{-2,0,2\}$ & 9 & $1536\times512$\\
    帧层3 & $\{-3,0,3\}$ & 15 & $1536\times512$\\
    帧层4、5 & $\{0\}$ & 15 & $512\times512$，$512\times1500$\\
    统计池化 & 全部$T$帧 & $T$ & $1500T\times3000$\\
    片段层6、7 & --- & $T$ & $3000\times512$，$512\times512$\\
    softmax & --- & $T$ & $512\times N$\\
    \bottomrule
  \end{tabular}
  \caption{x-vector网络结构，按Snyder等的原文整理。$T$为片段帧数，$N$为训练说话者数。帧层1的输入是5帧24维特征的拼接；x-vector取自片段层6的仿射输出。}
  \label{tab:nlp-speech-xvector}
\end{table}

\term{统计池化}（statistics pooling）把变长的帧级输出$\bm{h}_1,\ldots,\bm{h}_T\in\mathbb{R}^{1500}$汇总为固定长度：
\begin{equation}
 \bm{\mu}=\frac1T\sum_{t=1}^{T}\bm{h}_t,\qquad
 \bm{\sigma}=\Bigl(\frac1T\sum_{t=1}^{T}\bm{h}_t\odot\bm{h}_t-\bm{\mu}\odot\bm{\mu}\Bigr)^{\odot1/2},\qquad
 \bm{s}=[\bm{\mu};\bm{\sigma}]\in\mathbb{R}^{3000}.
 \label{eq:nlp-speech-stats-pooling}
\end{equation}
均值描述该说话者帧特征的平均位置，标准差描述它们的变化范围；两者都与片段长度无关。$\bm{s}$经两个片段级层后，对训练集中的$N$个说话者做softmax分类，损失为交叉熵。训练时还向录音中加入噪声与混响，这是原文提升鲁棒性的关键做法之一。使用时丢弃分类层，取片段层6的仿射输出作为嵌入；测试说话者不必出现在训练集中，因为比较的是嵌入而不是类别。后续的ECAPA-TDNN在同一框架下加入通道注意力、多层特征聚合和注意力加权的统计池化\scite{nlp-speech-desplanques2020}。

比较两个嵌入$\bm{e}_1,\bm{e}_2$最简单的方法是余弦相似度。更常用的是\term{概率线性判别分析}（probabilistic linear discriminant analysis, PLDA）\scite{nlp-speech-prince2007}。它把（经中心化、降维和长度归一化的）嵌入写成$\bm{e}=\bm{m}+\bm{\Phi}\bm{\beta}+\bm{\epsilon}$：$\bm{\beta}\sim\mathcal{N}(\bm{0},\bm{I})$是说话者共享的隐变量，$\bm{\epsilon}\sim\mathcal{N}(\bm{0},\bm{\Sigma}_w)$是同一说话者不同片段之间的变化。记说话者间协方差$\bm{\Sigma}_b=\bm{\Phi}\bm{\Phi}^{\top}$、总协方差$\bm{\Sigma}_t=\bm{\Sigma}_b+\bm{\Sigma}_w$。两段来自同一人时共享$\bm{\beta}$，来自不同人时独立，于是评分取对数似然比
\begin{equation}
 \mathrm{LLR}=\log\mathcal{N}\!\left(\begin{bmatrix}\bm{e}_1\\\bm{e}_2\end{bmatrix};
   \begin{bmatrix}\bm{m}\\\bm{m}\end{bmatrix},
   \begin{bmatrix}\bm{\Sigma}_t&\bm{\Sigma}_b\\\bm{\Sigma}_b&\bm{\Sigma}_t\end{bmatrix}\right)
 -\log\mathcal{N}\!\left(\begin{bmatrix}\bm{e}_1\\\bm{e}_2\end{bmatrix};
   \begin{bmatrix}\bm{m}\\\bm{m}\end{bmatrix},
   \begin{bmatrix}\bm{\Sigma}_t&\bm{0}\\\bm{0}&\bm{\Sigma}_t\end{bmatrix}\right).
 \label{eq:nlp-speech-plda}
\end{equation}
参数用带说话者标签的训练嵌入通过EM估计。与余弦相比，PLDA区分了“说话者之间的差异”与“同一人内部的变化”两种方向，后者在打分时被打折扣。

说话者验证用\term{等错误率}（equal error rate, EER）评价。给定阈值$\theta$，误接受率是不同人的配对中得分不低于$\theta$的比例，误拒绝率是同一人的配对中得分低于$\theta$的比例；阈值升高时前者下降、后者上升，两者相等时的值即EER。教学例子中，同一人配对的得分为$0.9,0.8,0.6,0.4$，不同人配对为$0.7,0.3,0.2,0.1$。取$\theta=0.5$，误拒绝只有$0.4$一项，误接受只有$0.7$一项，两者都是25\%，EER为25\%。

\term{说话者日志}（speaker diarization）把这些部件串起来。先由语音活动检测取得有声片段，再按固定长度（例如1.5秒、重叠一半）切成小段，提取并归一化嵌入，用凝聚层次聚类分组：从每段一组开始，反复合并相似度最高的两组，相似度低于开发集选定的阈值或达到已知人数时停止。设三段的相似度为$s_{1,2}=0.91$、$s_{1,3}=0.18$、$s_{2,3}=0.22$，用组间平均相似度、阈值$0.7$：先合并1、2，合并组与3的相似度为$(0.18+0.22)/2=0.20$，低于阈值，最终得到两组。把各组映射回起止时刻，就得到匿名的说话者A、B。评价用日志错误率$\mathrm{DER}=(T_{\mathrm{miss}}+T_{\mathrm{fa}}+T_{\mathrm{conf}})/T_{\mathrm{ref}}$，三项依次是漏检、虚警与身份混淆的时长，计算前须先求匿名标签与参考身份的最佳匹配，并约定交叠语音与边界容差的口径。

这条流水线的局限主要有三点。片段太短时统计池化不稳定，嵌入噪声大；两人同时说话时一个嵌入混合了两种身份，聚类无法给出重叠标签，需要额外的分离或重分段机制；聚类只给出匿名分组，要把A对应到某位真实用户，还需要已登记的声纹或其他证据。

\subsection{标点恢复与反向文本规范化}
\label{sec:nlp-speech-itn}

识别输出通常没有标点，数字也以口头形式出现，例如“明天可以吗不可以后天”。整理文字需要两步：恢复标点和句界，把口头形式改写为书面形式。

\term{标点恢复}（punctuation restoration）最常见的做法是逐词元标注：对识别输出的每个词元预测其后应接的标点，标签集如$\{$无，逗号，句号，问号$\}$。训练数据很容易构造，只需把带标点的书面文本去掉标点作为输入，原标点的位置与类别作为标签。早期模型用双向循环网络加注意力\scite{nlp-speech-tilk2016}；现在多用BERT类编码器加逐词元分类头（第\ref{sec:nlp-classical-bert-heads}小节），损失为逐位置交叉熵，因“无”类远多于其他类，常对标点类加权。是否加标点在很大程度上取决于后文：“明天可以吗不可以后天”既可以是“明天可以吗？不可以，后天。”，也可以是“明天可以吗？不，可以后天。”，句界改变了语义。在线场景只能看到有限的后文，通常规定右侧上下文$r$个词元，位置$i$的标点要等到第$i+r$个词元出现后才确定；只见当前前缀的在线模型不能与看见全文的离线模型直接比较。

\term{反向文本规范化}（inverse text normalization, ITN）把读出的形式改写为规定的书面形式，如“四月十八号下午两点”改为“4月18日14:00”。传统做法用WFST规则文法分两步完成：标注器先找出日期、时间、数量、金额等类型的片段并解析出字段，例如把上句解析为$\texttt{date}\{\texttt{month}{:}4,\texttt{day}{:}18\}$与$\texttt{time}\{\texttt{hour}{:}14,\texttt{minute}{:}00\}$；输出器再按规定格式把字段写成书面形式。两步都是第\ref{sec:nlp-speech-hybrid}小节意义下的加权转导器，规则可以审查，输出可控\scite{nlp-speech-zhang2021itn}。神经做法把ITN看作序列转换：一种是给每个口头词元预测一个改写操作，再由少量规则执行\scite{nlp-speech-pusateri2017}；另一种是用编码器—解码器直接生成书面形式。神经模型覆盖面更广，但可能生成规则不允许的格式，因此常用规则文法检查或约束其输出。

类型判断决定了逆映射。“二零一”可能是房间号201，也可能是一串编号，不能无条件按数量处理；“四月十六日十八点前”中的“前”改变了时间条件，必须保留。ITN通常不是一一可逆的：同一口头形式可能对应不同的简称、大小写或格式，需要上下文和规定的输出风格。它与第\ref{sec:nlp-speech-tts}节语音合成文本前端中的正向规范化方向相反，但两者并不对称：正向规范化把一种书面形式展开为一种读法，反向规范化要从口头形式中判断原本是哪一种书面形式。整理只改变允许的书面形式，不应凭语言流畅性把听成的“十七日”改为“十八日”。评价时既保留统一规范下的识别错误，也单独统计整理后的数量、单位、句界和格式错误，以免规范化掩盖内容差异。

\section{口语理解}
\label{sec:nlp-speech-understanding}

用户说“后天三个人”，系统要得到的是“修改预约，日期某天，人数3”这样的语义，而不一定是完整转写。\term{口语理解}（spoken language understanding, SLU）从声音及允许的对话历史恢复意图、参数与交流作用。第\ref{sec:nlp-dialogue-intent-crf}小节已在文字上解决了意图分类与槽位标注；换成声音输入后，又多出四项挑战。其一是识别错误传播：识别器把“后天”听成“后台”，后面的文本模型再准确也找不到日期。其二是口语现象：更正、重复和省略在书面语中少见，在口语里却很常见。其三是语境依赖：同一句“可以”在不同历史和语气下作用不同。其四是监督稀缺：同时带声音和语义标注的数据远少于纯文本标注。

方法谱系按“语义模型看到什么输入”推进。最简单的流水线把识别一最佳结果交给第\ref{chap:nlp-question-answering-dialogue}章的文本理解模型；为缓解错误传播，第二类方法把识别的多个候选或词格交给理解模型，让语义证据参与选择文字；第三类端到端模型直接读取声学编码，用语义目标训练，跳过文字接口；最后再把对话历史和韵律加入输入，判断交流作用。SLURP把声音、转写与场景、动作、实体标注配对，可用于分别检查识别误差与语义误差；它为评价组织提供依据，并不是下面各算法的出处\scite{nlp-speech-bastianelli2020}。全节的评价都把单个字段正确与完整记录正确分开：槽位F1高不等于每条请求都被正确执行。

\begin{table*}[htbp]
  \centering
  \begin{tabularx}{\linewidth}{@{}>{\raggedright\arraybackslash}p{28mm}>{\raggedright\arraybackslash}X>{\raggedright\arraybackslash}X>{\raggedright\arraybackslash}X@{}}
    \toprule
    方法 & 解决的问题 & 核心机制 & 新代价\\
    \midrule
    一最佳流水线 & 复用文本意图与槽位模型 & 识别一最佳$\to$联合意图—CRF模型 & 识别错误直接传入语义\\
    候选重排与词格输入 & 一最佳丢失正确文字 & 识别分数与语义分数联合选择；理解模型读取带后验的词格 & 需要识别器导出候选；计算随候选数增长\\
    直接声音分类 & 语义只需少量类别时跳过转写 & 自监督编码器$\to$池化$\to$softmax & 不给出参数；需声音—意图配对数据\\
    语义记录生成 & 参数、更正与撤回需要结构化输出 & 声学编码器$\to$交叉注意力解码器$\to$受约束记录 & 配对标注稀缺；证据定位需另设监督\\
    历史与韵律条件分类 & 同一文字的交流作用不同 & 拼接历史编码；韵律统计或声学表示与文本融合 & 标注一致性低；需按说话者和会话留出\\
    \bottomrule
  \end{tabularx}
  \caption{口语理解方法对照。自上而下，语义模型看到的输入从单一文字扩展到候选集合、声学序列，再到历史与韵律；每一步解决前一步丢失的信息，也带来新的数据或计算需求。}
  \label{tab:nlp-speech-slu-methods}
\end{table*}
\FloatBarrier

\subsection{意图、槽位与口语事实}
\label{sec:nlp-speech-semantics}

本小节另设一个与读书会无关的预约请求：2025年4月17日，用户说“明天两个人，改成后天三个人”。当前有效值是日期4月19日、人数3；4月18日与人数2是被撤回的旧值，需要另行记录，不能和有效值并列。意图说明请求做什么，槽位记录参数；每个槽位再分别保存规范值和证据区间，前者供业务使用，后者供核查。下面先介绍转写后的联合模型及其口语侧补救，再介绍直接从声音预测语义的模型。

\subsubsection{意图分类与CRF槽位标注}
\label{sec:nlp-speech-intent-crf}

文本侧的模型可以直接沿用。第\ref{sec:nlp-dialogue-jointbert}小节的联合模型让一个预训练文本编码器同时服务两个头：句级表示接softmax预测意图$c$，逐位置表示给BIO标签打分；在槽位头上加线性链CRF，可学习标签转移并排除句首\code{I-DATE}一类非法路径。联合目标为
\[
  L=-\log p(c\mid y)-\lambda\log p_{\mathrm{crf}}(z\mid y),
\]
其中$y$是输入文字，$z$是BIO标签序列，$\lambda$由开发集选定。CRF的定义、梯度与Viterbi解码见第\ref{sec:nlp-classical-crf}小节，合法转移掩码见第\ref{sec:nlp-analysis-paradigm-tagging}小节。对“明天／两／个／人／改成／后天／三／个／人”，模型标出两个日期跨度和两个人数跨度；“改成”触发更正关系，把后一组设为有效值、前一组标为撤回。日期按咨询时点解析，“后天”才得到4月19日。

口语侧的新问题在于$y$不再是人工写下的文字，而是识别结果。设识别器输出“明天两个人改成后台三个人”，CRF完全合法地给“后台”标O，日期槽位就缺失了。理解模型在人工转写上训练，从未见过这类替换，错误被原样传播。要定位误差来源，可以把同一语义模型分别接三种输入：人工转写给出语义模型自身的上限；识别一最佳给出流水线的实际表现；两者之差就是识别误差造成的损失。

第三种输入是识别候选集合，让语义证据参与选择文字。设识别器输出$N$个候选$y^{(1)},\ldots,y^{(N)}$，识别对数概率为$\log p_{\mathrm{asr}}(y^{(k)}\mid\bm{X})$，其中$\bm{X}$是声学特征序列。联合重排在候选与语义结果上同时取最大：
\begin{equation}
  (\hat k,\hat c,\hat z)=\argmax_{k,\,c,\,z}\;
  \mu\log p_{\mathrm{asr}}(y^{(k)}\mid\bm{X})
  +\log p(c,z\mid y^{(k)}),
  \label{eq:nlp-speech-slu-rerank}
\end{equation}
$\mu$平衡两种分数。教学示例取三条候选，识别对数概率为$-1.0$（“……后台三个人”）、$-1.4$（“……后天三个人”）、$-2.0$（“……后天四个人”）；语义模型对各自最优记录的对数概率为$-3.0$（日期槽缺失，记录不完整）、$-0.5$、$-0.6$。取$\mu=1$，三者总分为$-4.0$、$-1.9$、$-2.6$，重排选中第二条，而一最佳是第一条。这组数只用于说明计算；候选里有正确文字也不等于一定被选中，评价时要分别统计“候选含正确答案”与“最终选中正确答案”的比例。

候选列表仍有冗余：$N$条候选往往只在一两个词上不同。识别器解码时内部保留的\term{识别词格}（ASR lattice）把这些候选压缩成有向无环图，弧上是词及其后验概率；与第\ref{sec:nlp-analysis-pretrained-lattice}小节的中文分词词格结构相同，只是弧的来源是声学解码而不是词典。LatticeRNN把循环单元推广到词格上，每个节点汇合所有入弧的状态，并按弧后验加权\scite{nlp-speech-ladhak2016}；Transformer编码器也可以用可达关系构造注意力掩码，让每条弧只关注与它同在某条路径上的弧。词格输入保留了全部候选，计算量只随弧数增长；代价是需要识别器导出词格，且训练数据要用同一识别器解码得到，换识别器后分布会变化。另一种更简单的补救是直接在识别输出上训练语义模型，让它见过真实的识别错误。

无论接哪种输入，验收都不能止于槽位F1。标签序列完全合法时，“后天”仍可能被规范为4月18日，或“两个人”没有被撤回。应分别核查意图、跨度、规范值、更正关系和完整记录；SLURP为此还提出了容忍轻微转写差异的实体评分，用于把识别拼写误差与语义误差分开\cite{nlp-speech-bastianelli2020}。完整记录正确率要求一条请求的所有字段同时正确，与业务结果最接近。

\subsubsection{直接声音分类与语义记录生成}
\label{sec:nlp-speech-direct-slu}

流水线中识别器按逐字错误训练，把“后天”错成“后台”和把“的”错成“得”付出同样代价，而对语义而言前者是致命错误。端到端口语理解改用语义目标直接训练读取声音的模型，让对语义重要的声学差异得到更多关注。Serdyuk等最早演示了不经文字中间表示、直接从声音预测意图的模型\scite{nlp-speech-serdyuk2018}；在自监督语音表示出现后，这类模型通常以预训练编码器为主干。

先看只需判断意图的情形。设录音4秒，经第\ref{sec:nlp-speech-ssl}小节的自监督编码器（如wav2vec 2.0或HuBERT）得到每20毫秒一帧的表示，共$T=200$帧。按该小节末的使用方式，对编码器各层输出学一组softmax权重加权求和，得到$\bm{H}\in\mathbb{R}^{T\times d}$，$d=768$。意图是整句的属性，还要把$T$帧汇成一个向量，常用注意力池化：
\begin{equation}
  \bm{h}_{\mathrm{audio}}=\sum_{t=1}^{T}a_t\bm{h}_t,
  \qquad
  a_t=\frac{m_t\exp(\bm{q}^{\top}\bm{h}_t)}{\sum_{t'}m_{t'}\exp(\bm{q}^{\top}\bm{h}_{t'})}.
  \label{eq:nlp-speech-slu-pooling}
\end{equation}
$\bm{h}_t$是$\bm{H}$的第$t$行；$m_t\in\{0,1\}$是有效帧掩码，排除批内补齐的帧；$\bm{q}\in\mathbb{R}^{d}$是可学习查询向量。所有$a_t$相等时注意力池化退化为掩码均值池化。得到的$\bm{h}_{\mathrm{audio}}\in\mathbb{R}^{768}$经线性层和softmax输出$|\mathcal{C}|$个意图的概率，以交叉熵训练。数据较少时冻结编码器、只训练层权重、查询向量与分类头；数据较多时整体微调。

日期和人数还需要逐字段输出，单个池化向量装不下这些信息。语义记录生成保留整条声学序列$\bm{H}\in\mathbb{R}^{T\times d}$作为编码，接一个Transformer解码器，经式\eqref{eq:nlp-llm-cross-attention}的交叉注意力读取声音，逐词元生成序列化的语义记录$r$，以$-\sum_j\log p(r_j\mid\bm{H},r_{<j})$训练。上述请求的目标记录可以是
\begin{quote}
\code{intent=修改预约; date=2025-04-19; party=3;}\\
\code{retract=(date:2025-04-18, party:2)}
\end{quote}
解码时用第\ref{sec:nlp-inference-grammar}小节的文法约束只允许合法字段名、类型和结束位置，生成后再解析检查必填字段、日期范围与更正状态。一种常用变体让解码器先输出转写、再输出记录，并在编码器上加CTC辅助损失；这样仍保留文字用于核查，但转写与语义在同一模型中联合优化，语义部分可以利用声学证据纠正转写，而不是只能接受一最佳。

端到端模型的瓶颈是数据。声音—语义配对标注昂贵，常见的补救有两条：用预训练的自监督或识别编码器提供声学知识，只用少量配对数据学习语义头；用语音合成把纯文本语义数据转成声音，扩充训练集。合成声音的音色和噪声分布与真实录音不同，需与真实来源分开记录、分开报告。

与流水线相比，端到端模型避开了文字接口上的硬决策，延迟也只来自一个模型；代价是中间结果不可读，出错时难以判断是没听清还是没理解。若要求输出证据区间，必须说明监督对应声学帧区间还是转写字符区间，没有转写时不能称为文字跨度。比较两条路线时应固定声音、语义模式和训练资源，报告关键槽位与完整请求正确率。

\subsection{话语行为与情绪表达}
\label{sec:nlp-speech-pragmatics}

“可以”可能表示确认、许可或勉强接受。交流作用由话语行为刻画，其文本侧建模（HMM与CRF序列标注、分层对话编码）已在第\ref{sec:nlp-analysis-dialogue-acts}小节展开；情绪标签则按协议描述表达出的类别或唤醒度、效价量表，不等于直接测量心理状态。声音带来的新问题是：可用证据多了语调与节奏，文字又可能来自有误的识别。实验需要先固定输入范围（声音、文字、历史、说话者、回应对象各取哪些），标注协议允许不确定或多标签并报告标注者一致性，划分时按说话者和会话留出，以免模型记住某人的声音。类别不均衡时报告宏平均F1，量表任务另报告误差与相关性。

\subsubsection{历史条件下的话语行为分类}
\label{sec:nlp-speech-dialogue-act}

一个可复现的输入是“前轮内容、当前文字、当前说话者、回应对象、已确认状态”，用拼接历史的预训练编码器或第\ref{sec:nlp-analysis-dialogue-acts}小节的分层编码器（在线使用时改为单向）做分类，以交叉熵学习请求、告知、确认、否定、更正等标签；一句话兼有确认与补充告知时，改为多个二元标签并在开发集上定阈值。前轮是“周五下午两点方便吗”时，“可以”标为确认时间；前轮是“我可以修改人数吗”时，同样的回答表示许可。只看当前句的分类器面对完全相同的输入，无法区分两者，加入前轮与回应对象后依据才充分。

口语场景下历史本身也来自识别，前轮若被识别错误或被截断，分类器应输出不确定，系统据此发起澄清，例如“你是确认时间，还是允许修改人数”。评价把标签正确与对话后果分开：行为判为“确认”仍可能确认了错误的值，状态更新还需核对被确认的是哪一项。

\subsubsection{韵律特征、SVM与声文融合}
\label{sec:nlp-speech-prosody-classification}

同一句“可以”读得平缓或急促，节奏可能帮助解释表达。\term{韵律}指音高变化、能量、时长与停顿等跨越单个音素的组织。可复现的基线先从有效语音帧提取基频$F_0$、对数能量与倒谱特征，再统计均值、标准差、分位数、斜率、总时长与停顿比例，形成固定长度向量$\bm{v}$。无声段没有基频，应标为缺失并另记有声帧比例，不能按0\,Hz参与平均，否则停顿多的句子会被误算成音调低。

以二分类标签$z_i\in\{-1,1\}$为例，标准化参数只在训练集估计，支持向量机（support vector machine, SVM）最小化
\[
  \tfrac12\lVert\bm{w}\rVert_2^2+C\sum_i\max\bigl(0,\,1-z_i(\bm{w}^{\top}\bm{v}_i+b)\bigr),
\]
$C$由开发集选定\scite{cortes1995}。手工特征维数低、可解释，适合作为基线。更强的做法用式\eqref{eq:nlp-speech-slu-pooling}把自监督声学表示池化为$\bm{h}_{\mathrm{audio}}$，再与文本表示$\bm{h}_{\mathrm{text}}$融合。早融合把两者拼接为$[\bm{h}_{\mathrm{audio}};\bm{h}_{\mathrm{text}}]$后接分类头，可以学习声文交互；晚融合分别训练两个分类器，再以$\alpha\log p_{\mathrm{audio}}+(1-\alpha)\log p_{\mathrm{text}}$组合，某一路缺失时仍能工作；还可以让文字位置经交叉注意力读取声音序列。

教学对照固定“可以”的文字，只把时长从0.8\,s改为0.4\,s并改变音高走势，观察分类器输入怎样变化，而不预设短促就表示不耐烦。只有在留出的说话者与会话上比较纯文字、纯声音与融合模型，才能判断声音提供了多少增益；识别出某种情绪表达后，如何回应仍由对话目标决定。

\section{语音合成}
\label{sec:nlp-speech-tts}

\term{语音合成}（text-to-speech, TTS）把文字及规定的声音条件转换成可播放的波形，要求声音可懂、自然，并且内容正确。同一通知可以快读、慢读、强调新日期或在地点前停顿，这些都是允许的自由；改变日期、数量和否定则不允许。合成面对五项挑战。第一，读法不唯一：“10.5万元”“201”“重庆”在读出之前要先确定发音。第二，长度不匹配：几十个音素要对应几百个声学帧，模型必须决定每个音素持续多久。第三，一对多：同一句文字有许多种合理的语调和节奏，按均方误差回归会把它们平均成一种平淡模糊的声音。第四，相位缺失：梅尔频谱只保留各频带能量，从它恢复波形要补出相位与细节，而且每秒有数万个采样。第五，音色与风格控制：要用指定说话者的声音朗读，最好只给几秒参考录音就能做到。

方法谱系可以按“一个训练好的模型覆盖流水线的哪几段”来看，见图\ref{fig:nlp-speech-tts-lineage}。最早的神经合成系统是两阶段流水线：文本前端处理读法，声学模型把音素映射为梅尔频谱，声码器再把频谱变成波形；Tacotron 2用注意力隐式对齐，FastSpeech 2改为显式预测时长，以解决跳读、重复和串行推断。两阶段分别训练，中间频谱的误差会累积，VITS把声学模型与声码器合成一个条件变分自编码器，联合训练并在内部求对齐。回归式与对抗式目标仍难以表达一对多分布，扩散与流匹配模型从噪声逐步生成频谱，直接学习整个条件分布。最后，神经编解码器把声音变成离散词元，合成被改写为语言模型的续写任务，几秒参考录音就成了上下文提示。节末表\ref{tab:nlp-speech-tts-methods}比较这些方法的对齐方式、自回归性、并行度、零样本能力与内容错误风险。

\begin{figure*}[htbp]
  \centering
  % Four families of TTS systems placed on a common pipeline. Each box spans the
% stages that one trained model covers; arrows mark the interface between models.
\begin{tikzpicture}[
  x=1mm,y=1mm,figure picture,
  stage/.style={font=\figurefont\scriptsize\bfseries,text=FigureInk,align=center},
  fam/.style={font=\figurefont\scriptsize,text=FigureInk,anchor=east,align=right,text width=22mm},
  tm/.style={figure node,minimum height=9mm,inner ysep=2pt,font=\figurefont\scriptsize},
  front/.style={tm,draw=FigureInput,fill=FigureInputFill},
  acou/.style={tm,draw=FigureModel,fill=FigureModelFill},
  mid/.style={tm,draw=FigureStroke,fill=white},
  voc/.style={tm,draw=FigureLoss,fill=FigureLossFill},
  flow/.style={figure flow}
]
  % column headings
  \node[stage] at (14,8) {文本前端};
  \node[stage] at (55,8) {对齐与声学生成};
  \node[stage] at (93,8) {中间表示};
  \node[stage] at (124,8) {波形};
  \draw[figure link,draw=FigureGrid] (0,4) -- (140,4);

  % row 1: two-stage
  \node[fam] at (-2,-6) {两阶段流水线\\（\ref{sec:nlp-speech-acoustics}、\ref{sec:nlp-speech-waveform}）};
  \node[front,text width=24mm] (f1) at (14,-6) {规范化与G2P};
  \node[acou,text width=40mm] (a1) at (55,-6) {Tacotron 2：注意力隐式对齐\\FastSpeech 2：显式时长};
  \node[mid,text width=18mm] (i1) at (93,-6) {梅尔频谱};
  \node[voc,text width=26mm] (v1) at (124,-6) {WaveNet或\\HiFi-GAN};
  \draw[flow] (f1) -- (a1); \draw[flow] (a1) -- (i1); \draw[flow] (i1) -- (v1);

  % row 2: VITS
  \node[fam] at (-2,-20) {端到端VAE\\（\ref{sec:nlp-speech-vits}）};
  \node[front,text width=24mm] (f2) at (14,-20) {音素};
  \node[acou,text width=106mm] (a2) at (85,-20) {VITS：单调对齐搜索与随机时长；规范化流先验；隐变量经对抗训练的生成器直接得到波形};
  \draw[flow] (f2) -- (a2);

  % row 3: diffusion / flow matching
  \node[fam] at (-2,-34) {扩散与流匹配\\（\ref{sec:nlp-speech-diffusion-tts}）};
  \node[front,text width=24mm] (f3) at (14,-34) {音素或字符};
  \node[acou,text width=40mm] (a3) at (55,-34) {Grad-TTS：分数去噪\\F5-TTS：流匹配填充};
  \node[mid,text width=18mm] (i3) at (93,-34) {梅尔频谱};
  \node[voc,text width=26mm] (v3) at (124,-34) {声码器};
  \draw[flow] (f3) -- (a3); \draw[flow] (a3) -- (i3); \draw[flow] (i3) -- (v3);

  % row 4: codec language model
  \node[fam] at (-2,-48) {编解码词元\\语言模型（\ref{sec:nlp-speech-codec-lm-tts}）};
  \node[front,text width=24mm] (f4) at (14,-48) {音素$+$3秒提示};
  \node[acou,text width=40mm] (a4) at (55,-48) {VALL-E：首层自回归\\其余层非自回归};
  \node[mid,text width=18mm] (i4) at (93,-48) {离散声学\\词元};
  \node[voc,text width=26mm] (v4) at (124,-48) {编解码器\\解码器};
  \draw[flow] (f4) -- (a4); \draw[flow] (a4) -- (i4); \draw[flow] (i4) -- (v4);
\end{tikzpicture}

  \caption{四类语音合成系统在同一流水线上的覆盖范围。每个色块是一个单独训练的模型，箭头是模型之间的接口。两阶段流水线和扩散、流匹配模型以梅尔频谱为接口，另配声码器；VITS用一个模型从音素直接生成波形；编解码词元语言模型以离散声学词元为接口，由编解码器的解码器恢复波形。文本前端在端到端系统中可以被字符输入隐式吸收，但读法职责并未消失。}
  \label{fig:nlp-speech-tts-lineage}
\end{figure*}
\FloatBarrier

\subsection{文本前端}
\label{sec:nlp-speech-normalization}

书面文字到发音之间还隔着两步。\term{文本规范化}（text normalization）把日期、金额、编号等非标准写法展开成可以朗读的文字，例如“10.5万元”展开为“十点五万元”；字形到发音映射再决定每个字词的读音，例如“重庆”的“重”读\code{chong2}。直接读取字符的端到端合成没有显式前端，这两项职责由声学模型隐式完成，出错时更难定位，也更难由人工修正。下面依次讨论规则方法、神经规范化与G2P，比较维度是覆盖率、可控性与错误类型。

\subsubsection{规则词典与WFST规范化}
\label{sec:nlp-speech-wfst-normalization}

工业系统的规范化长期采用“先分类、后展开”的两步规则方法。Google的Kestrel系统是一个代表\scite{nlp-speech-ebden2015kestrel}：分类器把输入切成带\term{符号类别}（semiotic class）的单元，如日期、金额、基数、序数、编号、电话、缩写和普通词；展开器再按类别调用相应文法生成读法。两步都写成第\ref{sec:nlp-speech-hybrid}小节识别解码图所用的加权有限状态转导器（WFST）：弧上带输入符号、输出符号和代价，路径把输入串转换为输出串，代价为弧代价之和。

用WFST表达规范化的好处在于可以复合。设$T_{\mathrm{cls}}$把字符串映射为带类别标记的串，$T_{\mathrm{verb}}$把带标记的串映射为读法，两者的复合$T=T_{\mathrm{cls}}\circ T_{\mathrm{verb}}$仍是一个WFST，直接把书面文字映射到读法。对输入串$x$，规范化结果是所有输入为$x$的路径中代价最小者的输出：
\begin{equation}
  \hat y=\operatorname{out}\Bigl(\argmin_{\pi:\,\operatorname{in}(\pi)=x}\;\sum_{e\in\pi}w(e)\Bigr),
  \label{eq:nlp-speech-wfst}
\end{equation}
实现上先构造只接受$x$的线性自动机，与$T$组合后在得到的有向无环图上求最短路径。规则和词典来自显式编写的语言约定，代价可以人工设定，也可以用已标注的书面—读出对估计。

三个例子说明类别怎样决定读法。“10月5日”被日期文法消费为月份10、日期5，输出“十月五日”。“10.5万元”被金额文法输出“十点五万元”，其规范值仍为105000元。“201”有两条候选路径：编号读法“二零一”与基数读法“二百零一”。若前文是“海棠楼”或“房间”，上下文弧给编号路径较低代价，例如编号0.2、基数0.5，最短路径选中“二零一”；若前文是“共有”，代价反过来。同一数字串不能只凭局部形状决定类别。

规则方法可控、可审计：每个读法都能追溯到一条规则，修改一条规则只影响相应类别。缺点是词典缺项与规则冲突需要人工不断补充；规则只看局部窗口，难以利用整句上下文，例如“2/3”在“完成了2/3”中读作分数，在“2/3开会”中读作日期。

\subsubsection{句级神经文本规范化}
\label{sec:nlp-speech-neural-normalization}

规则需要显式枚举写法，整句上下文又会改变某个单元的读法。神经规范化把这件事交给序列到序列模型学习。Zhang等的系统以每个待规范化单元为中心取一个窗口，编码窗口内的字符，由解码器生成该单元的读法；大部分单元原样保留，模型只需输出“不变”标记\scite{nlp-speech-zhang2019textnorm}。Mansfield等改为以整句为输入，使用子词单位，一次生成整句读法，以读取窗口以外的上下文\scite{nlp-speech-mansfield2019}。

训练对$(x,r)$可以来自人工标注，也可以由规则系统生成后筛查；后一种数据会继承规则错误。模型按第\ref{sec:nlp-sft-masked-loss}小节的教师强制最小化$-\sum_j\log p(r_j\mid x,r_{<j})$，推理用贪心或束搜索生成读法。神经模型的典型失败是不可恢复的错误：输出流畅，却把“105000”读成“一万零五百”，听者无从察觉。Zhang等的对策是\term{覆盖文法}（covering grammar）：为每个类别写一个宽松的WFST，只描述“这个数字串可能有哪些读法”，不负责在其中选择；解码时把模型输出限制在覆盖文法接受的集合内，这与第\ref{sec:nlp-inference-grammar}小节的文法约束解码是同一机制。神经模型负责借上下文选择，规则负责保证不越界。

对“预算10.5万元，10月5日确认”，合格的读法应同时保持金额105000元与日期10月5日。可以把输入和输出中的类型、数值与单位分别解析成结构记录，要求两者逐项相等，解析失败的候选不自动接受。这一检查能发现漏读“万”，但发现不了读音错误。正向规范化与第\ref{sec:nlp-speech-itn}小节的反向文本规范化方向相反：“二零一”可能来自“201”也可能来自“二〇一”，两者都不保证可逆。

\subsubsection{字形到发音：G2P}
\label{sec:nlp-speech-g2p}

得到“重庆银行”后，系统还要决定“重”和“行”怎么读。\term{字形到发音}（grapheme-to-phoneme, G2P）把文字映射成发音单位序列。中文先要约定发音单位：可以是带调拼音音节（\code{chong2}），也可以拆为声母、韵母加声调（\code{ch}、\code{ong}、2）；轻声与变调是在前端标注还是交给声学模型实现，也需事先规定。

中文G2P的主要困难在多音字，绝大多数字只有一个读音，查词典即可。对多音字，模型在词典给出的候选集合内做分类。设第$i$个字$x_i$的候选读音集合为$\mathcal{R}(x_i)$，预训练字符编码器给出上下文表示$\bm{h}_i$，则
\begin{equation}
  p(r\mid x,i)=\frac{\exp(\bm{w}_r^{\top}\bm{h}_i)}{\sum_{r'\in\mathcal{R}(x_i)}\exp(\bm{w}_{r'}^{\top}\bm{h}_i)},
  \qquad r\in\mathcal{R}(x_i),
  \label{eq:nlp-speech-polyphone}
\end{equation}
分母只在候选集合上求和，模型不会给“行”输出\code{chong2}这类不可能的读音。g2pM用双向LSTM实现这一分类，并给出了一个人工标注的普通话多音字基准\scite{nlp-speech-park2020g2pm}。“重庆”为\code{chong2 qing4}，“银行”为\code{yin2 hang2}；把每个字最常见的读音当成唯一读音，就会读出“zhong4庆”“银xing2”。

词典之外的词，尤其是外文名和新词，需要生成而非分类。Yao与Zweig把G2P看作字母序列到音素序列的翻译，用编码器—解码器模型以交叉熵训练\scite{nlp-speech-yao2015g2p}，推理在音素词表上束搜索，再核对音节数与语种标识。拿到本调读音后，再按规则处理变调，例如两个上声相连时前一个读作阳平：“展览”本调为\code{zhan3 lan3}，实际读作\code{zhan2 lan3}。

通知里的“海棠楼”核为\code{hai3 tang2 lou2}，“201”先经规范化得到“二零一”再转为发音单位；\code{ML-0418}若要朗读，必须事先规定字母和数字的逐位读法。G2P只决定读哪个音，停顿、句重音和语速由后面的声学模型决定。

\subsection{两阶段声学模型}
\label{sec:nlp-speech-acoustics}

两阶段流水线中的声学模型把发音序列$x=(x_1,\ldots,x_N)$映射为梅尔频谱$\bm{M}\in\mathbb{R}^{T\times80}$，每行是一帧80维的对数梅尔能量（第\ref{sec:nlp-speech-features}小节）；声码器再从$\bm{M}$生成波形。声学模型的核心困难是长度差：$N$是几十，$T$是几百，而且每个音素对应多少帧事先未知。下面两个模型给出两种回答。Tacotron 2逐帧自回归生成，用注意力在生成过程中隐式决定当前帧读哪个字符；FastSpeech 2先显式预测每个音素的帧数，展开后并行生成全部帧。频谱提取参数（采样率、帧移、梅尔范围）须与声码器一致，不能只因两边都是80个频带就直接拼接。

\subsubsection{注意力声文对齐：Tacotron 2}
\label{sec:nlp-speech-tacotron}

Tacotron 2把合成看作序列到序列生成：编码器读字符，解码器每步生成一帧梅尔频谱，注意力决定这一帧对应哪些字符\scite{nlp-intro-shen2018}。图\ref{fig:nlp-speech-tacotron2}给出完整结构，下面按数据流逐层说明。

\begin{figure*}[tp]
  \centering
  % Tacotron 2 architecture: encoder (left), location-sensitive attention (middle),
% autoregressive decoder with post-net (right). Layer sizes follow Shen et al. (2018).
\begin{tikzpicture}[
  x=1mm,y=1mm,figure picture,
  enc/.style={figure node,draw=FigureInput,fill=FigureInputFill,text width=33mm,minimum height=9mm},
  dec/.style={figure node,draw=FigureModel,fill=FigureModelFill,text width=33mm,minimum height=9mm},
  att/.style={figure node,draw=FigureLoss,fill=FigureLossFill,text width=27mm,minimum height=11mm},
  io/.style={figure node,draw=FigureStroke,fill=white,text width=24mm,minimum height=8mm},
  flow/.style={figure flow},
  tstep/.style={figure flow,dash pattern=on 2pt off 1.5pt,draw=FigureMuted},
  lab/.style={figure edge label,font=\figurefont\scriptsize}
]
  % encoder
  \node[io] (chars) at (0,0) {字符序列，长度$N$};
  \node[enc] (emb) at (0,14) {字符嵌入\\$N\times512$};
  \node[enc] (conv) at (0,28) {3层卷积\\核宽5，512通道};
  \node[enc] (lstm) at (0,42) {双向LSTM\\每向256单元};
  \node[io,text width=30mm] (H) at (0,56) {编码$\bm H\in\mathbb R^{N\times512}$};
  \foreach \a/\b in {chars/emb,emb/conv,conv/lstm,lstm/H}{\draw[flow] (\a) -- (\b);}

  % attention
  \node[att] (att) at (50,42) {位置敏感注意力\\权重$\bm\alpha_i$};
  \draw[flow] (H.east) -- (att.north west);
  \draw[tstep] ([xshift=4mm]att.north) .. controls +(1,9) and +(9,7) .. (att.north east);
  \node[anchor=south,font=\figurefont\scriptsize,text=FigureMuted,align=center] at (57,55.5) {累积权重$\sum_{k<i}\bm\alpha_k$\\经卷积成位置特征};

  % decoder
  \node[io,text width=30mm] (prev) at (100,0) {前一帧$\bm m_{i-1}\in\mathbb R^{80}$};
  \node[dec] (pre) at (100,14) {预网络\\2层全连接256};
  \node[dec] (dlstm) at (100,28) {2层单向LSTM\\每层1024单元};
  \node[dec] (proj) at (100,42) {两个线性投影\\输入$[\bm s_i;\bm c_i]$};
  \draw[flow] (prev) -- (pre);
  \draw[flow] (pre) -- (dlstm);
  \draw[flow] (dlstm) -- (proj);
  \draw[flow] ([yshift=1.5mm]dlstm.west) -- (60,29.5) -- node[lab,pos=.45] {$\bm s_i$} (60,36.5);
  \draw[flow] (att.east) -- node[lab] {$\bm c_i$} (proj.west);
  \draw[tstep] (44,36.5) |- node[lab,pos=.75] {$\bm c_{i-1}$} ([yshift=-2mm]dlstm.west);

  % outputs
  \node[io,text width=20mm] (stop) at (80,62) {停止概率$\hat p_i$};
  \node[io,text width=20mm] (mel) at (114,62) {梅尔帧$\hat{\bm m}_i$};
  \draw[flow] (proj.north) ++(-6,0) |- ++(0,4) -| (stop.south);
  \draw[flow] (proj.north) ++(6,0) |- ++(0,4) -| (mel.south);
  \node[dec,text width=28mm] (post) at (114,74) {后处理网络\\5层卷积，预测残差};
  \node[circle,draw=FigureStroke,line width=.75pt,inner sep=1pt] (plus) at (114,84) {$+$};
  \node[io,text width=30mm] (melp) at (114,93) {修正频谱$\widehat{\bm M}^{+}$};
  \draw[flow] (mel) -- (post);
  \draw[flow] (post) -- (plus);
  \draw[flow] (plus) -- (melp);
  \draw[flow] (mel.west) -- (95,62) |- (plus.west);
  \draw[tstep] (mel.east) -- (134,62) |- node[pos=.25,fill=white,font=\figurefont\scriptsize,text=FigureMuted,align=center] {推断时\\反馈；\\训练时\\用真实帧} (prev.east);
\end{tikzpicture}

  \caption{Tacotron 2的结构。左列为编码器：长度为$N$的字符序列经嵌入、3层卷积和双向LSTM得到编码$\bm H$。中间为位置敏感注意力，它在打分时读取此前各步注意力权重的累积和。右列为解码器：前一帧经预网络，与上一步上下文$\bm c_{i-1}$一起送入2层LSTM，得到状态$\bm s_i$；$[\bm s_i;\bm c_i]$经两个线性投影分别输出梅尔帧与停止概率，后处理网络预测残差得到修正频谱。虚线表示跨时间步的连接。层数与宽度取自原论文。}
  \label{fig:nlp-speech-tacotron2}
\end{figure*}

编码器把字符序列变成上下文相关的表示。每个字符先查512维嵌入，得到$N\times512$的矩阵；接着经过3层一维卷积，每层512个核宽为5的滤波器，后接批归一化与ReLU，卷积让每个位置看到前后各若干字符，起到类似$n$元特征的作用；最后一层双向LSTM每个方向256单元，拼接后得到$\bm{H}=(\bm{h}_1,\ldots,\bm{h}_N)\in\mathbb{R}^{N\times512}$。

解码器第$i$步要预测第$i$帧$\bm{m}_i\in\mathbb{R}^{80}$。前一帧$\bm{m}_{i-1}$先经\term{预网络}（pre-net），即两层256单元的全连接层加ReLU，并以0.5的概率丢弃；丢弃在推断时也保持开启，原论文用它给输出引入随机性。预网络输出与上一步的上下文向量$\bm{c}_{i-1}$拼接，送入两层各1024单元的单向LSTM，得到解码器状态$\bm{s}_i$。

注意力用$\bm{s}_i$在$\bm{H}$上取上下文。普通加性注意力只比较$\bm{s}_i$与$\bm{h}_j$的内容，长句中容易跳到内容相似的远处字符。\term{位置敏感注意力}（location-sensitive attention）把“此前已经读到哪里”作为额外特征\scite{nlp-speech-chorowski2015}。设此前各步注意力权重的累积和为$\bm{\beta}_{i}=\sum_{k<i}\bm{\alpha}_k\in\mathbb{R}^{N}$，用$32$个长度为31的一维卷积核$\bm{F}$与它卷积，得到每个位置$j$的位置特征$\bm{f}_{i,j}\in\mathbb{R}^{32}$，则
\begin{equation}
\begin{aligned}
  e_{i,j}&=\bm{v}^{\top}\tanh\bigl(\bm{W}\bm{s}_i+\bm{V}\bm{h}_j+\bm{U}\bm{f}_{i,j}+\bm{b}\bigr),\qquad
  \bm{f}_{i}=\bm{F}*\bm{\beta}_{i},\\
  \alpha_{i,j}&=\frac{\exp(e_{i,j})}{\sum_{j'=1}^{N}\exp(e_{i,j'})},\qquad
  \bm{c}_i=\sum_{j=1}^{N}\alpha_{i,j}\bm{h}_j .
\end{aligned}
\label{eq:nlp-speech-location-attention}
\end{equation}
$\bm{W}$、$\bm{V}$、$\bm{U}$把三种输入投影到128维，$\bm{v}\in\mathbb{R}^{128}$给出标量分数。累积权重记录了每个字符已被读过多少：已读完的字符累积值接近1，尚未读到的接近0，卷积核可以学会“偏好累积值刚从1降到0的位置”，即紧接着已读部分的下一个字符。这样，对齐倾向于单调向前移动。

输出端把$[\bm{s}_i;\bm{c}_i]$经一个线性投影得到梅尔帧$\hat{\bm{m}}_i$，经另一个线性投影和sigmoid得到\term{停止标记}（stop token）概率$\hat p_i$，表示合成是否应在这一帧结束。全部帧拼成$\widehat{\bm{M}}$后，\term{后处理网络}（post-net）用5层一维卷积（每层512个核宽为5的滤波器，批归一化，除最后一层外接tanh）预测残差，得到$\widehat{\bm{M}}^{+}=\widehat{\bm{M}}+\operatorname{PostNet}(\widehat{\bm{M}})$。后处理网络能看到整段频谱，可以修正逐帧生成时只看过去造成的局部误差。训练目标为
\begin{equation}
  L_{\mathrm{Taco}}=\lVert\widehat{\bm{M}}-\bm{M}\rVert_2^2
  +\lVert\widehat{\bm{M}}^{+}-\bm{M}\rVert_2^2
  +\sum_{i=1}^{T}\operatorname{BCE}(\hat p_i,\,\mathbb{1}[i=T]),
  \label{eq:nlp-speech-tacotron-loss}
\end{equation}
两项频谱误差分别监督后处理前后的输出，第三项是停止标记的二元交叉熵，只有最后一帧的目标为1。训练时解码器输入真实的前一帧（教师强制），推断时输入自己生成的帧；推断在$\hat p_i$首次超过0.5时停止，实际系统还应设最大帧数，并把截断标记为异常而非正常完成。原系统的声码器是以预测频谱为条件的修改版WaveNet，见第\ref{sec:nlp-speech-wavenet}小节。

Tacotron 2的优点是结构统一，对齐不需要外部标注，音质在当时接近真人录音。它的缺点都来自隐式对齐与自回归。注意力可能卡在某个字符附近导致重复，也可能提前跳到句尾导致漏读，长句与训练集中少见的文本尤其容易出错。逐帧生成使推断时间与帧数成正比，无法并行。训练用真实帧、推断用自己生成的帧，误差会沿时间累积。对“四月十八日”，可以观察注意力是否依次经过各字，但注意力图只能提示问题，是否漏读仍需听觉核验。

\subsubsection{时长与韵律条件：FastSpeech 2}
\label{sec:nlp-speech-fastspeech}

Tacotron 2的跳读、重复与串行推断都源于“长度在生成中隐式形成”。若先给每个音素分配帧数，就能在长度确定后一次性并行生成全部帧，对齐也不会跳读。FastSpeech用自回归教师模型的注意力提取时长，并以教师合成的频谱作为蒸馏目标\scite{nlp-speech-ren2019fastspeech}；FastSpeech 2去掉教师，改用强制对齐器从真实录音提取时长，直接以真实频谱为目标，并把音高与能量也作为显式条件输入，以缓解一对多问题\scite{nlp-speech-ren2021}。

模型由三部分组成：音素编码器、方差适配器和梅尔频谱解码器。编码器与解码器各由4个\term{前馈Transformer块}（feed-forward Transformer block, FFT块）堆叠而成。FFT块与第\ref{sec:nlp-llm-decoder-layer}小节的Transformer层相同，只是把逐位置前馈网络换成两层一维卷积，以便利用相邻帧或相邻音素的局部关系。设输入$\bm{X}\in\mathbb{R}^{L\times256}$，
\begin{equation}
\begin{aligned}
  \bm{X}'&=\operatorname{LN}\bigl(\bm{X}+\operatorname{MHSA}(\bm{X})\bigr),\\
  \operatorname{FFT}(\bm{X})&=\operatorname{LN}\bigl(\bm{X}'+\operatorname{Conv}_{2}(\operatorname{ReLU}(\operatorname{Conv}_{1}(\bm{X}')))\bigr),
\end{aligned}
\label{eq:nlp-speech-fft-block}
\end{equation}
$\operatorname{MHSA}$为多头自注意力，$\operatorname{LN}$为层归一化；$\operatorname{Conv}_1$把256通道升到1024通道，$\operatorname{Conv}_2$再降回256通道。原文配置中自注意力为2头，两层卷积的核宽分别为9和1。编码器输入是音素嵌入加正弦位置编码，长度$L=N$；解码器输入是方差适配器展开后的序列，长度$L=T$，最后经线性层输出$T\times80$的频谱。

\term{方差适配器}（variance adaptor）依次加入时长、音高和能量三类信息。三个预测器结构相同：两层一维卷积（核宽3、256通道、ReLU），每层后接层归一化与丢弃，最后经线性层为每个位置输出一个标量。时长预测器作用在音素级编码$\bm{h}_1,\ldots,\bm{h}_N$上，预测对数时长$\hat r_i\approx\log(1+d_i)$，其中$d_i$是第$i$个音素的整数帧数；训练目标$d_i$由Montreal强制对齐器把音素对到录音后换算得到\cite{nlp-speech-mcauliffe2017}，强制对齐的原理见第\ref{sec:nlp-speech-forced-alignment}小节。\term{长度调节器}按$d_i$重复第$i$个向量，使隐藏长度与频谱一致：
\begin{equation}
 \widetilde{\bm{h}}_t=\bm{h}_i
 \quad\text{当}\quad
 \sum_{j<i}d_j<t\leq\sum_{j\leq i}d_j,
 \qquad T=\sum_{i=1}^{N}d_i.
 \label{eq:nlp-speech-length-regulator}
\end{equation}
教学音素序列/p, a, t, a/的时长为$(2,3,1,2)$帧，展开为
\[
  (\bm{h}_1,\bm{h}_1,\bm{h}_2,\bm{h}_2,\bm{h}_2,\bm{h}_3,\bm{h}_4,\bm{h}_4),
\]
共8帧；帧移10\,ms时对应80\,ms。推断时用$\hat d_i=\max(1,\operatorname{round}(\exp\hat r_i-1))$恢复整数，下界1保证非静音音素不被删掉，这是实现约定。语速控制把时长除以速度因子再取整：对$(2,3,1,2)$取因子1.5，得$(1.33,2,0.67,1.33)$，四舍五入并取下界1后为$(1,2,1,1)$，共5帧；$8/1.5$不是整数，实际加速比也不恰好是1.5。

长度调节之后，音高与能量预测器在帧级序列$\widetilde{\bm{H}}\in\mathbb{R}^{T\times256}$上工作。音高从训练录音估计基频得到，无声帧用插值补齐；能量取每帧STFT幅度向量的$L_2$范数。原方法还把连续音高曲线作连续小波变换，预测其多尺度表示后再逆变换恢复，以便刻画不同时间尺度的起伏。为把标量条件加回隐藏序列，FastSpeech 2把它们分桶量化后查嵌入表：取边界$b_0<b_1<\cdots<b_{K}$（音高在对数刻度上均匀取，能量在线性刻度上均匀取，$K=256$），第$t$帧的值$p_t$落入区间$[b_{k-1},b_k)$时取第$k$个嵌入$\bm{E}^{\mathrm{pitch}}_k\in\mathbb{R}^{256}$，
\begin{equation}
  \widetilde{\bm{h}}_t\leftarrow\widetilde{\bm{h}}_t
  +\bm{E}^{\mathrm{pitch}}_{\operatorname{bin}(p_t)}
  +\bm{E}^{\mathrm{energy}}_{\operatorname{bin}(\varepsilon_t)} .
  \label{eq:nlp-speech-variance-embed}
\end{equation}
以教学化的4个能量桶为例，边界取0、25、50、75、100，能量值12、48、51、80分别落入第1、2、3、4桶；48与51数值接近却落入不同桶，这是量化的代价，256个桶时相邻桶的嵌入在训练中会变得相近。训练时方差适配器使用从录音提取的真实时长、音高与能量（与教师强制同理），推断时全部改用预测值。

总目标是频谱的$L_1$误差加三项方差预测的均方误差：
\begin{equation}
  L_{\mathrm{FS2}}=\lVert\widehat{\bm{M}}-\bm{M}\rVert_1
  +\sum_{i}\bigl(\hat r_i-\log(1+d_i)\bigr)^2
  +\sum_{t}(\hat p_t-p_t)^2+\sum_{t}(\hat\varepsilon_t-\varepsilon_t)^2 .
  \label{eq:nlp-speech-fs2-loss}
\end{equation}
显式的音高与能量条件解决了一部分一对多问题：给定这些值后，频谱的不确定性大大降低，解码器不必再把不同语调平均成一种。但三个预测器本身仍以均方误差回归，预测出的音高曲线依然趋于平均，这是后面扩散与流匹配模型要解决的问题。

与Tacotron 2相比，FastSpeech 2的对齐由时长决定，不会跳读或重复，全部帧并行生成，语速、音高和能量可以直接调节，例如对通知强调“十八日”时延长相应音素并提高音高，而不改变“校外来宾”和“前”等内容。代价是依赖外部强制对齐器提供时长，对齐器的错误会直接成为训练目标；时长取整和分桶量化也引入了离散误差。

\subsection{神经声码器}
\label{sec:nlp-speech-waveform}

梅尔频谱只保留各频带的能量，丢掉了相位和频带内的细节；\term{声码器}（vocoder）要在给定频谱的条件下补出这些信息，生成波形。一个可复现的接口至少规定频谱提取参数、采样率、每帧对应的采样数（帧移）以及补边与裁剪方式；波形时长由采样数除以采样率得到，不能把帧数直接当成秒数。设采样率24\,kHz、帧移12.5\,ms，则每帧对应300个采样，80帧对应24000个采样，即1秒。声码器的谱系从逐采样自回归的WaveNet开始；为了并行生成，后续出现了知识蒸馏的并行WaveNet、基于规范化流的WaveGlow，以及基于生成对抗网络的MelGAN与HiFi-GAN。下面写出两端的代表：WaveNet与HiFi-GAN。

\subsubsection{条件逐采样生成：WaveNet}
\label{sec:nlp-speech-wavenet}

WaveNet把波形看成一维序列，按自回归分解其概率\scite{nlp-speech-oord2016wavenet}。设目标波形为$x_1,\ldots,x_N$，梅尔条件为$\bm{M}$，则
\begin{equation}
 p(x_{1:N}\mid\bm{M})
   =\prod_{n=1}^{N}p_\theta(x_n\mid x_{<n},\bm{M}),
 \qquad
 L_{\mathrm{wave}}=-\sum_{n=1}^{N}
       \log p_\theta(x_n\mid x_{<n},\bm{M}).
 \label{eq:nlp-speech-wavenet}
\end{equation}
困难在于$x_n$依赖很长的历史：一个音节持续数千个采样。WaveNet用\term{空洞因果卷积}（dilated causal convolution）扩大感受野。因果指输出只依赖当前及过去的采样；空洞率为$r$、核宽为$K$的卷积把相隔$r$的采样组合起来，$y_n=\sum_{k=0}^{K-1}w_k\,x_{n-kr}$。逐层把空洞率加倍，感受野呈指数增长。叠$L$层、第$\ell$层空洞率为$r_\ell$时，感受野为
\begin{equation}
  R=1+(K-1)\sum_{\ell=1}^{L}r_\ell .
  \label{eq:nlp-speech-receptive-field}
\end{equation}
取$K=2$，空洞率$1,2,4,\ldots,512$共10层为一组，一组的感受野为$1+1023=1024$个采样；重复3组得$1+3\times1023=3070$个采样，24\,kHz下约128毫秒。若不用空洞而要达到同样感受野，需要三千多层普通因果卷积。

每层的非线性是\term{门控激活单元}（gated activation unit）。设第$\ell$层输入为$\bm{z}$，上采样到采样率的梅尔条件为$\bm{y}$，则
\begin{equation}
  \bm{o}=\tanh(\bm{W}_{f}*\bm{z}+\bm{V}_{f}\bm{y})\odot\sigma(\bm{W}_{g}*\bm{z}+\bm{V}_{g}\bm{y}),
  \label{eq:nlp-speech-wavenet-gate}
\end{equation}
$*$是空洞因果卷积，$\bm{V}_f\bm{y}$、$\bm{V}_g\bm{y}$是条件的$1\times1$卷积，$\odot$为逐元素乘。tanh支路产生候选信号，sigmoid支路决定放行多少，与LSTM的门同理。$\bm{o}$经$1\times1$卷积后一路加回输入形成残差连接，一路作为跳跃连接汇总到输出端；所有层的跳跃输出相加，经两次ReLU与$1\times1$卷积后得到输出分布。梅尔条件$\bm{y}$由转置卷积把每帧上采样300倍，与采样对齐。

输出分布有两种。原WaveNet把16位采样先用\term[mu律压扩]{μ律压扩}（μ-law companding）压缩，再均匀量化为256类，用softmax建模：
\begin{equation}
  f(x)=\operatorname{sign}(x)\frac{\ln(1+\mu\lvert x\rvert)}{\ln(1+\mu)},\qquad
  x\in[-1,1],\ \mu=255 .
  \label{eq:nlp-speech-mulaw}
\end{equation}
它给小振幅更细的分辨率：$x=0.01$映射为$\ln3.55/\ln256\approx0.228$，$x=0.5$映射为$\ln128.5/\ln256\approx0.876$；在正半轴上，原始区间$[0,0.01]$只占1\%，压扩后却占约23\%的量化级。人耳对安静段的量化噪声更敏感，压扩把量化级更多地分给它们。Tacotron 2中的修改版改为输出10个分量的混合逻辑分布参数，直接对连续采样值计算似然，并以Tacotron 2在教师强制下预测的频谱为条件训练，使训练条件接近推断时的上游输出。

WaveNet的优点是音质高、概率模型明确；缺点是推断必须逐采样串行，生成1秒24\,kHz声音需要24000步，每步都要过一遍全部层，缓存可以避免重复计算，却不能把串行变成并行。另一方面，声码器只负责把频谱变成声音：上游频谱若对应“十七日”，再自然的波形也不会变成“十八日”。

\subsubsection{并行波形生成：HiFi-GAN}
\label{sec:nlp-speech-hifigan}

要并行生成波形，可以用一个卷积网络直接把整段梅尔频谱上采样到波形长度。难点在训练目标：只用频谱或波形的$L_1$误差训练，会得到相位模糊、带金属声的结果，因为满足同一频谱的波形有许多个，回归目标会把它们平均。第\ref{sec:nlp-speech-soundstream}小节的编解码器已经用对抗训练解决同一问题：判别器学习区分真实与生成的样本，生成器学习让判别器无法区分，于是生成样本被推向真实分布的某一个模式而非平均。HiFi-GAN把这一思路落实为一个生成器和两组针对语音周期结构设计的判别器\scite{nlp-speech-kong2020}。

生成器输入$\bm{M}\in\mathbb{R}^{T\times80}$，先经一层核宽7的卷积升到512通道，再经4次转置卷积上采样。以22.05\,kHz、帧移256个采样的配置为例，四次上采样倍率为$(8,8,2,2)$，乘积为256，通道依次减半为256、128、64、32，长度从$T$变为$256T$。每次上采样之后接一个\term{多感受野融合}（multi-receptive field fusion, MRF）模块：它并行使用核宽为3、7、11的三个残差块，每个残差块内部用空洞率1、3、5的空洞卷积，输出取三者的平均。不同核宽与空洞率覆盖不同长度的模式，MRF让每一层同时看到短周期与长周期结构。最后一层卷积接tanh，输出$[-1,1]$内的波形。

\term{多周期判别器}（multi-period discriminator, MPD）针对语音的周期性。浊音由基频决定的周期重复，同一相位上的采样之间有规律。MPD把长度为$N$的波形按周期$p$重排为$\lceil N/p\rceil\times p$的二维矩阵，第$a$行第$b$列为$x_{(a-1)p+b}$，再用沿第一维的二维卷积（核宽$(5,1)$）处理，每一列只包含间隔恰为$p$的采样。例如$x_1,\ldots,x_{12}$取$p=3$，三列分别为$(x_1,x_4,x_7,x_{10})$、$(x_2,x_5,x_8,x_{11})$、$(x_3,x_6,x_9,x_{12})$。HiFi-GAN取$p\in\{2,3,5,7,11\}$五个子判别器，用互质的周期尽量避免重叠。\term{多尺度判别器}（multi-scale discriminator, MSD）沿用MelGAN的设计\cite{nlp-speech-kumar2019melgan}，在原波形以及2倍、4倍平均池化后的波形上各用一个一维卷积判别器，检查连续的局部模式与较长的时间结构；MPD只看等间隔采样，MSD补上了相邻采样之间的关系。

训练采用最小二乘形式的对抗损失。令$\hat x=G(\bm{M})$，$D_k$为第$k$个子判别器（MPD与MSD共8个），$\phi$为固定的梅尔提取器，$D_k^{(l)}$为$D_k$第$l$层的特征、共$N_l$个元素：
\begin{equation}
\begin{aligned}
 L_D&=\sum_k\mathbb{E}\bigl[(D_k(x)-1)^2+D_k(\hat x)^2\bigr],\\
 L_G&=\sum_k\Bigl[\mathbb{E}(D_k(\hat x)-1)^2+\lambda_{\mathrm{fm}}L_{\mathrm{fm}}(G;D_k)\Bigr]\\
 &\quad+\lambda_{\mathrm{mel}}\,\mathbb{E}\lVert\phi(x)-\phi(\hat x)\rVert_1,\\
 L_{\mathrm{fm}}(G;D_k)&=\mathbb{E}\sum_{l}\frac{1}{N_l}\bigl\lVert D_k^{(l)}(x)-D_k^{(l)}(\hat x)\bigr\rVert_1,
\end{aligned}
\label{eq:nlp-speech-hifigan}
\end{equation}
系数取$\lambda_{\mathrm{fm}}=2$，$\lambda_{\mathrm{mel}}=45$。判别器希望真实样本输出1、生成样本输出0；生成器希望生成样本输出1。特征匹配损失与EnCodec中的同名项作用相同（第\ref{sec:nlp-speech-soundstream}小节），比较真实与生成波形在判别器各中间层的特征，提供比单个真假分数更稠密的梯度；梅尔损失约束生成波形的频谱与条件一致，保证内容不走样，并使训练初期稳定。两者交替更新，推断只保留生成器，一次前向即得到整段波形。

HiFi-GAN的优点是并行生成，在GPU上远快于实时，音质与自回归声码器相当。代价是训练依赖判别器设计与损失系数，对抗训练可能不稳定，生成波形偶尔出现周期性伪影；判别器认为逼真，也不说明每个数字都读对了。比较WaveNet与HiFi-GAN的速度时，应统一采样率、频谱定义、硬件、批量和计时边界，分别记录首段可播放时间与整段生成时间。

沿前端、声学模型与声码器贯通一份通知：“机器学习读书会改至4月18日14:00，地点为海棠楼201。校外来宾须在4月16日18:00前提交报名信息。”规范化保留活动时间与报名截止两个字段，产生“四月十八日下午两点”与“四月十六日十八点前”等读法；G2P核对“海棠楼”及编号；声学模型保留每个音节和断句；声码器按约定采样率恢复波形。最终声音按四项独立问题验收：是否听到4月18日14:00，是否听到海棠楼201，是否保留“校外来宾”这一对象，是否听到严格早于4月16日18:00的截止条件。回转识别可以自动比较这些字段，但识别器可能把发错的音纠正成常用词，因此只作筛查；真人盲听转写、配对偏好和平均意见分（mean opinion score, MOS）另行记录材料、评价者与量表，内容正确、可懂度和自然度分别报告。

\subsection{端到端合成：VITS}
\label{sec:nlp-speech-vits}

两阶段流水线有两个遗留问题。声学模型与声码器分别训练，声码器在训练时看到的是真实频谱，推断时却要处理声学模型的预测频谱，两者的误差累积；FastSpeech 2又依赖外部对齐器提供时长。VITS用一个模型同时解决这两个问题：它把从文字到波形的生成写成条件变分自编码器，以隐变量替代梅尔频谱作为内部接口，对齐由模型在训练中自己求得\scite{nlp-speech-kim2021vits}。

先回顾所需的两个工具。变分自编码器（VAE）为观测$\bm{x}$引入隐变量$\bm{z}$，用后验近似$q_\phi(\bm{z}\mid\bm{x})$与先验$p_\theta(\bm{z})$构造对数似然的下界：$\log p_\theta(\bm{x})\geq\mathbb{E}_{q_\phi}[\log p_\theta(\bm{x}\mid\bm{z})]-\operatorname{KL}(q_\phi(\bm{z}\mid\bm{x})\,\Vert\,p_\theta(\bm{z}))$，即\term{证据下界}，第一项要求能从$\bm{z}$重建$\bm{x}$，第二项要求后验贴近先验。规范化流用一个可逆映射$f$把简单分布变成复杂分布，换元后密度为$p(\bm{z})=\mathcal{N}(f(\bm{z}))\,\lvert\det\partial f(\bm{z})/\partial\bm{z}\rvert$，可精确计算。两者以及后两小节要用的扩散模型与流匹配，其一般推导可参考\BookChapterRef{chap:pgm-generative}{卷六的概率生成模型章}；本小节只用上述两个结论。

VITS的条件VAE由四部分组成。后验编码器读取目标语音的线性频谱$\bm{x}_{\mathrm{lin}}$（比梅尔频谱保留更多细节），经16个非因果WaveNet残差块输出每帧的均值与方差，给出$q_\phi(\bm{z}\mid\bm{x}_{\mathrm{lin}})$，$\bm{z}\in\mathbb{R}^{T\times192}$。解码器就是HiFi-GAN的生成器，把$\bm{z}$直接上采样为波形；为节省显存，训练时只解码随机截取的一小段$\bm{z}$。先验编码器由6层Transformer文本编码器和一个规范化流$f_\theta$组成：文本编码器为每个音素$i$输出均值$\bm{\mu}_i$与标准差$\bm{\sigma}_i$，对齐矩阵$A$把它们分配到各帧；流由4个仿射耦合层构成，把后验样本变换到文本给出的高斯分布上：
\begin{equation}
  p_\theta(\bm{z}\mid c,A)=\mathcal{N}\bigl(f_\theta(\bm{z});\,\bm{\mu}_{\theta}(c,A),\,\bm{\sigma}_{\theta}(c,A)\bigr)\,\Bigl\lvert\det\frac{\partial f_\theta(\bm{z})}{\partial\bm{z}}\Bigr\rvert,
  \label{eq:nlp-speech-vits-prior}
\end{equation}
$c$是音素序列。VITS把耦合层限定为只平移、不缩放的保体积变换，行列式恒为1，上式实际只剩高斯密度一项。只用对角高斯作先验时，同一音素各帧的分布完全相同，表达不出帧间变化；流让先验可以是复杂分布，缩小了与后验的差距。第四部分是随机时长预测器，下面单独说明。

对齐由\term{单调对齐搜索}（monotonic alignment search, MAS）求得，该方法最早用于Glow-TTS\cite{nlp-speech-kim2020glowtts}。对齐$A$把每帧$j$分配给一个音素$A(j)$，要求单调且不跳过：第1帧属于第1个音素，最后一帧属于最后一个音素，每帧要么留在前一帧的音素，要么前进一个。MAS在所有合法对齐中选使变换后隐变量在先验下似然最大的一个。记$\ell_{i,j}=\log\mathcal{N}(f_\theta(\bm{z}_j);\bm{\mu}_i,\bm{\sigma}_i)$，用动态规划
\begin{equation}
  Q_{i,j}=\max\bigl(Q_{i-1,j-1},\,Q_{i,j-1}\bigr)+\ell_{i,j},
  \qquad Q_{1,1}=\ell_{1,1},
  \label{eq:nlp-speech-mas}
\end{equation}
填表，再从$Q_{N,T}$回溯得到最优路径，复杂度$O(NT)$。它和第\ref{sec:nlp-classical-viterbi}小节的Viterbi算法同构，只是转移只允许“停留”和“前进一格”。教学示例取2个音素、3帧，$\ell_{1,\cdot}=(-1,-2,-4)$，$\ell_{2,\cdot}=(-5,-1,-1)$。第一行$Q_{1,\cdot}=(-1,-3,-7)$；第二行$Q_{2,1}=-\infty$（第1帧必须属于第1个音素），$Q_{2,2}=\max(-1,-\infty)-1=-2$，$Q_{2,3}=\max(Q_{1,2},Q_{2,2})-1=\max(-3,-2)-1=-3$。回溯得第1帧属于音素1、第2、3帧属于音素2，时长为$(1,2)$，总分$-3$；另一条合法路径（时长$(2,1)$）的总分为$-1-2-1=-4$。训练时每步先固定参数用MAS求对齐，再在该对齐下更新参数，与硬EM（Viterbi训练）的交替方式类似；对齐不参与梯度。

MAS给出的时长$d_i=\sum_j\mathbb{1}[A(j)=i]$还要用来训练一个推断时使用的时长预测器。FastSpeech 2的确定性预测器对同一文字总给出同一节奏。VITS的\term{随机时长预测器}（stochastic duration predictor）是一个基于流的条件生成模型，学习时长的分布$p(d\mid c)$。整数时长不能直接用连续密度建模，模型加入$[0,1)$内的均匀扰动$u$做变分去量化，并引入辅助变量$\nu$增加表达力，最大化
\begin{equation}
  \log p_\theta(d\mid c)\geq\mathbb{E}_{q_\phi(u,\nu\mid d,c)}\Bigl[\log\frac{p_\theta(d-u,\nu\mid c)}{q_\phi(u,\nu\mid d,c)}\Bigr],
  \label{eq:nlp-speech-vits-duration}
\end{equation}
其负值记为$L_{\mathrm{dur}}$。推断时从高斯噪声出发，经逆向流得到连续时长，再向上取整。同一文字采样不同噪声，就得到不同的节奏。

总目标把VAE的两项与对抗训练相加：
\begin{equation}
\begin{aligned}
  L_{\mathrm{VITS}}={}&\lambda_{\mathrm{mel}}\lVert\phi(x)-\phi(\hat x)\rVert_1
  +\underbrace{\mathbb{E}_{q_\phi}\bigl[\log q_\phi(\bm{z}\mid\bm{x}_{\mathrm{lin}})-\log p_\theta(\bm{z}\mid c,A)\bigr]}_{L_{\mathrm{kl}}}\\
  &+L_{\mathrm{dur}}+L_{\mathrm{adv}}(G)+L_{\mathrm{fm}}(G),
\end{aligned}
  \label{eq:nlp-speech-vits-loss}
\end{equation}
第一项是重建项，用梅尔$L_1$误差代替波形似然，相当于假设数据服从拉普拉斯分布；原文把各项直接相加，公开实现与HiFi-GAN一样取$\lambda_{\mathrm{mel}}=45$。第二项是单样本估计的KL项，连接后验与文本先验。对抗与特征匹配项沿用式\eqref{eq:nlp-speech-hifigan}；为节省训练开销，判别器保留HiFi-GAN的5个多周期子判别器，只从多尺度判别器中保留作用于原始波形的一个，相当于周期取$\{1,2,3,5,7,11\}$的多周期判别器。推断时不再需要后验编码器：文本编码器给出$\bm{\mu},\bm{\sigma}$，按随机时长展开到帧，从先验采样$f_\theta(\bm{z})$，经逆向流得到$\bm{z}$，再由解码器生成波形。

VITS的优点是单阶段、并行，训练时就让声码器部分见到模型自己的隐变量，消除了两阶段的不匹配，并且不需要外部对齐器，音质在单说话者数据上接近真人录音。缺点是模型和目标复杂，训练成本高；韵律主要由随机时长和隐变量采样决定，不如FastSpeech 2的显式音高、能量条件便于直接调节。

\subsection{基于扩散与流匹配的合成}
\label{sec:nlp-speech-diffusion-tts}

回归式声学模型的根本问题是一对多。设某句话的句末音高可以上扬（取值$+1$）或下降（取值$-1$），两种读法各占一半；以均方误差训练时，最优预测是条件均值0，一条两种读法都不是的平直曲线。FastSpeech 2的音高预测器、Tacotron 2的频谱输出都受此影响，声音因此显得平淡。对抗训练能避免平均，但不稳定且难以覆盖全部模式。扩散与流匹配模型直接学习条件分布本身：从高斯噪声出发，按学得的方向场逐步把样本移到数据分布上，推断时不同的初始噪声给出不同但都合理的读法。其一般推导见第\ref{sec:nlp-speech-vits}小节所引的卷六章，本小节只写条件构造、目标函数和采样步数的权衡。

\paragraph{Grad-TTS：以文本为均值的分数扩散}
Grad-TTS用扩散模型生成梅尔频谱\scite{nlp-speech-popov2021gradtts}。它的编码器部分与Glow-TTS相同：文本编码器为每个音素输出均值$\tilde{\bm{\mu}}_i\in\mathbb{R}^{80}$，训练时用MAS对齐，再展开为帧级均值$\bm{\mu}\in\mathbb{R}^{T\times80}$；推断时用时长预测器展开。与常见扩散模型把数据扩散到标准高斯不同，Grad-TTS的前向过程把频谱$\bm{X}_0$扩散到以文本均值为中心的高斯$\mathcal{N}(\bm{\mu},\bm{I})$：
\begin{equation}
  \mathrm{d}\bm{X}_t=\tfrac12(\bm{\mu}-\bm{X}_t)\,\beta_t\,\mathrm{d}t+\sqrt{\beta_t}\,\mathrm{d}\bm{W}_t,
  \qquad t\in[0,1],
  \label{eq:nlp-speech-gradtts-forward}
\end{equation}
$\beta_t$是噪声调度，$\bm{W}_t$是布朗运动。这个线性随机微分方程有闭式解：给定$\bm{X}_0$，$\bm{X}_t$服从高斯分布，均值为$\bm{\rho}_t=(1-e^{-\frac12\int_0^t\beta_s\mathrm{d}s})\bm{\mu}+e^{-\frac12\int_0^t\beta_s\mathrm{d}s}\bm{X}_0$，协方差为$\lambda_t\bm{I}$，$\lambda_t=1-e^{-\int_0^t\beta_s\mathrm{d}s}$。于是可以一步采样$\bm{X}_t=\bm{\rho}_t+\sqrt{\lambda_t}\bm{\xi}$，$\bm{\xi}\sim\mathcal{N}(\bm{0},\bm{I})$。网络$\bm{s}_\theta(\bm{X}_t,\bm{\mu},t)$（一个U-Net）学习对数密度的梯度，即分数函数；按去噪分数匹配，目标为
\begin{equation}
  L_{\mathrm{diff}}=\mathbb{E}_{\bm{X}_0,t,\bm{\xi}}\Bigl[\lambda_t\Bigl\lVert\bm{s}_\theta(\bm{X}_t,\bm{\mu},t)+\frac{\bm{\xi}}{\sqrt{\lambda_t}}\Bigr\rVert_2^2\Bigr],
  \label{eq:nlp-speech-gradtts-loss}
\end{equation}
$-\bm{\xi}/\sqrt{\lambda_t}$正是高斯条件密度在$\bm{X}_t$处的分数。总目标再加上编码器输出与频谱之间的高斯负对数似然（让$\bm{\mu}$本身就接近频谱）和时长预测损失。推断从$\mathcal{N}(\bm{\mu},\tau^{-1}\bm{I})$采样$\bm{X}_1$（$\tau$为温度），沿反向常微分方程
\[
  \mathrm{d}\bm{X}_t=\tfrac12\bigl(\bm{\mu}-\bm{X}_t-\bm{s}_\theta(\bm{X}_t,\bm{\mu},t)\bigr)\beta_t\,\mathrm{d}t
\]
从$t=1$积分到$t=0$。起点已经是接近目标的文本均值，而不是纯噪声，因此少量步数就能得到可用的频谱。步数$S$是质量与延迟的权衡：每步要完整运行一次U-Net，延迟随$S$线性增长。

\paragraph{F5-TTS：流匹配与填充式训练}
扩散模型的反向轨迹是弯曲的，需要较多步数才能积分准确。\term{条件流匹配}（conditional flow matching）改为直接回归一个把噪声沿直线运到数据的速度场\cite{lipman2023flowmatching}。取噪声$\bm{x}_0\sim\mathcal{N}(\bm{0},\bm{I})$与数据$\bm{x}_1$，定义直线路径与目标速度
\begin{equation}
\begin{aligned}
  \bm{x}_t&=\bigl(1-(1-\sigma_{\min})t\bigr)\bm{x}_0+t\,\bm{x}_1,\qquad
  \bm{u}_t=\bm{x}_1-(1-\sigma_{\min})\bm{x}_0,\\
  L_{\mathrm{CFM}}&=\mathbb{E}_{t,\bm{x}_0,\bm{x}_1}\bigl\lVert\bm{v}_\theta(\bm{x}_t,t,\text{条件})-\bm{u}_t\bigr\rVert_2^2,
\end{aligned}
  \label{eq:nlp-speech-cfm}
\end{equation}
$\sigma_{\min}$是很小的常数。一维教学示例：取$\sigma_{\min}=0$，$x_0=-0.6$，$x_1=1.0$，$t=0.25$时$x_t=0.75\times(-0.6)+0.25\times1.0=-0.2$，目标速度$u_t=1.6$在整条路径上不变；若网络恰好学到它，从$x_0$出发用欧拉法走两步（步长0.5）得$0.2$、$1.0$，正好到达数据点。真实数据的条件速度场是许多直线的平均，不再处处为常数，但仍比扩散轨迹平直，所需步数更少。

Voicebox把流匹配用于语音的\term{填充式训练}（infilling）：给定一段语音及其文本，遮住其中一部分频谱，让模型根据上下文和文本补全被遮的部分\cite{nlp-speech-le2023voicebox}。合成于是变成填空：把参考录音放在前面、目标部分整个遮住，模型续写出的声音自然沿用参考的音色与风格。F5-TTS进一步去掉了显式时长模型与音素对齐\scite{nlp-speech-chen2024f5tts}。它的输入构造如下：参考文本与目标文本的字符序列拼接后，用填充符补齐到与频谱相同的帧数$T$，经几层ConvNeXt V2块\cite{woo2023convnextv2}得到文本表示$\bm{C}\in\mathbb{R}^{T\times d}$；遮蔽标记$\bm{m}\in\{0,1\}^{T}$标出待生成的帧（训练时随机遮住70\%到100\%的帧），可见部分$(1-\bm{m})\odot\bm{x}_1$作为声音条件。三者与当前噪声频谱$\bm{x}_t$在特征维拼接，送入以时间$t$调制层归一化的扩散Transformer\cite{peebles2023dit}，输出速度$\bm{v}_\theta$。损失只在被遮蔽的帧上计算：
\begin{equation}
  L_{\mathrm{F5}}=\mathbb{E}\bigl\lVert\bm{m}\odot\bigl(\bm{v}_\theta(\bm{x}_t,t,(1-\bm{m})\odot\bm{x}_1,\bm{C})-\bm{u}_t\bigr)\bigr\rVert_2^2 .
  \label{eq:nlp-speech-f5-loss}
\end{equation}
字符与帧之间没有任何对齐监督，模型要在训练中自己学会哪些帧读哪个字；填充符只是让两条序列等长。推断时给出参考录音、参考文本和目标文本，总帧数按目标文本与参考文本的字符数之比由参考时长推算，再从噪声出发用欧拉法积分，原文主要报告32步；还可用无分类器引导，把有条件与无条件的速度外推$\bm{v}=\bm{v}_{\mathrm{cond}}+w(\bm{v}_{\mathrm{cond}}-\bm{v}_{\mathrm{uncond}})$，以加强对文本的服从。生成的频谱交给声码器。

扩散与流匹配模型解决了平均化问题，生成的韵律更富变化；填充式训练使零样本音色迁移只需一段参考录音。代价是多步采样增加了延迟，步数与质量需要折中；F5-TTS一类无显式对齐的模型在长句或罕见文本上可能漏读、重复，总时长靠字符数推算也会让语速偏离参考。

\subsection{基于编解码词元的语言模型合成}
\label{sec:nlp-speech-codec-lm-tts}

扩散与流匹配仍在连续频谱上建模。第\ref{sec:nlp-speech-codec}小节的神经编解码器把波形编码为若干并行的离散码流：残差向量量化（式\eqref{eq:nlp-speech-rvq}）的第一级码本刻画主要内容与音色，后续各级逐级补充残差细节。有了离散表示，合成可以改写为词元预测，直接沿用第\ref{chap:nlp-large-language-models}章大语言模型的训练与采样方法；参考录音的词元放在前面，就成了续写的提示。

VALL-E是这一思路的代表\scite{nlp-speech-wang2023valle}。它使用24\,kHz的EnCodec，每秒75帧，每帧8个码本、每个码本1024个码字，记一段语音的码矩阵为$\bm{C}\in\{1,\ldots,1024\}^{T\times8}$，$c_{t,j}$是第$t$帧第$j$级的码字。1秒声音有$75\times8=600$个词元，3秒提示即1800个。若把8级展平成一条序列自回归生成，序列太长；VALL-E按码本层级拆成两个模型。设目标文本的音素为$x$，提示录音的码矩阵为$\tilde{\bm{C}}$，

\begin{equation}
\begin{aligned}
  p(\bm{C}\mid x,\tilde{\bm{C}})
  ={}&\underbrace{\prod_{t=1}^{T}p_{\mathrm{AR}}\bigl(c_{t,1}\mid c_{<t,1},\,\tilde{\bm{C}}_{:,1},\,x\bigr)}_{\text{第一级：自回归}}
  \;\underbrace{\prod_{j=2}^{8}p_{\mathrm{NAR}}\bigl(\bm{C}_{:,j}\mid\bm{C}_{:,<j},\,\tilde{\bm{C}},\,x\bigr)}_{\text{其余各级：非自回归}} .
\end{aligned}
\label{eq:nlp-speech-valle}
\end{equation}
自回归模型是只含解码器的Transformer，输入序列为“音素$x$，提示的第一级码字$\tilde{\bm{C}}_{:,1}$，已生成的第一级码字”，逐帧预测$c_{t,1}$，直到输出结束符；生成多少帧由模型自己决定，因此不需要时长模型。非自回归模型对第$j$级一次性并行预测所有帧：每帧的输入是前$j-1$级码字嵌入之和，加上音素与完整提示，当前层级$j$通过自适应层归一化注入；8级只需额外7次前向。两个模型都以各级词元的交叉熵训练，与语言模型的下一词元预测相同。训练数据是约6万小时、以识别器自动转写的英语有声书录音。

推断时提供约3秒的参考录音及其转写，把参考转写与目标文本的音素拼接为$x$，参考录音的码矩阵作为$\tilde{\bm{C}}$。自回归模型从提示之后续写第一级码字，采用随机采样；束搜索容易陷入无限重复，原文因此改用采样。非自回归模型补全其余7级，最后由EnCodec解码器把码矩阵变回波形。续写出的声音沿用提示中的说话者、录音环境甚至情绪，这就是\term{零样本语音合成}（zero-shot text-to-speech）：目标说话者从未出现在训练集中，只靠上下文提示完成音色迁移。

与两阶段系统相比，编解码词元语言模型有三点不同。第一，它把音色控制变成上下文续写，不必为每个说话者训练或微调模型，也不必单独设计说话者嵌入。第二，它能沿用语言模型的规模化经验，数据从几十小时的精标录音扩大到数万小时的自动转写语音。第三，目标是离散分布上的交叉熵，采样天然给出多样的韵律，不存在回归平均化。代价同样来自自回归与隐式对齐：采样可能跳词、重复或在句末不停止，内容错误比FastSpeech 2一类显式时长方法更难控制；生成速度受第一级逐帧串行的限制；与参考录音高度相似的声音也带来冒用身份的风险，需要配合水印或说话者授权检查。此后的方案常把两类方法组合：CosyVoice先用语言模型生成从识别监督学到的语义词元，再由流匹配模型把语义词元转为频谱\scite{nlp-speech-du2024cosyvoice}，语言模型负责内容与韵律的离散规划，流匹配负责连续的声学细节。离散声学词元同样是第\ref{sec:nlp-speech-realtime-dialogue}节语音大模型的输出接口。

\begin{table*}[htbp]
  \centering
  \begin{tabularx}{\linewidth}{@{}>{\raggedright\arraybackslash}p{22mm}>{\raggedright\arraybackslash}p{26mm}>{\raggedright\arraybackslash}p{17mm}>{\raggedright\arraybackslash}X>{\raggedright\arraybackslash}p{18mm}>{\raggedright\arraybackslash}X@{}}
    \toprule
    方法 & 对齐方式 & 自回归 & 并行度与推断代价 & 零样本 & 内容错误风险\\
    \midrule
    Tacotron 2 & 位置敏感注意力，隐式 & 逐帧 & 帧数次串行；另需声码器 & 否 & 注意力跳读、重复、不停止\\
    FastSpeech 2 & 外部强制对齐的显式时长 & 否 & 一次前向并行；另需声码器 & 否 & 低；依赖对齐器质量\\
    VITS & 训练中单调对齐搜索 & 否 & 一次前向直接得到波形 & 否 & 低；单调约束保证不跳读\\
    Grad-TTS & 单调对齐搜索与时长 & 否 & 多步去噪，延迟随步数增长 & 否 & 低；对齐显式\\
    F5-TTS & 无显式对齐，填充符等长 & 否 & 约数十步流匹配积分 & 是，参考录音填充 & 中；长句可能漏读\\
    VALL-E & 无显式对齐，模型决定长度 & 首层逐帧 & 首层帧数次串行，其余7级并行 & 是，3秒提示 & 较高；跳词、重复、不停止\\
    \bottomrule
  \end{tabularx}
  \caption{语音合成方法对照。自上而下，对齐由隐式注意力走向显式时长，再由单调搜索走回不显式对齐的大规模生成；零样本能力来自填充式训练或上下文提示。内容错误风险一栏是按机制作出的定性判断，实际程度依赖数据规模与解码设置，需在同一测试文本上用内容核验比较。}
  \label{tab:nlp-speech-tts-methods}
\end{table*}
\FloatBarrier

不论采用哪一类方法，交付的声音都要按前文的四个字段核验内容。若只把已合成通知中的“十七日”改成“十八日”，还需声明修改区间与需要保持的前后区间：新日期时长不同，后续边界要重新对齐，不能把局部替换后的旧时间戳原样保留。音色相似、频谱接近或自然度高，都不能替代内容核验。

\section{语音翻译}
\label{sec:nlp-speech-translation}

\term{语音翻译}（speech translation）把源语言声音转换为目标语言文字或声音，核心约束仍是第\ref{chap:nlp-text-generation}章讨论的意义保持。描述一个语音翻译系统要分开三个维度：级联或直接是实现路线，目标文字或目标声音是交付对象，整段或流式是可见信息与时间条件。三者可以自由组合，例如直接模型也可以流式输出文字，级联系统也可以交付声音。

这一任务的挑战来自声音与翻译两方面。第一，识别错误会传进翻译，翻译器通常无法从错误转写中恢复源语原意。第二，世界上相当多的语言没有通行的书面形式，以源语转写为中间表示的路线对它们不可用。第三，“源声音—目标声音”的成对数据远少于平行文本和转写录音。第四，同一译文可以用许多种时长与韵律读出，直接回归目标频谱会把这些读法平均成模糊的声音。第五，流式译出必须在源句未完时决定等待还是提交。方法从级联转换出发：一最佳级联最简单，多候选传递缓解错误传播，增量级联处理时间条件。直接转换随后依次改变中间表示：Translatotron直接回归目标频谱；S2UT改为预测离散声音单位；UnitY先生成目标文字再生成单位，SeamlessM4T据此构建多语模型；SeamlessStreaming再加入单调注意力，形成流式直接翻译。表\ref{tab:nlp-speech-translation-routes}在节末对照这些路线。

\subsection{级联转换}
\label{sec:nlp-speech-cascade}

级联把识别、文本翻译与合成三个已有部件接起来，各部件可以分别利用大规模转写录音、平行文本与合成语料训练，中间文字也便于人工检查。代价同样来自串联：错误逐级传播，语气与情绪等副语言信息在文字接口处丢失，三段延迟相加。CoVoST 2提供多语种的声音—转写—译文三元组，可以同时监督识别、翻译与声音条件的翻译，常用来比较级联与直接路线\scite{nlp-intro-wang2021covost}

\subsubsection{一最佳识别–翻译–合成}
\label{sec:nlp-speech-one-best}

设源声音为$\bm{X}$，源语文字为$x$，译文为$y$。一最佳级联依次求
\[
 \hat x=\argmax_x p_{\mathrm{asr}}(x\mid\bm{X}),
 \qquad
 \hat y=\argmax_y p_{\mathrm{mt}}(y\mid\hat x),
\]
再把$\hat y$连同目标语言和音色条件交给第\ref{sec:nlp-speech-tts}节的合成器。两次$\argmax$都由束搜索近似。三个部件分别用转写录音、平行文本和目标语言的文字—声音配对训练，接口需要约定语言标记、分段、标点与专名写法。

通知中“改至四月十八日下午两点”应译为“rescheduled to April 18 at 2 p.m.”。若识别输出“四月十七日”，翻译器忠实译成“April 17”，翻译部件本身没有出错，端到端结果却错了。把人工正确转写与识别转写分别送进同一个翻译器，就能把误差归到识别还是翻译；合成器若把“2 p.m.”读成上午两点，则归为目标声音内容错误。三类错误的定位方法不同，验收时须分别检查译文和听者实际听到的声音。

\subsubsection{识别多候选传递与重排}
\label{sec:nlp-speech-nbest-translation}

一最佳传递在识别阶段就丢掉了其他解释。若声学上“十七”与“十八”都可能，可以保留识别的$n$个候选，每个候选各自翻译，再按识别分数、翻译分数与可核验约束的加权和重排：
\begin{equation}
 S(x,y)=\lambda_a\log p_{\mathrm{asr}}(x\mid\bm{X})
        +\lambda_m\log p_{\mathrm{mt}}(y\mid x)+r(x,y).
 \label{eq:nlp-speech-nbest}
\end{equation}
$r(x,y)$检查日期、实体与格式在源文和译文间是否一致，权重$\lambda_a,\lambda_m$在开发集上按相同的候选预算调定。也可以传递词格，让翻译器在更多候选上搜索。

完整的概率处理是对源语转写边缘化，$p(y\mid\bm{X})=\sum_xp_{\mathrm{mt}}(y\mid x)p_{\mathrm{asr}}(x\mid\bm{X})$。$n$最佳列表只覆盖少数$x$，而式\eqref{eq:nlp-speech-nbest}对每个$(x,y)$对单独打分，近似的是联合最大化而不是边缘化；两个源候选译成同一译文时，二者的分数并不相加。教学例子：“十八日”与“十七日”的识别概率为$0.55,0.45$，各自最佳译文的条件概率为$0.60,0.90$。取$\lambda_a=\lambda_m=1$、$r=0$时，乘积分别为$0.33$与$0.405$，后者获胜。翻译模型对某一说法的偏好压过了较好的声学证据；规范值检查可以拒绝源译不一致的组合，却无法判断哪个日期确实被说出。

\subsubsection{增量级联与等待策略}
\label{sec:nlp-speech-incremental-cascade}

用户尚未说完时，级联的每一环都只能看到前缀。识别器的输出分为已提交的稳定前缀和仍会改写的暂定尾部（第\ref{sec:nlp-speech-whisper-streaming}小节）；只有稳定前缀可以支撑不可撤回的译出，暂定文字只用于预计算。文本端的读写策略沿用第\ref{sec:nlp-generation-wait-k}小节：wait-$k$固定落后$k$个源词，第\ref{sec:nlp-generation-adaptive-policy}小节的自适应策略按已读内容决定读或写。在语音级联中，“读一个源词”要等识别器提交这个词，所以源词数不能直接换算成固定毫秒。

例如用户说“校外来宾可以直接参加……吗？不，须先报名”。过早译出并播放“可以直接参加”，就把尚未完成的询问当成了许可；问句标记和否定都在后面才到。翻译部件可以用前缀到前缀的数据训练，也可以在整句模型外加提交策略。目标文字未展示前可以重算，已播放的声音只能追加更正。

延迟须分段记账。设用户在墙钟3.0\,s说完，最终源文3.4\,s提交，译文3.7\,s确定，首段目标声音4.0\,s开始播放。用户结束到首段播放为1.0\,s，其中翻译只占0.3\,s，其余是识别等待和合成排队。这是计时演算，不是系统实测。

\subsection{直接转换}
\label{sec:nlp-speech-direct-translation}

直接转换用一个模型从源声音生成目标输出，不调用外部识别器取得源语转写。它要解决级联的三个问题：识别错误的传播，副语言信息的丢失，以及无书面形式的语言无法转写。新的困难是端到端配对数据少，目标声音又难以直接回归。下面四个模型依次改变目标表示与生成顺序：Translatotron直接预测频谱；S2UT把目标声音换成离散单位；UnitY在单位之前先生成目标文字；SeamlessStreaming加入单调读写决策。与级联比较时，须固定语言方向、训练资源、可见上下文和计算预算。

\subsubsection{直接语音到语音翻译：Translatotron与离散单位翻译}
\label{sec:nlp-speech-translatotron}

\term{直接语音到语音翻译}（direct speech-to-speech translation, S2ST）的第一个完整模型是Translatotron\scite{nlp-speech-jia2019translatotron}。它把第\ref{sec:nlp-speech-tacotron}小节的Tacotron 2从“文字到频谱”改成“源频谱到目标频谱”。输入是80维log-Mel帧，相邻3帧堆叠后送入8层双向LSTM编码器；解码器沿用Tacotron 2的预网络、自回归LSTM栈与后处理网络，用4头加性注意力读取编码器输出，每步预测2帧1025维对数线性频谱，最后由声码器（实验中主要用Griffin-Lim，听感评价用WaveRNN）恢复波形。可选的说话者编码器从源录音提取音色嵌入，使译文声音保留原说话者的音色。

只用频谱回归误差训练时，注意力很难学会对齐：源频谱与目标频谱属于两种语言，相邻帧之间没有单调对应。Translatotron因此在编码器中间层上接两个辅助解码器（各为2层LSTM与单头加性注意力），分别预测源语音素与目标音素序列，总损失为
\begin{equation}
 \mathcal L_{\mathrm{Tr}}=\mathcal L_{\mathrm{spec}}
 +\lambda_{\mathrm{src}}\mathcal L_{\mathrm{CE}}(\text{源音素})
 +\lambda_{\mathrm{tgt}}\mathcal L_{\mathrm{CE}}(\text{目标音素}),
 \label{eq:nlp-speech-translatotron}
\end{equation}
$\mathcal L_{\mathrm{spec}}$是预测频谱与目标频谱的回归误差。辅助解码器读取的层是超参数：在较小的Fisher西英数据上，源音素接第4层、目标音素接第6层，使较低层编码源语内容、较高层承担翻译；在较大的会话数据上两者都接最后一层。辅助解码器只在训练时使用。原文在西班牙语—英语会话数据上报告，不加辅助损失时BLEU只有0.4，加入两项辅助损失后为42.7，同条件的语音翻译加合成级联为48.7。直接模型可行，但仍落后于级联，并且训练依赖音素标注。

频谱回归的根本困难在于目标的多样性。同一译文可以快读、慢读或在不同位置停顿，回归误差的最优解是这些读法的条件均值，生成的频谱因而发糊；自回归频谱解码器还会出现跳读与重复。Lee等人的\term{语音到单位翻译}（speech-to-unit translation, S2UT）把目标换成离散符号，以分类代替回归\scite{nlp-speech-lee2022s2ut}。目标单位沿用第\ref{sec:nlp-speech-hubert}小节HuBERT构造伪标签的办法：用在LibriSpeech上预训练的HuBERT Base提取目标英语录音第6层的帧表示$\bm h_t$（每20\,ms一帧），做$K=100$的K均值聚类，每帧的单位就是最近中心的编号$z_t=\argmin_{c}\lVert\bm h_t-\bm\mu_c\rVert_2$。相邻帧常落在同一类，S2UT的“去重”做法把连续重复的单位合并为一个，例如帧单位序列$17,17,17,52,52,8$变为$17,52,8$，并记下时长$(3,2,1)$。去重序列更短，也不再携带时长，时长交给声码器预测。

S2UT模型是改造的Transformer翻译模型。编码器前是两层一维卷积（核宽5、1024通道、每层步长2并接GLU），把10\,ms一帧的源特征降采样4倍；其后是12层Transformer编码器（$d=256$，4头），6层Transformer解码器自回归预测去重单位，用带标签平滑的交叉熵训练：
\begin{equation}
 \mathcal L_{\mathrm{s2u}}=-\sum_{i=1}^{L}\log p(u_i\mid\bm X,u_{<i}).
 \label{eq:nlp-speech-s2ut}
\end{equation}
有目标文字时，可在单位解码器第3层上接CTC（第\ref{sec:nlp-speech-ctc}小节）预测目标子词，同一次解码同时得到单位与文字；源、目标音素或字符辅助任务沿用Translatotron的接法。单位声码器沿用Polyak等人改造的HiFi-GAN（第\ref{sec:nlp-speech-hifigan}小节）：输入单位嵌入序列，前置一个FastSpeech 2式时长预测器（两层一维卷积加线性层），在对数域上以均方误差拟合每个单位的真实帧数，并与HiFi-GAN的对抗损失联合训练\cite{nlp-speech-polyak2021resynthesis}。合成时先预测时长，把单位按时长展开，再生成波形。

与Translatotron相比，S2UT改动了两处。目标分布由连续回归变为每步$K$类的分类，类别分布可以多峰，不会把不同读法平均掉；时长与韵律的细节交给声码器，翻译模型只需决定说什么。单位由自监督模型从无标注录音得到，因此S2UT在完全不用目标文字的条件下也能训练，原文报告这时的效果与使用文字监督的频谱模型相当。代价是单位只保留语言内容，原说话者的音色与情绪基本丢失；单位声码器为每种目标语言或音色单独训练，单位序列的质量也受聚类数与所选层的限制。

\subsubsection{文本与声音单位生成：SeamlessM4T}
\label{sec:nlp-speech-seamless}

S2UT的单个解码器要同时完成翻译与发音，单位序列又比文字长得多，翻译的长距离依赖被摊在几百步上。Translatotron 2先生成音素再合成频谱，提出了两遍解码的思路\cite{nlp-speech-jia2022translatotron2}；UnitY把第一遍改为生成目标子词，把第二遍改为生成离散单位\scite{nlp-speech-inaguma2023unity}。UnitY由四个模块组成。Conformer语音编码器（第\ref{sec:nlp-speech-encoder}小节）把源声音编码为$\bm H$。第一遍文本解码器以交叉注意力读取$\bm H$，自回归生成目标子词$Y=(y_1,\ldots,y_M)$，其第$i$步输出投影前的状态为$\bm d^{\mathrm{text}}_i=\mathrm{TDec}(\bm H,y_{<i})\in\mathbb R^{d}$。双向的文字到单位编码器把这些状态改写为$\bm Z=\mathrm{T2UEnc}(\bm D^{\mathrm{text}})\in\mathbb R^{M\times d}$，长度不变，用来弥合文字与单位两种表示之间的差距。第二遍单位解码器只读取$\bm Z$，不再读取$\bm H$，自回归生成单位$U=(u_1,\ldots,u_L)$。训练目标为两遍交叉熵的加权和：
\begin{equation}
 \begin{aligned}
 \mathcal L_{\mathrm{UnitY}}
 &=\mathcal L_{\mathrm{s2u}}(U\mid\bm X,Y)+w_{\mathrm{s2t}}\mathcal L_{\mathrm{s2t}}(Y\mid\bm X),\\
 \mathcal L_{\mathrm{s2t}}&=-\frac1M\sum_{i=1}^{M}\log p(y_i\mid\bm d^{\mathrm{text}}_i),\\
 \mathcal L_{\mathrm{s2u}}&=-\frac1L\sum_{i=1}^{L}\log p(u_i\mid\mathrm{UDec}(\bm Z,u_{<i})).
 \end{aligned}
 \label{eq:nlp-speech-unity}
\end{equation}
训练时第一遍用真实译文作教师强制，第二遍读取的是真实译文对应的解码器状态。推断做两次束搜索：先以较宽的束得到最佳译文$\hat Y$，把其解码器状态送入文字到单位编码器，再以较窄的束（原文默认取1）生成单位，最后由单位声码器合成。第二遍只看$\bm Z$，使声音单位与已经输出的译文保持同步，代价是第二遍无法修正第一遍的翻译错误。

两遍结构的好处有三点。翻译在短的子词序列上完成，可以用无标注文本预训练第一遍解码器（UnitY用多语去噪自编码预训练的mBART初始化它）；单位解码器只做“已知译文怎样发音”，任务接近合成；第一遍的文字本身就是可读的译文。原文在相同数据上报告，UnitY比单遍S2UT高2.5至4.2个ASR-BLEU，解码快2.83倍。

SeamlessM4T用多任务的UnitY构建覆盖近百种语言的模型\scite{nlp-intro-seamless2023}。语音编码器是24层Conformer的w2v-BERT 2.0，约6亿参数，在约100万小时无标注录音上自监督预训练；文本编码器与第一遍解码器来自NLLB文本翻译模型，使语音输入与文字输入共享同一个译文解码器。目标单位取自XLS-R-1B第35层的50\,Hz表示，K均值聚类（$K=10000$）；多语单位声码器是加入语言嵌入与时长预测器的HiFi-GAN。成对的“源声音—目标单位”数据不足时，先由文本翻译数据的译文经文字到单位模型生成伪单位标签，再训练第二遍。各部分的监督来源因此不同：编码器来自无标注录音，翻译能力来自平行文本与声音—译文对，单位来自目标录音的聚类，声码器来自单位与真实波形的配对。

两遍输出需要分别核对。对“海棠楼201”，第一遍可在既定译名表下输出“Haitang Building, Room 201”；若第二遍把数字读成210，只看文字的评价会漏掉这个错误。ASR-BLEU先用识别器转写目标声音再算BLEU，它能发现读错的数字，但同时混入了评价识别器自身的错误。

\subsubsection{单调推进与流式译出：SeamlessStreaming}
\label{sec:nlp-speech-seamless-streaming}

SeamlessM4T等整句到齐后才开始解码。SeamlessStreaming在SeamlessM4T v2上加入\term{高效单调多头注意力}（efficient monotonic multihead attention, EMMA），让文字解码器边读源声音边决定何时写；生成的文字再交给非自回归的UnitY2文字到单位部件，按块产生声音单位\cite{nlp-speech-seamless2023streaming,nlp-speech-ma2023emma}。

EMMA沿用第\ref{sec:nlp-generation-adaptive-policy}小节的单调注意力与单调多头注意力：每个头对目标位置$i$和源位置$j$估计停止概率$p_{i,j}$，期望对齐$\alpha_{i,j}$满足式\eqref{eq:nlp-generation-monotonic-attention}的递推。递推原有的并行闭式解要除以$\prod_{\ell}(1-p_{i,\ell})$的累积乘积，源序列较长时这个乘积趋于零，训练中出现数值不稳定和对齐消失。语音的源序列比文字长得多，问题更突出。EMMA把递推改写为不含除法的矩阵形式：
\begin{equation}
 \begin{gathered}
 \bm\alpha_{i,:}=\bm p_{i,:}\odot\bigl(\bm\alpha_{i-1,:}\bm T^{(i)}\bigr),\\
 T^{(i)}_{m,n}=\begin{cases}\prod_{\ell=m}^{n-1}(1-p_{i,\ell}),&m<n,\\1,&m=n,\\0,&m>n.\end{cases}
 \end{gathered}
 \label{eq:nlp-speech-emma}
\end{equation}
$\bm\alpha_{i,:}\in\mathbb R^{S}$是第$i$个目标单位在$S$个源位置上的对齐分布，$T^{(i)}_{m,n}$是“上一步停在$m$、本步从$m$一路不停走到$n$”的概率。上三角矩阵$\bm T^{(i)}$可由$1-p_{i,:}$的累积乘积直接构造，各目标位置可以并行计算。

训练目标在翻译负对数似然上加两项正则。期望延迟$\bar d_i=\sum_j j\,\alpha_{i,j}$与对齐方差$\bar v_i=\sum_j j^2\alpha_{i,j}-\bar d_i^{\,2}$都可微，
\begin{equation}
 \mathcal L_{\mathrm{EMMA}}=-\log p(Y\mid\bm X)+\lambda_{\mathrm{lat}}\,\mathcal C(\bar d_1,\ldots,\bar d_M)+\lambda_{\mathrm{var}}\sum_{i}\bar v_i,
 \label{eq:nlp-speech-emma-loss}
\end{equation}
$\mathcal C$是某种延迟指标的可微形式。延迟项与第\ref{sec:nlp-generation-adaptive-policy}小节MILk的延迟损失作用相同，防止模型退化为等全部声音到齐再写；方差项让对齐集中，使推断时的确定性读写决策与训练时的期望对齐接近。

推断时，每当一块新声音到达，模型重算已见声音的编码，对所有头求停止概率，取其中的最小值与阈值比较：最小值超过阈值才写出一个目标词元，否则继续读。教学设定三个头的概率为$(0.7,0.6,0.3)$、阈值0.5，最小值0.3，模型等待；更多声音到达后变为$(0.8,0.7,0.6)$，才写出。最保守的头决定推进，这与单调多头注意力“所有头都停止才写”的规则一致。

写出文字不等于立即播放。新文字经文字到单位部件变成单位，单位块达到最小长度才送入声码器。延迟因此要分两种口径：以“写出第$i$个目标词元时已读入多少秒源声音”计的滞后，描述信息等待；以墙钟测得的首段播放与播放结束偏移，还包含编码重算、单位生成、网络与队列时间。

\begin{table*}[htbp]
  \centering
  \begin{tabularx}{\linewidth}{@{}>{\raggedright\arraybackslash}p{29mm}>{\raggedright\arraybackslash}p{26mm}>{\raggedright\arraybackslash}X>{\raggedright\arraybackslash}X>{\raggedright\arraybackslash}X@{}}
    \toprule
    路线 & 中间表示 & 监督来源 & 延迟来源 & 可检查性\\
    \midrule
    一最佳级联 & 源语转写、译文文字 & 转写录音、平行文本、合成语料分别训练 & 三段串联相加；流式时还要等识别提交 & 两处文字接口都可检查，错误可逐段归因\\
    多候选级联 & $n$最佳列表或词格 & 同上，另需开发集调权重 & 每个候选各翻译一次 & 可查看候选与重排分数\\
    Translatotron & 无显式中间表示，直接回归频谱 & 源—目标声音对，辅助音素标注 & 整句编码后自回归逐帧生成 & 只能借辅助解码器观察内部\\
    S2UT & 去重的HuBERT聚类单位 & 声音对；单位由无标注录音聚类得到，可不用文字 & 单位序列长，声码器另需时长预测 & 单位不可读；有文字时可加CTC输出\\
    UnitY与SeamlessM4T & 目标子词，再接声音单位 & 预训练编码器、平行文本、伪单位标签 & 两遍束搜索，第二遍可窄束 & 第一遍译文可读，声音需另行转写核对\\
    SeamlessStreaming & 同上，加单调对齐 & 同上，再加延迟与方差正则 & 由阈值和最保守的头决定读写；单位块需达到最小长度 & 读写时刻可记录，已播声音不可撤回\\
    \bottomrule
  \end{tabularx}
  \caption{语音翻译路线的对照。从上到下，中间表示由显式源语文字逐步变为目标文字加离散单位，监督也从各部件分别训练变为预训练表示与联合目标。表中延迟一列只列结构上的等待来源，不是实测数值。}
  \label{tab:nlp-speech-translation-routes}
\end{table*}

\section{语音大模型与实时口语对话}
\label{sec:nlp-speech-realtime-dialogue}

前面几节的识别、理解、合成与翻译模型可以串成一个口语对话系统，但这样的系统有三处缺口。其一，理解能力局限于各部件的训练任务，转写之外的语气、情绪和环境声在文字接口处丢失。其二，识别、文本生成、合成串联，延迟逐段相加。其三，系统按话轮交替工作，无法边听边说，被打断后也要靠外部检测才停下。\term{全双工}（full duplex）指双方可以同时发送和接收声音的交互条件，它要求模型在说话时仍持续读取用户的声音。

方法谱系沿两条线推进：声音怎样进入语言模型，以及声音怎样从模型产生。表\ref{tab:nlp-speech-dialogue-routes}列出三类路线。级联系统用识别器把声音变成文字，交给文本对话模型，再用合成器读出；\term{半级联}（half-cascade）系统把语音编码器直接接到大语言模型上，声音特征进入模型，输出仍是文字，再经外部合成；全双工联合模型把用户与系统两路声音都离散成词元，用同一个模型在共同时间轴上持续建模。半级联解决声音怎样进入模型；级联与半级联仍要靠外部的话轮检测与响应门控决定何时说，这两部分在全双工模型之前讨论；无论哪条路线，声音回复都要与业务动作协调。

\begin{table*}[htbp]
  \centering
  \begin{tabularx}{\linewidth}{@{}>{\raggedright\arraybackslash}p{20mm}>{\raggedright\arraybackslash}p{30mm}>{\raggedright\arraybackslash}X>{\raggedright\arraybackslash}X>{\raggedright\arraybackslash}X>{\raggedright\arraybackslash}X@{}}
    \toprule
    路线 & 组成 & 延迟来源 & 打断处理 & 副语言信息 & 可检查性\\
    \midrule
    级联 & 识别器—文本对话模型—合成器 & 等识别提交、文本生成、合成首段，三段相加 & 依赖外部VAD与轮次检测，停播由控制逻辑执行 & 在转写处丢失，合成端另设语气 & 转写与回复文字都可检查\\
    半级联 & 语音编码器—适配器—大语言模型，外接合成 & 编码整段声音后生成文字，再加合成 & 同级联，仍需外部检测 & 输入端保留，输出端仍由合成器决定 & 回复文字可检查，输入只能看模型表现\\
    全双工联合 & 两路声音词元与文字流，单一模型逐帧生成 & 每步一个编解码帧加声学延迟 & 用户声音持续输入，模型自行停止或交叠 & 输入与输出都在声音词元中 & 文字流可读；工具与业务状态需外接\\
    \bottomrule
  \end{tabularx}
  \caption{口语对话的三类路线。从上到下，声音在模型内部保留得越来越多，延迟中的串联部分越来越少，但对话行为也越来越依赖模型自身，越难用显式规则控制。表中延迟一列只列结构来源，不是实测数值。}
  \label{tab:nlp-speech-dialogue-routes}
\end{table*}

\subsection{语音编码器与语言模型的连接}
\label{sec:nlp-speech-llm-connector}

大语言模型已经具备指令理解与推理能力，转写却会丢掉语气、情绪与非语音事件。半级联路线的思路是不经过文字，直接把声音的连续表示接到语言模型的输入嵌入空间。困难在于两种表示的长度与分布都不匹配：Whisper编码器（第\ref{sec:nlp-speech-whisper}小节）每20\,ms输出一帧，30\,s录音有1500帧，而同样长度的转写只有一两百个词元；编码器的特征空间也与语言模型的词嵌入空间无关。

连接结构由三部分组成。预训练语音编码器把声音编码为$\bm H\in\mathbb R^{T\times d_a}$；适配器（adapter，指编码器与语言模型之间的连接模块，不同于第\ref{sec:nlp-sft-peft}节插在层内的适配器）把帧率降到每秒十余个到数十个向量，并投影到语言模型的宽度$d_{\mathrm{LM}}$；得到的音频嵌入与文字指令的词嵌入拼成一条序列送入语言模型。常见的降采样方式有两类。第一类是拼接或池化：把相邻$k$帧拼成一个$k d_a$维向量，再用两层MLP投影，
\begin{equation}
 \bm e_m=\mathrm{MLP}\bigl([\bm h_{(m-1)k+1};\ldots;\bm h_{mk}]\bigr)\in\mathbb R^{d_{\mathrm{LM}}},\qquad m=1,\ldots,\lceil T/k\rceil.
 \label{eq:nlp-speech-adapter-stack}
\end{equation}
第二类是查询读取：Q-Former用一组可学习查询向量，以交叉注意力从编码器特征中读出固定数目的输出向量。整段只用一组查询会丢失时间顺序，语音中常改为窗口级：把$\bm H$切成长$L$帧的窗口，每个窗口用$N$个查询读出$N$个向量，总长度为$N\lceil T/L\rceil$，保持单调的时间顺序。视觉—语言模型中的Q-Former与投影连接器采用同一框架：冻结的感知编码器后接一个小型可训练连接模块，把可变长特征变成语言模型可读的嵌入，可参考\BookChapterRef{chap:vision-multimodal-models}{卷十的多模态模型章}。

训练目标与第\ref{chap:nlp-supervised-finetuning}章的监督微调相同，只对回答词元计算下一词元交叉熵（式\eqref{eq:nlp-sft-masked-loss}），音频嵌入和指令只作条件。训练通常分阶段进行。先在识别、音频描述等大规模“声音—文字”配对上训练连接部分，使音频嵌入对齐到语言模型能读懂的空间；再在问答、翻译、情绪识别等多任务指令数据上微调，使模型按指令处理声音。

Qwen-Audio与SALMONN是这一路线的两种实现\cite{nlp-speech-chu2023qwenaudio,nlp-speech-tang2024salmonn}。Qwen-Audio用Whisper-large-v2初始化音频编码器（32层，约6.4亿参数），其后接步长2的池化层，使每个输出向量对应约40\,ms声音，30\,s录音得到750个向量，再接Qwen-7B语言模型。它在三十多种音频任务上联合训练，为减轻不同数据集标注格式的相互干扰，在输出前加入分层任务标签：转写或分析标记、音频语言、任务类型、输出文字语言、是否输出时间戳，最后是具体的输出格式说明。共享的标签促使相近任务共享知识，不同的标签把格式冲突的数据区分开。多任务预训练阶段冻结语言模型、只训练音频编码器；随后的对话版本Qwen-Audio-Chat冻结编码器、只微调语言模型。

SALMONN同时接入两个编码器：Whisper-large-v2擅长语音，通用音频编码器BEATs擅长环境声与音乐。两者都输出50\,Hz特征，逐帧拼接后送入窗口级Q-Former，窗口长$L=17$帧（约0.33\,s），每窗$N=1$个查询，30\,s录音得到88个向量。语言模型是冻结的Vicuna-13B，只在其自注意力的查询与值矩阵上加秩为8的LoRA（第\ref{sec:nlp-sft-peft}节），可训练参数约3300万。训练分三阶段：先用识别与音频描述数据预训练，再用十余种任务的指令数据微调，最后做“激活微调”。第三阶段针对的问题是任务过拟合：模型在前两阶段见到的多为短而固定的转写与描述，遇到开放指令时仍倾向于输出转写，忽略指令本身。SALMONN在推断时降低LoRA缩放系数，得到较长、较多样的回答，再用少量这样生成的样本微调。

两种设计的差别落在降采样位置与训练分工上。池化保留每一帧附近的全部信息，向量数较多；窗口级Q-Former把每0.33\,s压成一个向量，序列短，但细粒度的发音信息可能丢失。Qwen-Audio训练编码器而冻结语言模型，适合大规模多任务对齐；SALMONN冻结两侧主体，只训练连接器与LoRA，成本低。半级联模型擅长理解：可以回答“说话人是否生气”“背景里是什么声音”这类转写无法回答的问题。它的输出仍是文字，声音回复需要外部合成或第\ref{sec:nlp-speech-tts}节的词元化合成；Whisper类编码器以整段为单位编码，也不适合边听边答。

\subsection{输入流、轮次与打断}
\label{sec:nlp-speech-turns}

无论哪条路线，系统都要判断三件不同的事：现在有没有人说话，用户的话轮是否结束，这句话起什么作用。“帮我订……明天下午”中的停顿可能是思考；系统播报时用户说“等等”，是要求接管或更正；短促的“嗯”则可能只表示在听。在线输入需要保留双方声音的时间轴、系统自身的播放状态与对话历史，监督数据也分别标注声音活动区间与轮次事件。下面先介绍只看已发生声音的检测，再介绍预测未来说话者活动的模型。

\subsubsection{VAD与轮次结束检测}
\label{sec:nlp-speech-vad}

\term{语音活动检测}（voice activity detection, VAD）逐帧判断是否存在语音。最简单的能量基线对第$i$个短窗计算$E_i=w^{-1}\sum_{n=0}^{w-1}x_{ih+n}^2$，其中$w$为窗长、$h$为帧移。$E_i$超过启动阈值$\theta_{\mathrm{on}}$时进入活动状态，连续若干窗低于结束阈值$\theta_{\mathrm{off}}$才退出；取$\theta_{\mathrm{on}}>\theta_{\mathrm{off}}$形成滞回，避免在阈值附近来回抖动。能量高的可能是噪声，能量低的可能是轻声，所以实际系统多用神经VAD：以有声、无声帧标注训练一个小型帧分类器，再做同样的滞回平滑。

轮次结束检测要回答的是话是否说完，它需要声音活动以外的证据。一个常用实现是“听取—候选结束—等待—确认结束”状态机：VAD转为静音时进入候选结束，再结合静音时长、句末韵律下降与语义完整度决定是否确认；若用户在等待期间恢复发声，就回到听取，撤销尚未提交的结束判断。教学设定阈值为400\,ms：用户在“帮我订”后停顿300\,ms再说“明天下午”，静音未达阈值，系统继续听；即使停顿达到500\,ms，“帮我订”缺少预约对象，语义完整度低，系统应发起澄清而不是执行。评价时分别统计过早截断（用户没说完就开始回答）与结束等待（用户说完到系统确认的时间），两者此消彼长，由阈值调节。系统自己的播放声音也会回到麦克风，回声消除或双通道录音决定VAD实际接收什么。

\subsubsection{未来声音活动预测：VAP}
\label{sec:nlp-speech-vap}

静音时长只是已经发生的证据，人在对话中却会提前预判对方何时说完。\term{声音活动预测}（voice activity projection, VAP）把轮次控制改写为预测：根据已见声音与双方活动历史，预测未来2\,s内两人各自的活动模式\scite{nlp-speech-ekstedt2022vap}。原模型把单声道混合录音送入预训练CPC编码器，得到100\,Hz的256维帧表示；再加上双方当前的活动标记与过去若干时间窗的活动比例，送入4层、宽256的因果Transformer。

预测目标是离散的联合状态。未来2\,s划为200、400、600、800\,ms四个由细到粗的区间，每人每个区间按活动比例是否超过阈值记为0或1，两人共8位，对应$2^8=256$个联合状态。模型输出256类分布，用交叉熵训练；标签由录音的后续活动自动构造，不需要人工标注，模型输入只含当前与过去，不会泄漏未来。与对8个区间分别做二分类相比，联合分布保留了两人活动之间的依赖，例如“A停、B接”与“A、B同时停”可以有不同的概率。

推断时，把与某个事件相符的状态概率相加，就得到该事件的概率，不需要为每种事件单独训练。在双方都静音的时刻，比较“只有A将说话”与“只有B将说话”两组状态的总概率，就得到保持话轮或移交话轮的判断；对“对方只在近期区间短暂活动”的状态求和，得到短回应的概率。教学例子：设符合移交、保持与其他模式的总概率为$0.55,0.35,0.10$，在“移交或保持”二选一的条件下，移交概率为$0.55/(0.55+0.35)\approx0.611$，而不加条件的移交概率仍是$0.55$；报告时须说明用的是哪一种。VAP给出的是何时可能轮到系统说话，它不判断预约信息是否齐全，后者要由下一小节的语义状态决定。

\subsection{增量理解与响应选择}
\label{sec:nlp-speech-incremental-understanding}

流式识别会把“明天”改成“后天”，转写不是只追加的日志。每条局部转写或语义记录都带版本号、对应的源声音区间，以及已确认前缀与暂定尾部。状态更新按新增、替换与撤回三种操作进行：版本1暂记4月18日；版本2听到“不，后天”后改为4月19日，依赖旧日期且尚未执行的动作同时失效。原话可以保留供追溯，当前有效值只有一份。在这一状态上，系统要在等待、短确认、澄清与正式回答之间选择响应。

\subsubsection{状态驱动的响应门控}
\label{sec:nlp-speech-response-gate}

响应门控按状态与事件选择下一个动作。规则实现依次检查五个条件：语义是否完整，关键字段是否确认，用户是否交出话轮（来自上一小节的轮次检测或VAP），工具是否返回，以及剩余延迟预算。算法\ref{alg:nlp-speech-response-gate}给出一个教学版本。规则来自业务前置条件与交互协议；有标注的历史对话时，也可以在同一动作集上训练分类器，或让语言模型提出动作，再经过同样的字段、权限与结果检查。

\begin{algorithm}[状态驱动的响应门控（教学版本）]
  \label{alg:nlp-speech-response-gate}
  \begin{algorithmic}[1]
    \Require 语义状态$\mathcal S$（字段、版本、确认标记），轮次信号$q\in\{\text{未结束},\text{结束}\}$，在途工具请求集合$R$，剩余等待预算$b$
    \Ensure 动作$a\in\{\text{继续听},\text{短确认},\text{澄清},\text{请求授权},\text{正式回答}\}$
    \If{有新版本替换了字段}\Comment{更正优先}
      \State 取消依赖旧版本且未执行的派生动作；迟到结果按版本号标为过期
    \EndIf
    \If{$q=\text{未结束}$}
      \State \Return 继续听
    \EndIf
    \If{必需字段缺失}
      \State \Return 澄清（询问缺失字段）
    \EndIf
    \If{涉及写操作且未获授权}
      \State \Return 请求授权
    \EndIf
    \If{$R$中有当前版本的请求尚未返回}
      \State \Return 若$b$即将耗尽则短确认（“正在查询”），否则继续等待
    \EndIf
    \State \Return 正式回答（以当前版本的已返回结果为依据）
  \end{algorithmic}
\end{algorithm}

对“明天……不，后天”，第一次只创建暂定日期，第二次替换日期并取消依赖旧版本的待发请求；已经发出的查询返回时先核对版本，不把旧结果写进新状态。\term{短应答}（backchannel）如“嗯，我在听”只表示收到，不当作业务许可。教学门控看到“日期已确认、人数缺失、轮次结束”时，输出澄清人数并进入等待；看到“字段完整、查询在途”时，至多短确认，不编造可预约时段。评价除动作标签的准确率外，还要统计实际的抢话次数、把暂定值提交给工具的次数，以及澄清后缺失信息是否真的补回。

\subsection{全双工语音对话：Moshi}
\label{sec:nlp-speech-moshi}

级联与半级联系统都按话轮交替工作，何时说、何时停取决于外部检测与门控。Moshi换了一种建模方式：用户与系统两路声音都离散成词元，同一个模型在共同的80\,ms时间网格上，每步同时读取用户的新声音并生成系统的下一帧声音\scite{defossez2024moshi}。说话、沉默、交叠和短回应都只是两路声音在时间轴上的不同模式，不再需要单独的轮次决策模块。实现这一点需要三个部件：低帧率且带语义的音频编解码器Mimi，能在每步生成多个码流的RQ-Transformer，以及让文字先于声音出现的内心独白。

第\ref{sec:nlp-speech-codec}小节的神经编解码器以残差向量量化（式\eqref{eq:nlp-speech-rvq}）把每帧压成若干级码本索引。语言模型直接建模这些码流时，帧率决定序列长度，所以Mimi把帧率压到12.5\,Hz。编码器是因果卷积网络，各层步长依次为4、5、6、8、2，总步长1920，把24\,kHz波形变为每80\,ms一帧的512维潜表示；量化前后各有一个8层因果Transformer，增强帧间建模。量化用8个码本、每个2048项，码率为$12.5\times8\times\log_2 2048=1100$\,bit/s，约1.1\,kbps。

Mimi的第一级码本经过\term{语义蒸馏}（semantic distillation）：用WavLM（第\ref{sec:nlp-speech-hubert}小节）提取同一段录音的表示，平均池化到12.5\,Hz，把第一级量化输出线性投影后与之计算余弦距离，作为附加损失。若第一级仍处在普通残差量化链的开头，后续各级要量化它留下的残差，语义约束会与声学重建相互牵制。Mimi因此拆分量化器：一个单独的向量量化器负责语义码，一个7级残差量化器并行负责声学细节，两者输出相加后送入解码器。第一级因而成为语义词元，其余7级是声学词元。训练只用对抗损失与特征匹配损失，不用频谱重建损失；原文报告这样客观指标下降，主观音质反而提高。

每个80\,ms步要生成17个子词元：1个系统文字词元，系统声音的8级码，用户声音的8级码。若把它们展平成一条序列，7B规模的语言模型每秒要跑$12.5\times17=212.5$步。RQ-Transformer借用残差量化图像生成中的做法，把时间与码本深度分给两个Transformer\cite{nlp-speech-lee2022rqtransformer}。记第$s$步的子词元为$V_{s,1},\ldots,V_{s,K}$，$K=17$。时间Transformer读取此前所有步，每步的输入是$K$个子词元嵌入之和，输出上下文向量
\begin{equation}
 \bm z_s=\mathrm{Tr}_{\mathrm{Temp}}(V_0,\ldots,V_{s-1})\in\mathbb R^{d}.
 \label{eq:nlp-speech-rq-temporal}
\end{equation}
深度Transformer在步内沿$k$自回归，输出第$k$个子词元的logit：
\[
 \bm l_{s,1}=\mathrm{Lin}(\bm z_s),\qquad
 \bm l_{s,k}=\mathrm{Tr}_{\mathrm{Depth}}(\bm z_s,V_{s,1},\ldots,V_{s,k-1}),\quad k>1,
\]
$\operatorname{softmax}(\bm l_{s,k})$近似$V_{s,k}$在全部历史与本步前$k-1$个子词元条件下的分布。联合分布于是分解为
\begin{equation}
 p(V)=\prod_{s}\prod_{k=1}^{K}p\bigl(V_{s,k}\mid V_{<s},V_{s,<k}\bigr).
 \label{eq:nlp-speech-rq-factorization}
\end{equation}
时间Transformer与7B文本模型Helium结构相同（32层、宽4096），并由Helium初始化，上下文取3000步（约4分钟），从而继承文本知识；深度Transformer只有6层、宽1024，各$k$使用独立的投影参数。大模型每80\,ms只运行一次，小模型在步内运行17次，计算量远小于展平序列。图\ref{fig:nlp-speech-moshi-rq}画出了一个时间步内的计算。

\begin{figure*}[htbp]
  \centering
  % 第九卷 15.7.4 节：Moshi 在一个 80 ms 时间步内的 RQ-Transformer 计算（通栏，约 145 mm）。
% 蓝色：用户侧输入；黄色描边：两个 Transformer；红色：系统生成的词元与声音。
% 左侧时间 Transformer 每步运行一次，右侧深度 Transformer 在步内依次预测 K=17 个子词元。
\begin{tikzpicture}[
  x=1mm,y=1mm,
  figure picture,
  model/.style={figure node,draw=FigureModel,fill=FigureModelFill,
    minimum height=12mm,font=\figurefont\scriptsize},
  inp/.style={figure node,draw=FigureInput,fill=FigureInputFill,
    minimum height=8mm,font=\figurefont\scriptsize},
  tok/.style={figure node,draw=FigureOutput,fill=FigureOutputFill,
    minimum height=7mm,minimum width=12mm,inner xsep=1.5pt,
    font=\figurefont\scriptsize},
  utok/.style={tok,draw=FigureInput,fill=FigureInputFill},
  lab/.style={font=\figurefont\scriptsize,text=FigureMuted,align=center},
  fl/.style={figure flow},
  back/.style={figure flow,draw=FigureMuted,dash pattern=on 2pt off 1.5pt}
]
  % 时间维：每 80 ms 运行一次的大模型。
  \node[inp,minimum width=36mm] (hist) at (20,0) {历史步$V_0,\ldots,V_{s-1}$\\每步17个词元};
  \node[model,minimum width=36mm] (emb) at (20,15) {每步17个嵌入相加\\得到$\mathbb R^{4096}$输入};
  \node[model,minimum width=36mm,minimum height=16mm] (temp) at (20,34)
    {时间Transformer\\Helium初始化\\32层，$d=4096$};
  \draw[fl] (hist) -- (emb);
  \draw[fl] (emb) -- (temp);

  % 深度维：步内依次预测 17 个子词元的小模型。
  \node[model,minimum width=84mm,minimum height=14mm] (depth) at (100,34)
    {深度Transformer（6层，$d=1024$，各$k$有独立投影）\\读取$\bm z_s$与本步已生成的$V_{s,1},\ldots,V_{s,k-1}$，输出$V_{s,k}$的分布};
  \draw[fl] (temp.east) -- node[figure edge label,above=0.5mm] {$\bm z_s$} (depth.west);

  % 第 s 步的 17 个子词元。
  \node[tok] (w) at (64,58) {$W_s$\\文字};
  \node[tok] (a1) at (79,58) {$A_{s,1}$\\语义};
  \node[tok,minimum width=19mm] (aq) at (97,58) {$A_{s-\tau,2:8}$\\声学7级};
  \node[utok] (u1) at (116,58) {$A'_{s,1}$\\用户语义};
  \node[utok,minimum width=19mm] (uq) at (134,58) {$A'_{s-\tau,2:8}$\\用户声学};
  \draw[fl] (w) -- (a1);
  \draw[fl] (a1) -- (aq);
  \draw[fl] (aq) -- (u1);
  \draw[fl] (u1) -- (uq);
  \foreach \n in {w,a1,aq,u1,uq} {\draw[fl] (depth.north -| \n) -- (\n.south);}
  \node[figure edge label,font=\figurefont\scriptsize] at (64,47) {$k=1$};
  \node[figure edge label,font=\figurefont\scriptsize] at (79,47) {$k=2$};
  \node[figure edge label,font=\figurefont\scriptsize] at (97,47) {$k=3,\ldots,9$};
  \node[figure edge label,font=\figurefont\scriptsize] at (125,47) {$k=10,\ldots,17$};

  % 输出与输入的声音。
  \node[figure node,draw=FigureOutput,fill=FigureOutputFill,minimum width=32mm,
    font=\figurefont\scriptsize] (dec) at (88,82) {Mimi解码器\\系统声音，80\,ms一帧};
  \draw[fl] (a1.north) -- (dec.south -| a1);
  \draw[fl] (aq.north) -- (dec.south -| aq);
  \node[inp,minimum width=36mm] (enc) at (125,82) {用户声音$\to$Mimi编码器\\推断时取实际输入};
  \draw[fl,draw=FigureInput] (enc.south -| u1) -- (u1.north);
  \draw[fl,draw=FigureInput] (enc.south -| uq) -- (uq.north);

  % 本步结果写回历史。
  \draw[back] (w.west) -- (50,58) |- (-4,72) -- (-4,0) -- (hist.west);
  \node[lab,anchor=south west] at (-3,72.5) {第$s$步的17个词元写回历史，供下一步读取};
\end{tikzpicture}

  \caption{Moshi在第$s$个80\,ms时间步内的计算。左侧时间Transformer读取此前所有步（每步17个子词元的嵌入之和），输出$\bm z_s$；右侧深度Transformer以$\bm z_s$为条件，按$k=1,\ldots,17$依次预测文字词元、系统语义与声学词元、用户语义与声学词元，声学词元相对语义词元延迟$\tau$步。系统词元经Mimi解码为声音；推断时用户词元取自实际收到的声音，不由模型采样。本步的17个词元写回历史，供下一步读取。}
  \label{fig:nlp-speech-moshi-rq}
\end{figure*}

只建模声音词元时，生成的语句常常语言质量差。\term{内心独白}（inner monologue）让系统在每步先生成与自身声音对齐的文字词元，作为声音的先导。训练数据用Whisper转写系统一侧的录音，得到带时间戳的词；第$i$个词的起始帧$t_i$取其起始时间乘以12.5，它的子词依次放在$t_i,t_i+1,\ldots$；其余步填PAD，词前一步放EPAD标记词的开始。英语对话中约65\%的文字位置是PAD。用户一侧不转写，文字流只描述系统自己要说的话。

声学码依赖语义码：同一帧的细节要等内容确定后才好预测。Moshi让声学码相对语义码延迟$\tau$步，第$s$步的子词元为
\begin{equation}
 V_{s}=\bigl(W_s,\;A_{s,1},\;A_{s-\tau,2},\ldots,A_{s-\tau,8},\;A'_{s,1},\;A'_{s-\tau,2},\ldots,A'_{s-\tau,8}\bigr),
 \label{eq:nlp-speech-moshi-delay}
\end{equation}
$W_s$是文字词元，$A_{s,q}$是系统声音第$q$级码，$A'_{s,q}$是用户声音的对应码，$s-\tau<1$时置为零。于是预测第$s-\tau$帧的声学码时，模型已经看到该帧之后$\tau$步的语义码。原文预训练取$\tau=2$，后续微调取$\tau=1$。文字流与声音之间的相对延迟还可以改变模型的用途：文字领先声音时模型相当于流式合成，声音领先文字时相当于流式识别。

训练目标是全部子词元的交叉熵。文字词元与全部声音词元各占一半权重，声音词元内部再按类型加权：
\[
 L=\frac1S\sum_{s=1}^{S}\Bigl(\operatorname{CE}(\bm l_{s,1},V_{s,1})+\frac{\sum_{k=2}^{K}\alpha_k\operatorname{CE}(\bm l_{s,k},V_{s,k})}{\sum_{k=2}^{K}\alpha_k}\Bigr),
\]
语义码取$\alpha_k=100$、声学码取$\alpha_k=1$，使有限的模型容量优先保证说出的内容。训练时用户一路的词元也计入损失；推断时只采样系统的文字与声音，用户一路直接填入Mimi对实际用户声音的编码。这样模型说话时用户流仍在输入，交叠与打断都表示在两路词元之中。

12.5\,Hz帧率规定了表示的节拍：零延迟时最小延迟为一帧80\,ms，每多一步声学延迟加80\,ms。原文据此给出$\tau=1$时160\,ms的理论延迟，实际约200\,ms。这个数字不含网络传输与播放缓冲，端到端延迟须按墙钟另测。对系统说“按明天……”时用户插话“不是，后天”的情形，要分别检查用户声音是否进入了模型，旧回复是否停下，后续声音是否改用正确日期，以及业务状态是否同步。

与级联系统相比，Moshi的优势在于自然的交叠、短回应与低延迟，语气与情绪在两侧都保留在声音词元中。代价是对话行为由模型自行决定，很难像算法\ref{alg:nlp-speech-response-gate}那样显式规定何时必须澄清或请求授权；文字流虽然可读，工具调用、事实核验与业务状态仍需外接，并通过下一小节的协调机制与声音输出对齐。

\subsection{声音回复与行为协调}
\label{sec:nlp-speech-coordination}

声音一旦播出就无法撤回。实时回复的每个片段因此按待生成、已生成未播、已播三种状态管理，并附上任务版本、依赖字段与播放位置。更正或新的工具结果使旧版本失效时，取消过时的生成任务，清空未播队列，并记录已经播出的范围；已播的错误内容只能明确更正。用户确实试图接管时停止或衰减当前播放，短应答则不触发停播，否则用户每说一次“嗯”，系统都会丢掉有用的回复。

工具请求同样附请求标识、任务版本与执行状态。取消本地等待不等于取消服务端的写操作；收到迟到结果时先查询实际执行记录，不直接重发写请求。幂等、结果查询与失败恢复由第\ref{chap:nlp-solving-and-execution}章的执行闭环处理，声音层保存这些证据并决定何时播报。

构造一条墙钟记录：用户在3.0\,s说完，系统3.8\,s开始播放“正在查询”，工具4.6\,s返回，结果声音5.0\,s开始播放。结束到首段声音为0.8\,s，结束到结果声音为2.0\,s；短确认缩短了反馈等待，业务完成仍在4.6\,s。再设用户在4.0\,s插话改日期，系统4.15\,s停止播放，打断停止延迟为150\,ms；4.6\,s返回的是旧日期的查询结果，只能存为过期版本，不能当作新日期的答案播出。

完整评价需要业务结果与交互质量同时计分。$\tau$-Voice在模拟业务环境中按环境终态检查任务是否完成，同时记录回应、让出话轮、打断与短回应处理等行为；它的模拟时间与墙钟解耦，适合组织故障测试，部署延迟仍要在真实墙钟下测量\scite{nlp-intro-ray2026}。实际测试应保存用户结束、稳定识别、首段播放、结果播放、打断停止与业务完成等事件，口音、噪声、交叠与接口失败都应进入完整交互测试；只有这些事件共享同一时间基准，表\ref{tab:nlp-speech-contracts}中的时间比较才成立。


\section*{文献说明}
\label{sec:nlp-speech-literature}

下列一手来源按本章顺序分组，说明各自支持的论述。正文中的中文通知、格点数值、状态机与业务核验例子都是本章的教学构造，不属于这些论文的实验；引用原文结果时，各节只保留说明机制所需的个别数字。

表示部分，SpecAugment给出时间扭曲、频率遮蔽与时间遮蔽的增强策略\cite{nlp-speech-park2019}；Conformer给出卷积模块与马卡龙式前馈的编码器结构\cite{nlp-speech-gulati2020conformer}；wav2vec 2.0给出乘积量化目标上的遮蔽对比学习\cite{nlp-speech-baevski2020wav2vec2}，HuBERT与WavLM给出基于聚类伪标签的遮蔽预测及其去噪扩展\cite{nlp-speech-hsu2021hubert,nlp-speech-chen2022wavlm}；VQ-VAE给出向量量化与直通梯度的训练方式\cite{nlp-speech-oord2017vqvae}，SoundStream与EnCodec给出残差向量量化编解码器、量化器丢弃与多尺度判别器\cite{nlp-speech-zeghidour2022soundstream,nlp-speech-defossez2023encodec}，AudioLM给出语义词元与声学词元的分工\cite{nlp-speech-borsos2023audiolm}。

识别与整理部分，CTC给出对齐路径求和的序列损失\cite{nlp-speech-graves2006}，LAS给出注意力编码器—解码器识别器\cite{nlp-speech-chan2015}，联合CTC/注意力训练把两种目标结合起来\cite{nlp-speech-kim2017}；RNN-T给出带输出历史的转导器\cite{nlp-speech-graves2012}，剪枝RNN-T降低其训练显存\cite{nlp-speech-kuang2022}；Whisper给出大规模弱监督的多任务识别\cite{radford2023whisper}，Whisper-Streaming给出在其上构造流式转写的策略\cite{nlp-speech-machacek2023}；x-vector给出说话者嵌入\cite{nlp-speech-snyder2018}。

理解与合成部分，SLURP提供口语语义理解的数据与评价\cite{nlp-speech-bastianelli2020}。Tacotron 2给出字符到梅尔谱再接条件WaveNet的合成\cite{nlp-intro-shen2018}，FastSpeech 2给出带方差适配器的非自回归声学模型\cite{nlp-speech-ren2021}；WaveNet与HiFi-GAN分别代表自回归与对抗式神经声码器\cite{nlp-speech-oord2016wavenet,nlp-speech-kong2020}；VITS给出变分推断、规范化流与对抗训练结合的端到端合成\cite{nlp-speech-kim2021vits}；Grad-TTS与F5-TTS分别代表扩散与流匹配合成\cite{nlp-speech-popov2021gradtts,nlp-speech-chen2024f5tts}；VALL-E给出以编解码词元为目标的语言模型式合成\cite{nlp-speech-wang2023valle}。

翻译与对话部分，Translatotron与S2UT分别给出直接频谱生成与离散单位生成的语音到语音翻译\cite{nlp-speech-jia2019translatotron,nlp-speech-lee2022s2ut}，UnitY把直接转换拆成先生成文字、再生成单位的两阶段解码\cite{nlp-speech-inaguma2023unity}，两篇Seamless论文分别给出多任务两阶段模型与后续的流式策略\cite{nlp-intro-seamless2023,nlp-speech-seamless2023streaming}，EMMA给出流式翻译所用的高效单调多头注意力\cite{nlp-speech-ma2023emma}。Qwen-Audio与SALMONN给出语音编码器接入大语言模型的方式\cite{nlp-speech-chu2023qwenaudio,nlp-speech-tang2024salmonn}，VAP给出未来声音活动预测\cite{nlp-speech-ekstedt2022vap}，Moshi给出平行声音流与全双工对话模型\cite{defossez2024moshi}。

\section{本章小结}
\label{sec:nlp-speech-summary}

回到章首的问题：每个字都转写正确，系统仍可能抢答或把日期读错。这说明声音任务至少要分开检查三件事：听见了什么，文字怎样对应声音的时间，哪些结果已经稳定到可以交付。

表示层面，log-Mel与MFCC把波形压成每10\,ms一帧的低维特征，SpecAugment在其上模拟局部缺失；卷积前端把帧率降到40\,ms，Conformer用深度卷积建模局部模式、用相对位置自注意力建模全局依赖。自监督表示用无标注录音换取对转写的需求：wav2vec 2.0在量化目标上做遮蔽对比，HuBERT与WavLM改为预测离线聚类得到的伪标签。神经编解码器以残差向量量化把声音变成多路离散码流，与自监督聚类得到的语义词元一起，使声音可以作为语言模型的输入与输出。

识别层面，HMM、CTC、注意力编码器—解码器与RNN-T的差别落在两点：怎样处理声音与文字之间的对齐，以及每一步能否读取已输出的历史。流式能力还取决于编码器允许看见多少未来声音，并不只由损失决定。识别之后，还要恢复当前有效的日期、数量与更正，并整理标点、数字与说话者。

生成层面，合成从文本前端、声学模型与声码器的流水线，发展到联合训练与编解码词元的语言模型生成；翻译从级联走到直接生成目标语言的离散单位。无论采用哪条路线，内容都必须在最终的声音上核验，目标文字正确并不保证目标声音正确。

语音大模型与全双工模型把声音并入语言模型，但话轮是否合适、打断是否及时、业务是否完成，仍要依靠录音时间轴、墙钟事件与执行证据来判定。图像以及音画联合的处理，由卷十展开。
